% Les alcanes — Chimie 1ère TI — Cours signé OnBuch+
% Programme officiel MINESEC, Première C, D, E, TI. Compiler avec : tectonic N01-les-alcanes.tex
\newcommand{\DOCMATIERE}{Chimie}
\newcommand{\DOCNIVEAU}{1ère TI}
\newcommand{\DOCMODULE}{Module 1 : Chimie organique}
\newcommand{\DOCLECON}{01}
\newcommand{\DOCTITRE}{Les alcanes}
% OnBuch+ — préambule « cours signés » ENRICHI (compilé avec tectonic / XeLaTeX)
% Système visuel « Pop » d'OnBuch+ : fond crème (jamais blanc), bordures encre
% épaisses, ombres « dures » décalées, titres Archivo Black en capitales.
% Version enrichie pour les cours de chimie : chimie (mhchem, chemfig),
% graphes (pgfplots), boîtes pédagogiques supplémentaires (définition,
% propriété, méthode, attention, le savais-tu, exercice résolu, exercice,
% TP, bilan), page de garde refaite et signature OnBuch+ ancrée en bas.
%
% Macros attendues dans main.tex AVANT \input{preamble} :
%   \DOCMATIERE  \DOCNIVEAU  \DOCMODULE  \DOCLECON (numéro)  \DOCTITRE
\documentclass[11pt]{article}

\usepackage{fontspec}
\usepackage[french]{babel}
\usepackage[a4paper,margin=2cm,top=3.4cm,bottom=2.8cm,headheight=1.4cm,footskip=1.3cm]{geometry}
\usepackage{amsmath,amssymb,amsfonts}
\usepackage[table]{xcolor} % [table] : fournit aussi \rowcolors (alternance de lignes)
\usepackage{pagecolor}
\usepackage{tikz}
\usetikzlibrary{arrows.meta,positioning,calc,shapes.geometric,shapes.misc,shapes.arrows,decorations.pathmorphing,decorations.markings,patterns,shadows,mindmap,trees,fit,backgrounds,matrix,intersections}
\usepackage{pgfplots}
\pgfplotsset{compat=1.18}
\usepackage[version=4]{mhchem}
\usepackage{chemfig}
\usepackage{enumitem}
\usepackage{fancyhdr}
% [most] charge toutes les bibliothèques tcolorbox (listings, minted, raster,
% documentation...) et gonfle inutilement la mémoire TeX sur un document long
% (~300 boîtes) : on ne charge que ce qui est réellement utilisé.
\usepackage[skins,breakable]{tcolorbox}
\usepackage{graphicx}
\usepackage{colortbl}
\usepackage{booktabs}
\usepackage{tabularx}
\usepackage{array}
\usepackage{multicol}
\usepackage{siunitx}
% parse-numbers=false : accepte aussi \SI{2,5 \times 10^{-3}}{...}
\sisetup{locale=FR,output-decimal-marker={,},group-minimum-digits=4,parse-numbers=false}
\DeclareSIUnit\ohm{\ensuremath{\symup{\Omega}}} % U+2126 absent de la police
\usepackage{qrcode}
% \degree : commande fréquemment utilisée par les modèles pour un angle ou une
% température en dehors de \SI{}{} (non fournie par les paquets chargés).
\providecommand{\degree}{^{\circ}}
\usepackage{lastpage}
\usepackage{hyperref}
\hypersetup{hidelinks}

% ============================================================
% Palette « Pop » OnBuch+ — mêmes hex que la classe P (plus_ui.dart)
% ============================================================
\definecolor{poporange}{HTML}{F59321}
\definecolor{popdark}{HTML}{E07A0C}
\definecolor{popcream}{HTML}{FFF6E8}
\definecolor{popink}{HTML}{1C1714}
\definecolor{poppurple}{HTML}{7A5AE0}
\definecolor{popgreen}{HTML}{2E9E6A}
\definecolor{poppink}{HTML}{E0457B}
\definecolor{popblue}{HTML}{2F7DE1}
\definecolor{popgold}{HTML}{C98A12}
\definecolor{popmuted}{HTML}{8A7F76}
\definecolor{popline}{HTML}{EADFD2}
\definecolor{popgray}{HTML}{8A7F76}
\colorlet{popgrey}{popgray}
% Alias tolérants (le modèle nomme parfois les couleurs différemment)
\colorlet{popwhite}{white}
\colorlet{popblack}{popink}
\colorlet{popred}{poppink}
\colorlet{popyellow}{popgold}
% Teintes claires pour fonds de tableaux / zones
\colorlet{popblueL}{popblue!10!white}
\colorlet{poporangeL}{poporange!14!white}
\colorlet{popgreenL}{popgreen!12!white}
\colorlet{poppurpleL}{poppurple!12!white}
\colorlet{poppinkL}{poppink!12!white}
\colorlet{popgoldL}{popgold!14!white}

\pagecolor{popcream}

% ============================================================
% Typographie
% ============================================================
\defaultfontfeatures{Ligatures=TeX}
\setmainfont{PlusJakartaSans-Medium.ttf}[Path=../../../fonts/,BoldFont=PlusJakartaSans-Bold.ttf]
% Maths en sans-serif (Fira Math) avec chiffres et lettres droites de Plus
% Jakarta, pour que les formules se fondent dans le texte.
\usepackage[math-style=ISO,bold-style=ISO]{unicode-math}
\setmathfont{FiraMath-Regular.otf}[Path=../../../fonts/]
\setmathfont{PlusJakartaSans-Medium.ttf}[Path=../../../fonts/,range={up/num,up/{Latin,latin},\mathup}]
\usepackage{icomma} % virgule décimale sans espace en mode math
% Caractères absents de la police de texte (sinon carré vide) : rendus par
% leur équivalent mathématique, en texte comme en formule.
\usepackage{newunicodechar}
\newunicodechar{α}{\ensuremath{\alpha}}
\newunicodechar{Α}{\ensuremath{\alpha}}
\newunicodechar{β}{\ensuremath{\beta}}
\newunicodechar{δ}{\ensuremath{\delta}}
\newunicodechar{✓}{\popcheck{popgreen}}
\newfontfamily\popdisplay{ArchivoBlack-Regular.ttf}[Path=../../../fonts/]
\newfontfamily\poplogo{SpaceGrotesk-Bold.ttf}[Path=../../../fonts/]
\newfontfamily\popsemi{PlusJakartaSans-SemiBold.ttf}[Path=../../../fonts/]
\newfontfamily\popxbold{PlusJakartaSans-ExtraBold.ttf}[Path=../../../fonts/]
\newfontfamily\popxboldls{PlusJakartaSans-ExtraBold.ttf}[Path=../../../fonts/,LetterSpace=3.0]

\setlist[enumerate]{leftmargin=1.7em,itemsep=5pt,topsep=4pt}
\setlist[itemize]{leftmargin=1.7em,itemsep=4pt,topsep=4pt,label={\color{poporange}\textbullet}}
\setlength{\parindent}{0pt}
\setlength{\parskip}{5pt}
\renewcommand{\arraystretch}{1.25}

% chemfig : liaisons un peu plus épaisses, encre Pop
\setchemfig{atom sep=2em,bond style={line width=0.9pt,popink},cram width=3pt}

% pgfplots : style Pop par défaut
\pgfplotsset{
  every axis/.append style={axis line style={popink,line width=1pt},tick style={popink},
    grid style={popline,line width=0.5pt},label style={font=\small},tick label style={font=\footnotesize}},
  popcurve/.style={poporange,line width=1.8pt,no markers,smooth},
}

% ============================================================
% Pill / squiggle
% ============================================================
\newcommand{\poppill}[3]{%
  \tikz[baseline=-0.55ex]\node[draw=popink,line width=1.2pt,rounded corners=2.6pt,%
    fill=#2,inner xsep=6.5pt,inner ysep=3pt]{%
    \popxboldls\fontsize{7}{7}\selectfont\color{#3}\MakeUppercase{#1}};%
}
\newcommand{\popsquiggle}[1]{%
  \tikz[baseline=-0.4ex]{
    \draw[#1,line width=2.6pt,line cap=round]
      plot [smooth] coordinates {(0,0.09) (0.34,0.01) (0.68,0.21) (1.02,0.01) (1.36,0.10)};
  }%
}
% Mot-clé surligné (marqueur orange sous le texte)
\newcommand{\cle}[1]{{\popxbold\color{popdark}#1}}

% ============================================================
% En-tête / pied
% ============================================================
\pagestyle{fancy}
\fancyhf{}
\renewcommand{\headrulewidth}{0pt}
\renewcommand{\footrulewidth}{0pt}
\lhead{\raisebox{-0.15ex}{\poplogo\fontsize{13}{13}\selectfont\color{popink}OnBuch\color{poporange}+}}
\rhead{\poppill{\DOCMATIERE\ \textbullet\ \DOCNIVEAU\ \textbullet\ Leçon \DOCLECON}{white}{popink}}
\lfoot{\popxbold\fontsize{7.6}{7.6}\selectfont\color{popmuted}\MakeUppercase{Cours signé OnBuch+}\ \textbullet\ {\popsemi\DOCTITRE}}
\rfoot{\tikz[baseline=-0.6ex]\node[rounded corners=8.5pt,fill=popink,minimum height=17pt,minimum width=17pt,inner xsep=5pt,inner sep=0]{\popxbold\fontsize{8}{8}\selectfont\color{popcream}\thepage\,/\,\pageref*{LastPage}};}

% ============================================================
% Titres
% ============================================================
\newcommand{\chaptitre}[2]{%
  \begin{tcolorbox}[enhanced,colback=poporange,colframe=popink,boxrule=2.6pt,arc=11pt,
      drop shadow=popink,left=15pt,right=15pt,top=12pt,bottom=11pt,before skip=3pt,after skip=18pt]
    \poppill{#2}{popink}{popcream}\par\vspace{9pt}
    {\raggedright\hyphenpenalty=10000\exhyphenpenalty=10000\popdisplay\fontsize{22}{24}\selectfont\color{popcream}\MakeUppercase{#1}\par}
  \end{tcolorbox}
}
% Section numérotée : pastille orange avec le numéro + titre Archivo
\newcounter{popsec}
\newcommand{\coursec}[1]{%
  \stepcounter{popsec}\setcounter{popsub}{0}%
  \par\vspace{16pt}\noindent
  \begin{minipage}{\linewidth}
  \tikz[baseline=-0.7ex]\node[draw=popink,line width=1.6pt,fill=poporange,rounded corners=4pt,minimum size=22pt,inner sep=0]{\popdisplay\fontsize{12}{12}\selectfont\color{popcream}\thepopsec};%
  \hspace{8pt}{\popdisplay\fontsize{15}{17}\selectfont\color{popink}\MakeUppercase{#1}}\par
  \vspace{4pt}\noindent{\color{poporange}\rule{36pt}{3.6pt}}
  \end{minipage}\par\nopagebreak\vspace{8pt}
}
\newcounter{popsub}
\newcommand{\courssub}[1]{%
  \stepcounter{popsub}%
  \par\vspace{9pt}\noindent{\popxbold\fontsize{11.5}{13}\selectfont\color{popink}{\color{poporange}\thepopsec.\thepopsub}\hspace{6pt}#1}\par\nopagebreak\vspace{4pt}
}

% ============================================================
% Boîtes pédagogiques — même recette : fond blanc, bordure épaisse colorée,
% ombre dure encre, badge plein. Titre optionnel [#1].
% ============================================================
\tcbset{popbase/.style={breakable,enhanced,colback=white,boxrule=2.3pt,arc=9pt,
  shadow={3pt}{-3pt}{0pt}{popink},left=14pt,right=14pt,top=11pt,bottom=11pt,before skip=11pt,after skip=15pt,
  fontupper=\popsemi\fontsize{10.6}{14.5}\selectfont\color{popink},
  fontlower=\popsemi\fontsize{10.6}{14.5}\selectfont\color{popink}}}
% Coche dessinée (le glyphe ✓ n'existe pas dans les polices du document)
\newcommand{\popcheck}[1]{\tikz[baseline=-0.5ex]\draw[#1,line width=1.6pt,line cap=round,line join=round](-0.13,0.0)--(-0.04,-0.09)--(0.13,0.1);}
\newcommand{\popboxhead}[3]{\poppill{#1}{#2}{white}\ifx\relax#3\relax\else\hspace{6pt}{\popxbold\fontsize{10.6}{12}\selectfont\color{popink}#3}\fi\par\vspace{7pt}}

\newtcolorbox{aretenir}[1][]{popbase,colframe=popgreen,before upper={\popboxhead{À retenir}{popgreen}{#1}}}
\newtcolorbox{exemplebox}[1][]{popbase,colframe=poporange,before upper={\popboxhead{Exemple}{poporange}{#1}}}
\newtcolorbox{definition}[1][]{popbase,colframe=popblue,before upper={\popboxhead{Définition}{popblue}{#1}}}
\newtcolorbox{propriete}[1][]{popbase,colframe=poppurple,before upper={\popboxhead{Propriété}{poppurple}{#1}}}
\newtcolorbox{methode}[1][]{popbase,colframe=popgold,before upper={\popboxhead{Méthode}{popgold}{#1}}}
\newtcolorbox{attention}[1][]{popbase,colframe=poppink,before upper={\popboxhead{Attention}{poppink}{#1}}}
\newtcolorbox{experience}[1][]{popbase,colframe=popblue,colback=popblueL,before upper={\popboxhead{Expérience}{popblue}{#1}}}
% « Le savais-tu ? » : carte violette pleine (contexte, culture, Cameroun)
\newtcolorbox{savaistu}[1][]{popbase,colback=poppurple,colframe=popink,
  fontupper=\popsemi\fontsize{10.4}{14.2}\selectfont\color{white},
  before upper={\poppill{Le savais-tu\,?}{white}{poppurple}\ifx\relax#1\relax\else\hspace{6pt}{\popxbold\color{white}#1}\fi\par\vspace{7pt}}}
% Exercice résolu : énoncé en haut, \tcblower puis la solution sur fond crème
\newcounter{popexo}\newcounter{popexr}
\newtcolorbox{exoresolu}[1][]{popbase,colframe=poporange,colbacklower=popcream,
  segmentation style={popink,line width=1.2pt,dashed},
  before upper={\stepcounter{popexr}\popboxhead{Exercice résolu \thepopexr}{poporange}{#1}},
  before lower={\poppill{Solution}{popink}{popcream}\par\vspace{6pt}}}
% Exercice (non résolu en place) — la correction est dans la partie Corrigés
\newtcolorbox{exercice}[1][]{popbase,colframe=popink,
  before upper={\stepcounter{popexo}\popboxhead{Exercice \thepopexo}{popink}{#1}}}
% Corrigé
\newtcolorbox{corrige}[1][]{popbase,colframe=popgreen,colback=popgreenL,
  before upper={\popboxhead{Corrigé}{popgreen}{#1}}}
% Objectifs / prérequis / bilan
\newtcolorbox{objectifs}{popbase,colframe=popink,colback=poporangeL,
  before upper={\popboxhead{Objectifs de la leçon}{popink}{}}}
\newtcolorbox{prerequis}{popbase,colframe=popink,colback=popblueL,
  before upper={\popboxhead{Prérequis}{popblue}{}}}
\newtcolorbox{bilan}{popbase,colframe=popink,colback=white,boxrule=2.8pt,
  before upper={\popboxhead{Fiche bilan}{poporange}{L'essentiel à connaître pour l'examen}}}
% Figure encadrée avec légende
\newtcolorbox{popfigure}[1][]{enhanced,breakable=false,colback=white,colframe=popline,boxrule=1.4pt,arc=8pt,
  left=10pt,right=10pt,top=10pt,bottom=8pt,before skip=10pt,after skip=12pt,halign=center,
  lower separated=false,halign lower=center,
  fontlower=\popsemi\fontsize{9}{11.5}\selectfont\color{popmuted},
  #1}
\newcommand{\legende}[1]{\par\vspace{6pt}{\popsemi\fontsize{9}{11.5}\selectfont\color{popmuted}#1}}

% ============================================================
% Signature OnBuch+
% ============================================================
\newcommand{\poplogoinline}[1]{{\poplogo\fontsize{#1}{#1}\selectfont\color{popink}OnBuch\color{poporange}+}}
\newcommand{\signatureonbuch}{%
  \begin{tcolorbox}[enhanced,colback=popink,colframe=popink,boxrule=2.2pt,arc=10pt,
      drop shadow=poporange,left=14pt,right=14pt,top=10pt,bottom=10pt,before skip=0pt,after skip=0pt]
    \tikz[baseline=-0.7ex]\node[circle,fill=popgreen,draw=popcream,line width=1.3pt,minimum size=22pt,inner sep=0]{\popcheck{white}};%
    \hspace{9pt}\begin{minipage}[c]{0.62\linewidth}
      {\popxbold\fontsize{12}{13}\selectfont\color{popcream}Signé\ \textbullet\ {\poplogo OnBuch\color{poporange}+}}\par\vspace{2pt}
      {\raggedright\popsemi\fontsize{8}{10.5}\selectfont\color{popline}\MakeUppercase{Équipe pédagogique OnBuch+}\\\MakeUppercase{Conforme au programme officiel MINESEC}\par}
    \end{minipage}\hfill
    {\poplogo\fontsize{22}{22}\selectfont\color{popcream}OnBuch\color{poporange}+}
  \end{tcolorbox}%
}

% ============================================================
% Page de garde
% #1 titre · #2 matière · #3 niveau · #4 module · #5 n° leçon · #6 durée ·
% #7 description / objectifs (texte ou itemize)
% ============================================================
\newcommand{\pagegarde}[7]{%
  \thispagestyle{empty}
  \newgeometry{a4paper,margin=1.7cm,top=1.6cm,bottom=1.4cm}%
  \begin{tikzpicture}[remember picture,overlay]
    \fill[poporange!18] ([xshift=-3.2cm,yshift=-3.6cm]current page.north east) circle (4.6cm);
    \fill[poppurple!14] ([xshift=2.4cm,yshift=7.6cm]current page.south west) circle (3.8cm);
  \end{tikzpicture}%
  {\poplogo\fontsize{30}{30}\selectfont\color{popink}OnBuch\color{poporange}+}\hfill
  \raisebox{0.6ex}{\poppill{Cours signé \textbullet\ Premium}{popink}{popcream}}\par
  \vspace{1pt}\popsquiggle{poppurple}\par
  \vspace{8pt}
  \begin{tcolorbox}[enhanced,colback=poporange,colframe=popink,boxrule=2.8pt,arc=13pt,
      drop shadow=popink,left=18pt,right=18pt,top=12pt,bottom=13pt,after skip=10pt]
    \poppill{#2 \textbullet\ #3}{popink}{popcream}\hspace{5pt}\poppill{#4}{white}{popink}\par\vspace{8pt}
    {\popdisplay\fontsize{14}{14}\selectfont\color{popink}LEÇON #5}\par\vspace{4pt}
    {\raggedright\hyphenpenalty=10000\exhyphenpenalty=10000\popdisplay\fontsize{25}{27}\selectfont\color{popcream}\MakeUppercase{#1}\par}
    \vspace{8pt}
    {\popxbold\fontsize{9.5}{9.5}\selectfont\color{popink}\MakeUppercase{Durée conseillée : #6}}
  \end{tcolorbox}
  \begin{tcolorbox}[enhanced,colback=white,colframe=popink,boxrule=2.2pt,arc=10pt,
      drop shadow=popink,left=16pt,right=16pt,top=10pt,bottom=10pt,after skip=8pt]
    \poppill{Dans cette leçon}{poppurple}{white}\par\vspace{5pt}
    \popsemi\fontsize{9.3}{12.6}\selectfont\color{popink}#7
  \end{tcolorbox}
  \noindent\begin{minipage}[t]{0.487\linewidth}
    \begin{tcolorbox}[enhanced,colback=popgreen,colframe=popink,boxrule=2.2pt,arc=10pt,
        drop shadow=popink,left=11pt,right=11pt,top=10pt,bottom=10pt,height=3.1cm,valign=center]
      \tcbox[colback=white,colframe=popink,boxrule=1.3pt,arc=5pt,boxsep=0pt,nobeforeafter,
        left=3pt,right=3pt,top=3pt,bottom=3pt,valign=center]{\qrcode[height=1.9cm]{https://play.google.com/store/apps/details?id=com.luvvix.onbuch&hl=fr}}%
      \hspace{8pt}\begin{minipage}[c]{0.5\linewidth}
        {\popdisplay\fontsize{11}{12.5}\selectfont\color{white}\MakeUppercase{Télécharge OnBuch}}\par\vspace{4pt}
        {\raggedright\popsemi\fontsize{8.4}{11}\selectfont\color{white}Gratuit \textbullet\ Google Play\\Tuteur IA, annales, concours, résultats\par}
      \end{minipage}
    \end{tcolorbox}
  \end{minipage}\hfill
  \begin{minipage}[t]{0.487\linewidth}
    \begin{tcolorbox}[enhanced,colback=poppurple,colframe=popink,boxrule=2.2pt,arc=10pt,
        drop shadow=popink,left=11pt,right=11pt,top=10pt,bottom=10pt,height=3.1cm,valign=center]
      \tikz[baseline=-0.7ex]\node[circle,fill=white,draw=popink,line width=1.3pt,minimum size=1.6cm,inner sep=0]{\popdisplay\fontsize{24}{24}\selectfont\color{poppurple}+};%
      \hspace{8pt}\begin{minipage}[c]{0.6\linewidth}
        {\popdisplay\fontsize{11}{12.5}\selectfont\color{white}\MakeUppercase{Passe à OnBuch+}}\par\vspace{4pt}
        {\raggedright\popsemi\fontsize{8.4}{11}\selectfont\color{white}Tous les cours signés, Tuteur IA illimité, fiches PDF — onglet OnBuch+ de l'app\par}
      \end{minipage}
    \end{tcolorbox}
  \end{minipage}
  \vspace{8pt}
  \popcontenu
  \vfill
  \signatureonbuch
  \restoregeometry
  \newpage
}

% Rangée « Ce cours contient » (page de garde)
\newcommand{\poptile}[3]{% #1 couleur #2 gros texte #3 légende
  \begin{tcolorbox}[enhanced,colback=#1,colframe=popink,boxrule=2pt,arc=9pt,shadow={2.5pt}{-2.5pt}{0pt}{popink},
      left=6pt,right=6pt,top=9pt,bottom=9pt,halign=center,valign=center,height=1.75cm,width=\linewidth,nobeforeafter]
    {\popdisplay\fontsize{12.5}{13}\selectfont\color{white}\MakeUppercase{#2}}\par\vspace{3pt}
    {\centering\popsemi\fontsize{7.8}{9.5}\selectfont\color{white}#3\par}
  \end{tcolorbox}}
\newcommand{\popcontenu}{%
  \par\noindent{\popxboldls\fontsize{7.5}{7.5}\selectfont\color{popmuted}CE COURS SIGNÉ CONTIENT}\par\vspace{7pt}
  \noindent\begin{minipage}[t]{0.235\linewidth}\poptile{popblue}{Cours}{complet, illustré, pas à pas}\end{minipage}\hfill
  \begin{minipage}[t]{0.235\linewidth}\poptile{poporange}{Méthodes}{et exercices résolus}\end{minipage}\hfill
  \begin{minipage}[t]{0.235\linewidth}\poptile{poppink}{Exercices}{type examen + corrigés}\end{minipage}\hfill
  \begin{minipage}[t]{0.235\linewidth}\poptile{popgreen}{Fiche bilan}{l'essentiel à retenir}\end{minipage}\par}

% Sommaire
\newtcolorbox{sommaire}{enhanced,colback=white,colframe=popink,boxrule=2.2pt,arc=10pt,
  drop shadow=popink,left=18pt,right=18pt,top=13pt,bottom=13pt,before skip=6pt,after skip=20pt,
  before upper={\poppill{Sommaire}{poppurple}{white}\par\vspace{9pt}\popsemi\fontsize{10.6}{14.5}\selectfont\color{popink}}}

% Signature de fin de cours — poussée tout en bas de la dernière page
\newcommand{\findecours}{%
  \par\vfill\nopagebreak
  \begin{center}\popsquiggle{poporange}\end{center}\vspace{4pt}
  \signatureonbuch
}
% Glyphes absents de Fira Math : redéfinis à partir de glyphes disponibles
\AtBeginDocument{%
  \def\setminus{\mathbin{\backslash}}\let\smallsetminus\setminus
  \def\vdots{\mathord{\vbox{\baselineskip3.2pt\lineskiplimit0pt\kern4pt\hbox{.}\hbox{.}\hbox{.}}}}%
  \def\ddots{\mathinner{\mkern1mu\raise6.5pt\hbox{.}\mkern2mu\raise3.6pt\hbox{.}\mkern2mu\raise0.7pt\hbox{.}\mkern1mu}}%
  \def\triangle{\mathord{\scalebox{0.9}{$\symup{\Delta}$}}}%
  \def\nabla{\mathord{\raisebox{\depth}{\scalebox{1}[-1]{$\symup{\Delta}$}}}}%
  \def\checkmark{\popcheck{popdark}}%
  \def\star{\mathbin{\ast}}\def\bigstar{\mathord{\ast}}%
  \def\diamond{\mathbin{\scalebox{0.6}{$\Diamond$}}}%
  \def\bigcup{\mathop{\mathchoice{\raisebox{-0.3em}{\scalebox{1.7}{$\cup$}}}{\scalebox{1.25}{$\cup$}}{\cup}{\cup}}\displaylimits}%
  \def\bigcap{\mathop{\mathchoice{\raisebox{-0.3em}{\scalebox{1.7}{$\cap$}}}{\scalebox{1.25}{$\cap$}}{\cap}{\cap}}\displaylimits}%
  \def\lesseqqgtr{\mathrel{\lessgtr}}\def\bigoplus{\mathop{\oplus}\displaylimits}%
}
\providecommand{\ssi}{si et seulement si} % abréviation fréquente

%% FIGFIX-BEGIN v4 — correctifs globaux des figures (docs/onbuch-generation/tools/figfix)
\usepackage{adjustbox}\usepackage{etoolbox}
% toute figure TikZ/pgfplots plus large que la ligne est réduite à la largeur disponible
\BeforeBeginEnvironment{tikzpicture}{\begin{adjustbox}{max width=\linewidth}}
\AfterEndEnvironment{tikzpicture}{\end{adjustbox}}
% légendes pgfplots lisibles par-dessus les courbes
\pgfplotsset{every axis/.append style={legend style={fill=white,fill opacity=0.92,text opacity=1,draw=black!15,inner sep=3pt}}}
% étiquettes d'arêtes (syntaxe edge["..."]) avec fond blanc
\tikzset{every edge quotes/.append style={fill=white,inner sep=1.2pt}}
%% FIGFIX-END
\begin{document}

\pagegarde{Les alcanes}{Chimie}{1ère TI}{Module 1 : Chimie organique}{01}{6 h}{Cette leçon introduit la chimie organique à travers les alcanes, hydrocarbures saturés de formule générale CnH2n+2. De la structure tétragonale du carbone à la combustion du butane, en passant par la nomenclature et l'isomérie, elle fournit toutes les bases indispensables pour la suite du programme et le Probatoire.
\vspace{4pt}
\begin{multicols}{2}\small
\begin{itemize}
\item À la fin de cette leçon, je sais décrire la structure tétragonale de l'atome de carbone et représenter spatialement le méthane et l'éthane
\item À la fin de cette leçon, je sais distinguer les différentes conformations d'un alcane à l'aide de modèles moléculaires
\item À la fin de cette leçon, je sais écrire la formule générale des alcanes et celle des cyclanes
\item À la fin de cette leçon, je sais écrire les formules développées et semi-développées des alcanes de formule CnH2n+2 avec n ≤ 5
\item À la fin de cette leçon, je sais identifier les isomères de constitution et nommer les alcanes selon les règles de nomenclature
\item À la fin de cette leçon, je sais décrire les propriétés physiques des alcanes (état, solubilité, densité)
\end{itemize}\end{multicols}}

\begin{sommaire}
\begin{multicols}{2}
\begin{enumerate}
\item Structure des alcanes : le carbone tétragonal
\item Formule générale, formules développées et isomérie
\item Nomenclature des alcanes
\item Propriétés physiques des alcanes
\item Propriétés chimiques : combustion et substitution
\item Préparation du méthane au laboratoire
\item Travaux pratiques
\item Méthodes et exercices résolus
\item Exercices
\item Corrigés des exercices
\item Fiche bilan
\end{enumerate}
\end{multicols}
\end{sommaire}

\begin{objectifs}
\begin{itemize}
\item À la fin de cette leçon, je sais décrire la structure tétragonale de l'atome de carbone et représenter spatialement le méthane et l'éthane
\item À la fin de cette leçon, je sais distinguer les différentes conformations d'un alcane à l'aide de modèles moléculaires
\item À la fin de cette leçon, je sais écrire la formule générale des alcanes et celle des cyclanes
\item À la fin de cette leçon, je sais écrire les formules développées et semi-développées des alcanes de formule CnH2n+2 avec n ≤ 5
\item À la fin de cette leçon, je sais identifier les isomères de constitution et nommer les alcanes selon les règles de nomenclature
\item À la fin de cette leçon, je sais décrire les propriétés physiques des alcanes (état, solubilité, densité)
\item À la fin de cette leçon, je sais écrire et équilibrer les équations de combustion complète et incomplète d'un alcane
\item À la fin de cette leçon, je sais écrire les équations des réactions de substitution d'un alcane avec le dichlore ou le dibrome et citer l'importance des dérivés halogénés
\item À la fin de cette leçon, je sais décrire et schématiser le dispositif de préparation du méthane à partir du carbure d'aluminium
\end{itemize}
\end{objectifs}

\begin{prerequis}
\begin{itemize}
\item Structure de l'atome de carbone et notion de valence (Seconde)
\item Liaison covalente et représentation de Lewis (Seconde)
\item Écriture et équilibrage d'une équation-bilan (Seconde)
\item Notion de mole et calculs de quantités de matière (début de Première)
\item Hydrocarbures et distillation du pétrole (notions vues en Seconde)
\end{itemize}
\end{prerequis}

\newpage

\coursec{Structure des alcanes : le carbone tétragonal}

À Yaoundé comme à Douala, des milliers de familles cuisinent chaque jour grâce au gaz butane vendu en bouteilles. Ce gaz, qui brûle avec une flamme bleue, appartient à une famille de composés organiques appelés \cle{alcanes}. Mais pourquoi le butane brûle-t-il si bien ? Pourquoi doit-on l'utiliser dans un endroit aéré ? Et comment les chimistes nomment-ils et représentent-ils ces molécules ? Cette leçon répond à ces questions en explorant la structure, les propriétés et les réactions des alcanes.

Tout commence par une question d'apparence simple : quelle est la forme d'une molécule d'alcane ? Tu vas voir que la réponse — un petit tétraèdre autour de chaque atome de carbone — explique presque tout le reste du chapitre.

\courssub{Le carbone, atome à quatre liaisons}

\textbf{Intuition.} Imagine une main avec quatre doigts : chaque doigt peut saisir un objet, mais pas deux à la fois. L'atome de carbone fonctionne exactement ainsi : il possède \emph{quatre} « bras » pour se lier à d'autres atomes, ni trois, ni cinq. C'est cette règle absolue qui gouverne toute la chimie organique.

L'atome de carbone a pour numéro atomique $Z = 6$. Sa configuration électronique est $1s^2\,2s^2\,2p^2$ : il possède donc \cle{4 électrons de valence} sur sa couche externe. Pour acquérir la structure stable du gaz noble le plus proche (le néon, octet complet), il doit mettre en commun 4 électrons supplémentaires. Il forme donc \textbf{quatre liaisons covalentes}.

\begin{definition}[Alcane]
Un \cle{alcane} est un \cle{hydrocarbure saturé} : sa molécule n'est constituée que d'atomes de carbone et d'hydrogène, et ne comporte \textbf{que des liaisons covalentes simples} \ce{C-C} et \ce{C-H}. On dit « saturé » car chaque carbone est entouré du nombre maximal d'atomes d'hydrogène possible.
\end{definition}

Dans les alcanes, chaque carbone forme donc quatre liaisons \emph{simples}. Un tel carbone porte un nom particulier.

\begin{definition}[Carbone tétragonal]
Un \cle{carbone tétragonal} (ou tétraédrique) est un atome de carbone engagé dans \textbf{quatre liaisons covalentes simples}, dirigées vers les quatre sommets d'un \cle{tétraèdre régulier} dont le carbone occupe le centre.
\end{definition}

\courssub{Géométrie tétraédrique : l'angle de 109°28'}

\textbf{Intuition.} Suppose que tu tiens quatre ballons identiques gonflés ensemble par leurs nœuds. Les quatre ballons se repoussent mutuellement (ils occupent tous de l'espace) et s'écartent spontanément jusqu'à occuper les positions les plus éloignées possibles les unes des autres : ils pointent alors vers les sommets d'un tétraèdre. Les quatre doublets d'électrons des liaisons du carbone font de même : se repoussant électriquement, ils s'orientent vers les sommets d'un tétraèdre régulier.

\begin{propriete}[Géométrie du carbone tétragonal]
Autour d'un carbone tétragonal, les quatre liaisons simples sont dirigées vers les sommets d'un tétraèdre régulier. L'angle entre deux liaisons quelconques vaut :
\[ \widehat{\ce{H-C-H}} = 109^{\circ}28' \approx 109,5^{\circ} \]
Cette valeur se démontre par la géométrie du tétraèdre régulier : $\cos\theta = -\dfrac{1}{3}$, d'où $\theta \approx 109,47^{\circ}$.
\end{propriete}

\begin{attention}[Le piège du dessin à plat]
Erreur très fréquente au Probatoire : dessiner le méthane \ce{CH4} à plat, en croix, avec quatre angles droits de $90^{\circ}$. C'est \textbf{faux} ! La molécule de méthane n'est \textbf{pas plane} : elle est \textbf{tétraédrique}. Un dessin en croix n'est qu'une formule de Lewis (qui montre les liaisons, pas la géométrie). Dès qu'on parle de forme dans l'espace, il faut penser au tétraèdre et à l'angle de $109^{\circ}28'$.
\end{attention}

\courssub{Étude du méthane \ce{CH4}}

Le méthane est le plus simple des alcanes : un seul atome de carbone, quatre atomes d'hydrogène. C'est le principal constituant du gaz naturel et du biogaz produit par la fermentation des déchets organiques (marécages, fosses à compost, digesteurs agricoles).

\textbf{Formule de Lewis.} Le carbone met en commun chacun de ses 4 électrons de valence avec l'électron d'un atome d'hydrogène :

\begin{center}
\chemfig{H-C(-[2]H)(-[6]H)-H}
\end{center}

Chaque trait représente un doublet liant (une liaison covalente simple \ce{C-H}). Le carbone est entouré de 8 électrons (octet satisfait) ; chaque hydrogène est entouré de 2 électrons (structure de l'hélium).

\textbf{Représentation spatiale.} Les quatre atomes d'hydrogène occupent les sommets d'un tétraèdre régulier, le carbone étant au centre.

\begin{popfigure}
\begin{center}
\begin{tikzpicture}[scale=1.6, line join=round]
  % Sommets du tétraèdre
  \coordinate (C) at (0,0);
  \coordinate (H1) at (1.1,0.9);    % avant droit
  \coordinate (H2) at (-1.2,0.7);   % avant gauche
  \coordinate (H3) at (0.2,-1.3);   % bas
  \coordinate (H4) at (0.1,1.5);    % arrière haut
  % Liaisons pleines (dans le plan)
  \draw[line width=1.2pt, popblue] (C) -- (H1);
  \draw[line width=1.2pt, popblue] (C) -- (H2);
  \draw[line width=1.2pt, popblue] (C) -- (H3);
  % Liaison pointillée (en arrière)
  \draw[line width=1.2pt, poporange, dashed] (C) -- (H4);
  % Atomes
  \fill[popdark] (C) circle (0.16) node[below left=1pt, text=popdark]{\textbf{C}};
  \fill[popblue] (H1) circle (0.11) node[right, text=popblue]{\textbf{H}};
  \fill[popblue] (H2) circle (0.11) node[left, text=popblue]{\textbf{H}};
  \fill[popblue] (H3) circle (0.11) node[below, text=popblue]{\textbf{H}};
  \fill[poporange] (H4) circle (0.11) node[above, text=poporange]{\textbf{H}};
  % Angle 109°28'
  \draw[popink, ->] (0.55,0.45) arc (39:150:0.68);
  \node[popink] at (0,0.62) {\small $109^{\circ}28'$};
\end{tikzpicture}
\end{center}
\legende{Structure tétraédrique du méthane \ce{CH4} : le carbone est au centre du tétraèdre, les quatre hydrogènes aux sommets. L'angle entre deux liaisons \ce{C-H} vaut $109^{\circ}28'$.}
\end{popfigure}

\textbf{Modèles moléculaires.} Deux types de modèles décrivent la même molécule :
\begin{itemize}
\item le \cle{modèle éclaté} (boules reliées par des tiges) : il montre les liaisons et les angles, mais exagère les distances ;
\item le \cle{modèle compact} (sphères qui se touchent) : il montre l'encombrement réel des atomes, mais masque les liaisons.
\end{itemize}

\begin{methode}[Représenter un alcane en 3D : la convention de Cram]
Pour dessiner une molécule tétraédrique sur une feuille plane, on utilise la \cle{représentation de Cram} :
\begin{enumerate}
\item un \textbf{trait plein fin} : la liaison est dans le plan de la feuille ;
\item un \textbf{trait plein élargi} (coin) : la liaison sort du plan, vers le lecteur ;
\item un \textbf{trait hachuré ou pointillé} : la liaison rentre derrière le plan.
\end{enumerate}
Autour de chaque carbone tétragonal, on dessine en général deux liaisons dans le plan, une en avant et une en arrière.
\end{methode}

\begin{popfigure}
\begin{center}
\chemfig{H-C(-[2]H)(<:[6]H)(<[7]H)-H}
\hspace{2cm}
\chemfig{CH_3-[:-30]C(-[2]H)(<[6]H)(<:[-60]H)}
\end{center}
\legende{Représentations de Cram : à gauche, le méthane \ce{CH4} ; à droite, un carbone de l'éthane \ce{C2H6} (liaison \ce{C-C} dans le plan vers la gauche, trois liaisons \ce{C-H} : une dans le plan vers le haut, une en coin vers l'avant, une en hachuré vers l'arrière).}
\end{popfigure}

\courssub{Étude de l'éthane \ce{C2H6} : la liaison C–C et la libre rotation}

L'éthane contient deux atomes de carbone. Chaque carbone, tétragonal, forme trois liaisons \ce{C-H} et une liaison \ce{C-C} avec l'autre carbone. Sa formule semi-développée est \ce{CH3-CH3}.

\textbf{Intuition.} Pense à deux tétraèdres soudés par un sommet commun. La liaison \ce{C-C} qui les relie agit comme un \emph{axe} : les deux groupes \ce{CH3} peuvent \textbf{tourner librement l'un par rapport à l'autre} autour de cet axe, comme deux roues montées sur un même essieu. C'est la \cle{libre rotation} autour d'une liaison simple.

\begin{definition}[Conformation]
Les différentes positions relatives que prennent les groupes \ce{CH3} de l'éthane lors de la rotation autour de la liaison \ce{C-C} sont appelées \cle{conformations}. Ce ne sont \textbf{pas} des molécules différentes : on passe de l'une à l'autre par simple rotation, sans rompre aucune liaison.
\end{definition}

Pour visualiser ces conformations, on utilise la \cle{projection de Newman} : on regarde la molécule \emph{le long} de l'axe \ce{C-C}. Le carbone avant est figuré par un point (centre du cercle), le carbone arrière par le cercle lui-même. Deux conformations extrêmes existent :

\begin{itemize}
\item la conformation \cle{éclipsée} : les hydrogènes du carbone arrière sont cachés derrière ceux du carbone avant ; les liaisons \ce{C-H} sont face à face, les répulsions sont maximales ;
\item la conformation \cle{décalée} : les hydrogènes de l'arrière se placent entre ceux de l'avant (rotation de $60^{\circ}$) ; les répulsions sont minimales, c'est la conformation la plus stable.
\end{itemize}

\begin{popfigure}
\begin{center}
\begin{tikzpicture}[scale=1.15]
  % --- Conformation éclipsée ---
  \begin{scope}
    \draw[popmuted, line width=1pt] (0,0) circle (1.05);
    % Carbone avant : liaisons depuis le centre
    \foreach \a in {90,210,330}{
      \draw[line width=1.2pt, popblue] (0,0) -- (\a:0.62);
      \fill[popblue] (\a:0.62) circle (0.09) node[shift={(\a:0.28)}, text=popblue]{\scriptsize H};
    }
    % Carbone arrière : éclipsé (mêmes directions, jusqu'au cercle)
    \foreach \a in {90,210,330}{
      \draw[line width=1.2pt, poporange, dashed] (\a:0.72) -- (\a:1.05);
      \fill[poporange] (\a:1.05) circle (0.09) node[shift={(\a:0.3)}, text=poporange]{\scriptsize H};
    }
    \fill[popdark] (0,0) circle (0.07);
    \node at (0,-1.7) {\small \textbf{Conformation éclipsée}};
    \node at (0,-2.05) {\small (répulsions maximales)};
  \end{scope}
  % --- Conformation décalée ---
  \begin{scope}[xshift=5.2cm]
    \draw[popmuted, line width=1pt] (0,0) circle (1.05);
    % Carbone avant
    \foreach \a in {90,210,330}{
      \draw[line width=1.2pt, popblue] (0,0) -- (\a:0.62);
      \fill[popblue] (\a:0.62) circle (0.09) node[shift={(\a:0.28)}, text=popblue]{\scriptsize H};
    }
    % Carbone arrière : décalé de 60°
    \foreach \a in {30,150,270}{
      \draw[line width=1.2pt, poporange] (\a:0.45) -- (\a:1.05);
      \fill[poporange] (\a:1.05) circle (0.09) node[shift={(\a:0.3)}, text=poporange]{\scriptsize H};
    }
    \fill[popdark] (0,0) circle (0.07);
    \node at (0,-1.7) {\small \textbf{Conformation décalée}};
    \node at (0,-2.05) {\small (la plus stable)};
  \end{scope}
  % Flèche de rotation
  \draw[->, line width=1.2pt, popink] (1.7,0) -- (3.3,0) node[midway, above, text=popink]{\small rotation de $60^{\circ}$};
\end{tikzpicture}
\end{center}
\legende{Projections de Newman de l'éthane \ce{CH3-CH3}, vue le long de l'axe \ce{C-C}. En bleu : les trois hydrogènes du carbone avant ; en orange : ceux du carbone arrière. Une rotation de $60^{\circ}$ fait passer de la conformation éclipsée à la conformation décalée.}
\end{popfigure}

\begin{savaistu}[Pourquoi les conformations ne sont-elles pas des isomères ?]
À température ordinaire, l'agitation thermique fournit largement l'énergie nécessaire (environ \SI{12}{kJ.mol^{-1}} seulement) pour franchir la petite barrière entre deux conformations : les groupes \ce{CH3} tournent des milliards de fois par seconde ! Il est donc \textbf{impossible d'isoler} la forme éclipsée de la forme décalée. C'est pourquoi les conformations ne sont \emph{pas} des isomères séparables, contrairement aux isomères de constitution que tu rencontreras bientôt (comme le butane et le 2-méthylpropane), qui, eux, sont des molécules différentes que l'on peut mettre dans des flacons distincts.
\end{savaistu}

\courssub{Longueurs et énergies des liaisons : ordres de grandeur}

Deux grandeurs caractérisent une liaison covalente :
\begin{itemize}
\item sa \cle{longueur} : distance entre les noyaux des deux atomes liés ;
\item son \cle{énergie} : énergie qu'il faut fournir pour la rompre.
\end{itemize}

\begin{center}
\begin{tabularx}{\linewidth}{|l|c|c|X|}
\hline
\rowcolor{popblueL} \textbf{Liaison} & \textbf{Longueur (pm)} & \textbf{Énergie (kJ/mol)} & \textbf{Remarque} \\
\hline
\ce{C-H} & 109 & 415 & liaison courte et forte \\
\hline
\ce{C-C} & 154 & 348 & liaison plus longue, un peu moins forte \\
\hline
\end{tabularx}
\end{center}

\begin{exemplebox}[Lire les ordres de grandeur]
La liaison \ce{C-H} ($\ell = \SI{109}{pm} = \SI{0,109}{nm}$) est plus courte et plus forte que la liaison \ce{C-C} ($\ell = \SI{154}{pm}$, $E = \SI{348}{kJ.mol^{-1}}$). Ces valeurs élevées expliquent que les alcanes sont des molécules \textbf{stables}, peu réactives à froid : on les appelait autrefois « paraffines » (du latin \emph{parum affinis}, « peu d'affinité »). Il faut une flamme ou une étincelle pour déclencher leur combustion — c'est pourquoi le gaz butane ne s'enflamme pas tout seul dans sa bouteille !
\end{exemplebox}

\courssub{Généralisation : la chaîne carbonée en zigzag}

\textbf{Intuition.} Si chaque carbone est un tétraèdre avec des angles de $109^{\circ}28'$, alors une suite de carbones reliés par des liaisons simples ne peut \emph{pas} former une ligne droite : la chaîne carbonée se replie naturellement en \cle{zigzag}, comme un serpentin. C'est exactement ce que l'on observe dans le butane de nos bouteilles de gaz, \ce{CH3-CH2-CH2-CH3}.

\begin{propriete}[Squelette des alcanes]
Dans un alcane à chaîne linéaire de $n$ atomes de carbone, la chaîne carbonée adopte une forme en \textbf{zigzag} dans l'espace, car chaque carbone est tétragonal (angles voisins de $109^{\circ}28'$). Grâce à la libre rotation autour de chaque liaison \ce{C-C}, la chaîne peut en outre se replier et adopter de nombreuses conformations.
\end{propriete}

\begin{center}
\chemfig{CH_3-[::30]-[::-60]-[::60]CH_3}
\end{center}

Dans cette écriture, chaque segment représente une liaison \ce{C-C} et chaque sommet (ou extrémité) un atome de carbone tétragonal portant ses hydrogènes. Le butane n'est donc pas une molécule « droite », même si on l'écrit commodément \ce{CH3-CH2-CH2-CH3} sur une ligne.

\begin{exoresolu}[Compter les liaisons dans le butane]
\textbf{Énoncé.} Le butane a pour formule brute \ce{C4H10}. (1) Combien de liaisons \ce{C-C} et de liaisons \ce{C-H} contient une molécule de butane ? (2) Vérifier que chaque carbone est bien tétragonal.
\tcblower
\textbf{Solution.} (1) La chaîne de 4 carbones comporte $4 - 1 = 3$ liaisons \ce{C-C}. Chaque carbone formant 4 liaisons au total, le nombre de liaisons \ce{C-H} est égal au nombre d'atomes d'hydrogène, soit \textbf{10 liaisons \ce{C-H}}. Vérification par les doublets : 4 carbones $\times$ 4 liaisons $= 16$ « bras » ; les 3 liaisons \ce{C-C} consomment $2 \times 3 = 6$ bras ; il reste $16 - 6 = 10$ bras pour 10 hydrogènes. (2) Les deux carbones des extrémités (\ce{CH3}) forment 3 liaisons \ce{C-H} et 1 liaison \ce{C-C} : 4 liaisons simples. Les deux carbones centraux (\ce{CH2}) forment 2 liaisons \ce{C-H} et 2 liaisons \ce{C-C} : 4 liaisons simples. Chaque carbone est donc bien \textbf{tétragonal}, entouré de 4 doublets liants dirigés vers les sommets d'un tétraèdre.
\end{exoresolu}

\begin{attention}[Zigzag ne veut pas dire « angles droits »]
Quand tu dessines une chaîne carbonée en zigzag, ne trace jamais des angles de $90^{\circ}$ ou de $180^{\circ}$ entre deux liaisons \ce{C-C} : l'angle $\widehat{\ce{C-C-C}}$ vaut environ $109,5^{\circ}$. De même, écrire \ce{CH3-CH2-CH2-CH3} sur une ligne droite n'est qu'une \emph{commodité d'écriture} (formule semi-développée) : la molécule réelle est coudée dans l'espace.
\end{attention}

Cette tétravalence du carbone, source de la géométrie tétraédrique, conduit directement à la formule générale des alcanes : \textbf{$C_nH_{2n+2}$} (pour les chaînes non cycliques). Nous verrons dans la prochaine section comment cette formule permet de prévoir l'isomérie et de nommer systématiquement ces composés.

\begin{aretenir}[L'essentiel de la section]
\begin{itemize}
\item Le carbone ($Z = 6$) possède 4 électrons de valence et forme \textbf{4 liaisons covalentes simples} dans les alcanes : c'est un \cle{carbone tétragonal}.
\item Les 4 liaisons sont dirigées vers les sommets d'un \textbf{tétraèdre régulier} ; l'angle de liaison vaut $109^{\circ}28'$.
\item Le méthane \ce{CH4} est tétraédrique (et non plan !) ; on le représente en 3D par la convention de Cram.
\item Dans l'éthane \ce{C2H6}, la \textbf{libre rotation} autour de la liaison \ce{C-C} donne des \cle{conformations} (éclipsée, décalée), visibles en projection de Newman ; ce ne sont pas des isomères.
\item Liaisons : \ce{C-H} ($\ell \approx \SI{109}{pm}$, $E \approx \SI{415}{kJ.mol^{-1}}$) ; \ce{C-C} ($\ell \approx \SI{154}{pm}$, $E \approx \SI{348}{kJ.mol^{-1}}$).
\item La chaîne carbonée d'un alcane adopte une forme en \textbf{zigzag} dans l'espace.
\item Formule générale des alcanes acycliques : \textbf{$C_nH_{2n+2}$}.
\end{itemize}
\end{aretenir}

\coursec{Formule générale, formules développées et isomérie}

Dans la section précédente, nous avons découvert le \cle{carbone tétragonal} : chaque atome de carbone d'un alcane forme quatre liaisons covalentes simples dirigées vers les sommets d'un tétraèdre. Nous avons aussi vu que les molécules de méthane \ce{CH4} et d'éthane \ce{C2H6} peuvent tourner autour de la liaison C--C (conformations). Une question se pose alors naturellement : si l'on continue à ajouter des atomes de carbone les uns à la suite des autres, peut-on prévoir la formule de la molécule obtenue ? Et surtout, une même formule brute peut-elle correspondre à plusieurs molécules différentes ? C'est tout l'objet de cette section, qui te donnera les outils pour écrire et reconnaître \emph{tous} les alcanes contenant jusqu'à cinq atomes de carbone.

\courssub{La formule générale des alcanes acycliques : $C_nH_{2n+2}$}

\textbf{Intuition : compter les hydrogènes}\par

Prenons les premiers alcanes et comptons leurs atomes :

\begin{center}
\begin{tabularx}{\linewidth}{|l|c|c|X|}
\hline
\rowcolor{popblueL} \textbf{Alcane} & \textbf{Nombre de C ($n$)} & \textbf{Nombre de H} & \textbf{Formule brute} \\
\hline
Méthane & 1 & 4 & \ce{CH4} \\
\hline
Éthane & 2 & 6 & \ce{C2H6} \\
\hline
Propane & 3 & 8 & \ce{C3H8} \\
\hline
Butane & 4 & 10 & \ce{C4H10} \\
\hline
Pentane & 5 & 12 & \ce{C5H12} \\
\hline
\end{tabularx}
\end{center}

Observe la colonne des hydrogènes : 4, 6, 8, 10, 12... À chaque fois qu'on ajoute \emph{un} atome de carbone, on ajoute \emph{deux} atomes d'hydrogène. C'est une suite arithmétique de raison 2. Pour $n$ atomes de carbone, le nombre d'hydrogènes vaut donc $2n + 2$ (vérifie : pour $n = 1$, $2 \times 1 + 2 = 4$ ; pour $n = 5$, $2 \times 5 + 2 = 12$).

\textbf{Démonstration par la valence du carbone}\par

Ce résultat n'est pas un hasard : il découle directement du caractère tétragonal du carbone. Considérons une chaîne \emph{ouverte} (acyclique) de $n$ atomes de carbone :

\begin{itemize}
\item les $n$ carbones forment $n - 1$ liaisons C--C entre eux ;
\item chaque carbone forme au total 4 liaisons, soit $4n$ liaisons au départ de l'ensemble des carbones ;
\item les $n-1$ liaisons C--C « consomment » chacune 2 de ces liaisons (une sur chaque carbone), donc $2(n-1)$ liaisons ;
\item il reste donc $4n - 2(n-1) = 2n + 2$ liaisons disponibles pour les hydrogènes.
\end{itemize}

\begin{propriete}[Formule générale des alcanes acycliques]
Tout alcane à chaîne ouverte (acyclique), ramifiée ou non, possède la formule brute générale :
\[ C_nH_{2n+2} \qquad (n \geq 1) \]
Réciproquement, tout hydrocarbure de formule $C_nH_{2n+2}$ est un alcane acyclique : il ne contient que des liaisons simples C--C et C--H.
\end{propriete}

\begin{exemplebox}[Vérification avec le gaz domestique]
Le butane, gaz contenu dans les bouteilles de gaz domestique utilisées dans les cuisines camerounaises, possède $n = 4$ atomes de carbone. Sa formule brute est donc $C_4H_{2\times 4 + 2} = \ce{C4H10}$. De même, l'hexane ($n = 6$), solvant d'extraction de l'huile de palme, a pour formule \ce{C6H14}.
\end{exemplebox}

\begin{attention}[Piège classique]
La formule $C_nH_{2n+2}$ ne vaut que pour les alcanes \textbf{acycliques}. Dès que la molécule contient un cycle, on « perd » deux hydrogènes (voir ci-dessous). Ne conclus jamais qu'un composé est un alcane acyclique sans avoir vérifié l'absence de cycle !
\end{attention}

\courssub{Les cyclanes (cycloalcanes) : formule générale $C_nH_{2n}$}

Que se passe-t-il si les deux extrémités d'une chaîne carbonée se rejoignent pour former un \cle{cycle} ? Pour créer cette liaison C--C supplémentaire, il faut retirer un hydrogène sur chacun des deux carbones d'extrémité : la molécule perd donc \ce{H2} par rapport à l'alcane correspondant.

\begin{propriete}[Formule générale des cyclanes]
Un cycloalcane (monocyclique) possède la formule brute générale :
\[ C_nH_{2n} \qquad (n \geq 3) \]
Il faut au moins 3 atomes de carbone pour fermer un cycle.
\end{propriete}

Les quatre premiers cyclanes sont :

\begin{itemize}
\item le \cle{cyclopropane} \ce{C3H6} : cycle à 3 carbones, très « tendu » (angles de 60° au lieu de 109,5°), utilisé autrefois comme anesthésique ;
\item le \cle{cyclobutane} \ce{C4H8} : cycle à 4 carbones ;
\item le \cle{cyclopentane} \ce{C5H10} : cycle à 5 carbones ;
\item le \cle{cyclohexane} \ce{C6H12} : cycle à 6 carbones, le plus stable et le plus répandu (solvant industriel).
\end{itemize}

\begin{popfigure}
\begin{center}
\begin{tabular}{cccc}
\chemfig{*3(-CH_2-CH_2-CH_2-)} &
\chemfig{*4(-CH_2-CH_2-CH_2-CH_2-)} &
\chemfig{*5(-CH_2-CH_2-CH_2-CH_2-CH_2-)} &
\chemfig{*6(-CH_2-CH_2-CH_2-CH_2-CH_2-CH_2-)} \\[4pt]
Cyclopropane & Cyclobutane & Cyclopentane & Cyclohexane \\
\ce{C3H6} & \ce{C4H8} & \ce{C5H10} & \ce{C6H12} \\
\end{tabular}
\end{center}
\legende{Formules développées des quatre premiers cyclanes. Chaque sommet et chaque extrémité de trait représente un atome de carbone portant deux hydrogènes (non écrits explicitement dans cette représentation polygonale).}
\end{popfigure}

\begin{attention}[Ne pas confondre]
Un cycloalcane de formule $C_nH_{2n}$ a la \emph{même} formule brute qu'un alcène (hydrocarbure à double liaison, que tu étudieras plus tard). Par exemple \ce{C3H6} peut être le cyclopropane \emph{ou} le propène. La formule brute seule ne suffit donc pas à identifier une molécule : c'est toute la richesse de l'isomérie.
\end{attention}

\courssub{Formule brute, formule développée, formule semi-développée}

Pour décrire une molécule organique, le chimiste dispose de trois niveaux d'écriture, du plus compact au plus détaillé.

\begin{definition}[Les trois types de formules]
\begin{itemize}
\item La \cle{formule brute} indique la nature et le nombre des atomes de la molécule, sans préciser leur enchaînement. Exemple : \ce{C4H10}.
\item La \cle{formule développée plane} représente \emph{tous} les atomes et \emph{toutes} les liaisons de la molécule, disposés dans le plan de la feuille. Exemple pour le butane : \chemfig{H-C(-[2]H)(-[6]H)-C(-[2]H)(-[6]H)-C(-[2]H)(-[6]H)-C(-[2]H)(-[6]H)-H}.
\item La \cle{formule semi-développée} est un compromis : on écrit les groupes d'atomes liés à chaque carbone sans représenter les liaisons C--H, mais en conservant l'enchaînement des carbones. Exemple : \ce{CH3-CH2-CH2-CH3}.
\end{itemize}
\end{definition}

\begin{methode}[Passer de la formule développée à la semi-développée]
\begin{enumerate}
\item Repérer chaque atome de carbone dans la formule développée.
\item Regrouper avec lui les hydrogènes qui lui sont directement liés : un carbone terminal devient \ce{CH3}, un carbone intermédiaire devient \ce{CH2}, un carbone ramifié devient \ce{CH}.
\item Écrire les groupes obtenus dans l'ordre de la chaîne, séparés par des tirets représentant les liaisons C--C.
\item Vérifier que le compte total des atomes redonne bien la formule brute.
\end{enumerate}
\end{methode}

\begin{exemplebox}[Application au propane]
Formule développée du propane : \chemfig{H-C(-[2]H)(-[6]H)-C(-[2]H)(-[6]H)-C(-[2]H)(-[6]H)-H}. Le premier carbone porte 3 H : groupe \ce{CH3}. Le deuxième porte 2 H : groupe \ce{CH2}. Le troisième porte 3 H : groupe \ce{CH3}. Formule semi-développée : \ce{CH3-CH2-CH3}. Vérification : $3 + 2 + 3 = 8$ hydrogènes, soit \ce{C3H8}. C'est correct.
\end{exemplebox}

\courssub{L'isomérie de constitution : une même formule brute, plusieurs molécules}

\textbf{Intuition : le jeu de construction}\par

Imagine que tu disposes de 4 boules noires (les carbones) et de 10 boules blanches (les hydrogènes) d'un modèle moléculaire. Tu peux aligner les 4 boules noires en file indienne... mais tu peux aussi en aligner 3 et fixer la quatrième sur la boule du milieu ! Dans les deux cas, tu as utilisé exactement les mêmes boules, donc la même formule brute \ce{C4H10} ; pourtant les deux édifices sont différents, et ils ne peuvent pas se superposer, même en tournant les liaisons. Tu viens de découvrir l'\cle{isomérie}.

\begin{definition}[Isomères de constitution]
Des \cle{isomères de constitution} sont des composés qui possèdent la \textbf{même formule brute} mais des \textbf{formules développées différentes}, c'est-à-dire un enchaînement d'atomes différent. Ils ont des propriétés physiques et chimiques distinctes (températures d'ébullition différentes, par exemple).
\end{definition}

\textbf{Les deux isomères du butane \ce{C4H10}}\par

\begin{itemize}
\item le \cle{butane} (chaîne linéaire) : \ce{CH3-CH2-CH2-CH3} ;
\item le \cle{2-méthylpropane} (chaîne ramifiée) : \ce{CH3-CH(CH3)-CH3}, c'est-à-dire \chemfig{CH_3-CH(-[2]CH_3)-CH_3}.
\end{itemize}

Ces deux gaz ont des températures d'ébullition différentes : \SI{-0,5}{\celsius} pour le butane, \SI{-11,7}{\celsius} pour le 2-méthylpropane. Ce sont bien deux substances distinctes.

\begin{attention}[Erreur très fréquente au Probatoire]
Beaucoup d'élèves croient que \ce{CH3-CH(CH3)-CH3} et \ce{CH3-CH2-CH2-CH3} sont deux écritures de la même molécule « parce qu'elles ont la même formule brute ». C'est faux : ce sont \textbf{deux isomères distincts}. Dans le premier, le carbone central est lié à \emph{trois} autres carbones ; dans le second, aucun carbone n'est lié à plus de deux carbones. L'enchaînement des atomes diffère, donc les molécules diffèrent. À l'inverse, deux dessins qui semblent différents mais qui se superposent par rotation autour des liaisons C--C représentent la \emph{même} molécule.
\end{attention}

\courssub{Les trois isomères du pentane \ce{C5H12}}

Appliquons une démarche systématique pour trouver \emph{tous} les isomères du pentane, sans en oublier et sans en compter deux fois.

\begin{methode}[Trouver tous les isomères d'un alcane]
\begin{enumerate}
\item Écrire d'abord l'isomère à \textbf{chaîne linéaire} (tous les carbones bout à bout).
\item \textbf{Raccourcir} la chaîne principale d'un carbone et placer ce carbone comme ramification (groupe méthyle \ce{CH3}) sur chaque position possible de la chaîne, en évitant les extrémités (cela rallongerait la chaîne !) et les doublons par symétrie.
\item Raccourcir encore la chaîne et placer \textbf{deux ramifications}, en variant leurs positions.
\item Pour chaque formule obtenue, vérifier que chaque carbone forme bien 4 liaisons et que la formule brute est respectée.
\item S'assurer qu'aucune formule n'est superposable à une autre par rotation ou retournement.
\end{enumerate}
\end{methode}

Appliquons cette méthode à \ce{C5H12} :

\begin{itemize}
\item \textbf{Étape 1} : chaîne de 5 carbones $\Rightarrow$ le \cle{pentane} : \ce{CH3-CH2-CH2-CH2-CH3}.
\item \textbf{Étape 2} : chaîne de 4 carbones + 1 méthyle. Le méthyle ne peut se placer qu'en position 2 (la position 3 est équivalente par symétrie) $\Rightarrow$ le \cle{2-méthylbutane} : \ce{CH3-CH(CH3)-CH2-CH3}.
\item \textbf{Étape 3} : chaîne de 3 carbones + 2 méthyles, forcément sur le carbone central $\Rightarrow$ le \cle{2,2-diméthylpropane} : \ce{C(CH3)4}, soit \ce{CH3-C(CH3)2-CH3}.
\end{itemize}

Il n'existe aucune autre possibilité : \ce{C5H12} possède \textbf{exactement 3 isomères}.

\begin{popfigure}
\begin{center}
\begin{tabular}{ccc}
\chemfig{CH_3-CH_2-CH_2-CH_2-CH_3} &
\chemfig{CH_3-CH(-[2]CH_3)-CH_2-CH_3} &
\chemfig{CH_3-C(-[2]CH_3)(-[6]CH_3)-CH_3} \\[6pt]
Pentane & 2-méthylbutane & 2,2-diméthylpropane \\
\multicolumn{3}{c}{\textbf{Formule brute commune : \ce{C5H12}}}
\end{tabular}
\end{center}
\legende{Formules semi-développées des trois isomères de constitution du pentane \ce{C5H12}. Remarque que plus la molécule est ramifiée, plus elle est compacte, et plus sa température d'ébullition est basse (\SI{36,1}{\celsius} ; \SI{27,8}{\celsius} ; \SI{9,5}{\celsius} respectivement).}
\end{popfigure}

\begin{exemplebox}[Comptage des isomères pour $n \leq 5$]
\ce{C4H10} possède exactement \textbf{2 isomères} (butane et 2-méthylpropane) ; \ce{C5H12} en possède exactement \textbf{3}. Pour $n \leq 3$, il n'y a qu'une seule structure possible : le méthane, l'éthane et le propane n'ont pas d'isomère.
\end{exemplebox}

\begin{popfigure}
\begin{center}
\begin{tabularx}{\linewidth}{|c|c|c|X|}
\hline
\rowcolor{popblueL} \textbf{Formule brute} & \textbf{$n$} & \textbf{Nombre d'isomères} & \textbf{Exemples de structures} \\
\hline
\ce{CH4} & 1 & 1 & méthane \\
\hline
\ce{C2H6} & 2 & 1 & éthane \\
\hline
\ce{C3H8} & 3 & 1 & propane \\
\hline
\ce{C4H10} & 4 & 2 & butane, 2-méthylpropane \\
\hline
\ce{C5H12} & 5 & 3 & pentane, 2-méthylbutane, 2,2-diméthylpropane \\
\hline
\end{tabularx}
\end{center}
\legende{Nombre d'isomères de constitution des alcanes pour $n$ allant de 1 à 5. Ce nombre n'augmente qu'à partir de $n = 4$, puis il explose littéralement.}
\end{popfigure}

\begin{savaistu}[Une croissance explosive]
Le nombre d'isomères croît extraordinairement vite avec le nombre de carbones : \ce{C6H14} en possède 5, \ce{C7H16} en possède 9, \ce{C10H22} en possède déjà \textbf{75}, \ce{C15H32} plus de 4 000, et \ce{C20H42} plus de 366 000 ! C'est cette prolifération qui explique l'immense diversité de la chimie organique : le pétrole traité à la SONARA de Limbé contient des milliers d'hydrocarbures différents, que le raffinage sépare en fractions (essence, kérosène, gasoil...) selon leurs températures d'ébullition.
\end{savaistu}

\courssub{Activité de manipulation : modèles moléculaires des alcanes ($n \leq 5$)}

\begin{experience}[Construction des modèles moléculaires]
\textbf{Matériel} : boîte de modèles moléculaires (boules noires tétravalentes pour le carbone, boules blanches monovalentes pour l'hydrogène, tiges pour les liaisons simples).

\textbf{Protocole} :
\begin{enumerate}
\item Construire successivement les modèles de \ce{CH4}, \ce{C2H6}, \ce{C3H8} : vérifier qu'il n'existe qu'\emph{une} façon de les assembler.
\item Construire le modèle du butane \ce{CH3-CH2-CH2-CH3}, puis le transformer en 2-méthylpropane en déplaçant un carbone d'extrémité sur le carbone central : constater que les deux modèles ne sont pas superposables.
\item Construire les trois modèles de \ce{C5H12} et vérifier qu'aucun quatrième assemblage non superposable n'est possible.
\item Pour chaque modèle, faire tourner les groupes \ce{CH3} autour des liaisons C--C : constater que la molécule reste la même (conformations).
\end{enumerate}

\textbf{Observations} : pour $n \leq 3$, un seul assemblage possible ; pour $n = 4$, deux assemblages ; pour $n = 5$, trois assemblages. Toute tentative de construire un quatrième isomère du pentane redonne, par rotation ou retournement, l'un des trois déjà obtenus.

\textbf{Interprétation} : l'isomérie de constitution provient de l'ordre d'enchaînement des atomes de carbone, et non de l'orientation spatiale des liaisons, qui est libre autour des liaisons simples.
\end{experience}

\begin{exoresolu}[Dénombrer les isomères de l'hexane]
Énoncé. Un élève affirme : « L'hexane \ce{C6H14} possède 4 isomères. » En appliquant la méthode de raccourcissement progressif de la chaîne, déterminer le nombre exact d'isomères de \ce{C6H14} et écrire leurs formules semi-développées.
\tcblower
Solution détaillée. Appliquons la méthode :

\textbf{Chaîne de 6 carbones} : un seul isomère linéaire, l'hexane :
\[ \ce{CH3-CH2-CH2-CH2-CH2-CH3} \]

\textbf{Chaîne de 5 carbones + 1 méthyle} : le méthyle peut se placer en position 2 ou en position 3 (les positions 4 et 5 sont équivalentes aux positions 2 et 1 par symétrie ; la position 1 est interdite car elle redonnerait l'hexane) :
\[ \ce{CH3-CH(CH3)-CH2-CH2-CH3} \quad \text{(2-méthylpentane)} \]
\[ \ce{CH3-CH2-CH(CH3)-CH2-CH3} \quad \text{(3-méthylpentane)} \]

\textbf{Chaîne de 4 carbones + 2 méthyles} : les deux méthyles peuvent être sur le même carbone (position 2) ou sur deux carbones différents (positions 2 et 3) :
\[ \ce{CH3-C(CH3)2-CH2-CH3} \quad \text{(2,2-diméthylbutane)} \]
\[ \ce{CH3-CH(CH3)-CH(CH3)-CH3} \quad \text{(2,3-diméthylbutane)} \]

\textbf{Chaîne de 3 carbones} : il faudrait placer 3 méthyles sur 3 carbones, mais le carbone central porterait alors 5 liaisons : impossible.

Conclusion : \ce{C6H14} possède exactement \textbf{5 isomères}, et non 4. L'élève avait probablement oublié le 2,3-diméthylbutane, l'oubli le plus fréquent. Vérification finale : chaque formule compte bien 6 carbones et $2 \times 6 + 2 = 14$ hydrogènes.
\end{exoresolu}

\begin{exercice}[À toi de jouer]
\begin{enumerate}
\item Un hydrocarbure gazeux a pour formule brute \ce{C5H12}. Justifier, à l'aide de la formule générale, qu'il s'agit d'un alcane acyclique.
\item Écrire les formules semi-développées de tous ses isomères et donner le nom de chacun.
\item Un autre hydrocarbure a pour formule \ce{C5H10}. Peut-il s'agir d'un alcane acyclique ? Proposer une explication.
\end{enumerate}
\end{exercice}

\begin{corrige}[Exercice]
\begin{enumerate}
\item Pour $n = 5$, la formule générale des alcanes donne $2n + 2 = 12$ hydrogènes : \ce{C5H12} correspond exactement à $C_nH_{2n+2}$. C'est donc un alcane acyclique.
\item Les trois isomères sont : \ce{CH3-CH2-CH2-CH2-CH3} (pentane) ; \ce{CH3-CH(CH3)-CH2-CH3} (2-méthylbutane) ; \ce{CH3-C(CH3)2-CH3} (2,2-diméthylpropane).
\item \ce{C5H10} possède deux hydrogènes de moins que \ce{C5H12} : ce n'est pas un alcane acyclique. Il peut s'agir d'un cycloalcane ($C_nH_{2n}$, par exemple le cyclopentane) ou d'un alcène (hydrocarbure à double liaison).
\end{enumerate}
\end{corrige}

\begin{aretenir}[L'essentiel de la section]
\begin{itemize}
\item Les alcanes acycliques ont pour formule générale $C_nH_{2n+2}$ ; les cycloalcanes ont pour formule générale $C_nH_{2n}$ ($n \geq 3$).
\item Trois niveaux d'écriture : formule brute (composition), formule développée (tous les atomes et toutes les liaisons), formule semi-développée (enchaînement des groupes \ce{CH3}, \ce{CH2}, \ce{CH}).
\item Des isomères de constitution ont la même formule brute mais des enchaînements d'atomes différents : ce sont des substances différentes, aux propriétés différentes.
\item \ce{C4H10} possède 2 isomères ; \ce{C5H12} en possède 3 ; \ce{C6H14} en possède 5.
\item Pour dénombrer les isomères sans erreur : raccourcir progressivement la chaîne principale et déplacer méthodiquement les ramifications, en éliminant les doublons par symétrie.
\end{itemize}
\end{aretenir}

Maintenant que tu sais écrire toutes les structures possibles d'un alcane, il reste à apprendre à les \emph{nommer} sans ambiguïté : c'est l'objet de la section suivante, consacrée à la nomenclature officielle des alcanes.

\coursec{Nomenclature des alcanes}

Dans la section précédente, tu as découvert que plusieurs molécules différentes peuvent partager la même formule brute : ce sont les \cle{isomères}. Le pentane $\ce{C5H12}$, par exemple, existe sous trois formes. Comment alors désigner chacune d'elles sans ambiguïté ? Imagine un marché de Douala où trois vendeuses portent le même prénom : pour les distinguer, on ajoute des précisions. La chimie procède exactement ainsi : la \cle{nomenclature} est l'art de donner à chaque molécule un nom unique, construit selon des règles précises. C'est ce langage universel que nous allons apprendre dans cette section.

\courssub{Les noms des alcanes linéaires : racine et suffixe}

Commençons par le cas le plus simple : les alcanes à chaîne \emph{linéaire} (non ramifiée). Leur nom se construit toujours de la même façon :

\begin{center}
\textbf{racine} (qui indique le nombre d'atomes de carbone) \textbf{+ suffixe -ane} (qui indique la famille des alcanes)
\end{center}

\begin{definition}[Racines et suffixe]
Le nom d'un alcane linéaire est formé d'une \cle{racine} grecque ou latine indiquant le nombre $n$ d'atomes de carbone de la chaîne, suivie du suffixe \cle{-ane}, signature de la famille des alcanes (formule générale $C_nH_{2n+2}$).
\end{definition}

Voici les racines à connaître par cœur — elles reviendront dans toutes les familles de composés organiques :

\begin{center}
\begin{tabularx}{\linewidth}{|c|l|l|X|}
\hline
\rowcolor{popblueL} \textbf{Nombre de C} & \textbf{Racine} & \textbf{Nom de l'alcane} & \textbf{Formule brute} \\
\hline
1 & méth- & méthane & \ce{CH4} \\
\hline
2 & éth- & éthane & \ce{C2H6} \\
\hline
3 & prop- & propane & \ce{C3H8} \\
\hline
4 & but- & butane & \ce{C4H10} \\
\hline
5 & pent- & pentane & \ce{C5H12} \\
\hline
6 & hex- & hexane & \ce{C6H14} \\
\hline
7 & hept- & heptane & \ce{C7H16} \\
\hline
8 & oct- & octane & \ce{C8H18} \\
\hline
\end{tabularx}
\end{center}

\begin{aretenir}[À mémoriser impérativement]
Les quatre premières racines (\textbf{méth-}, \textbf{éth-}, \textbf{prop-}, \textbf{but-}) sont historiques et ne suivent pas la numération grecque. À partir de 5 carbones, on retrouve les préfixes grecs classiques : \textbf{pent-} (5), \textbf{hex-} (6), \textbf{hept-} (7), \textbf{oct-} (8), comme dans \emph{pentagone} ou \emph{octogone}. Le gaz butane des bouteilles de gaz domestiques utilisées dans nos cuisines doit son nom à ses 4 atomes de carbone.
\end{aretenir}

\courssub{La chaîne principale et les groupes alkyles}

\courssub{La chaîne principale : la plus longue chaîne carbonée}

Quand un alcane est ramifié, sa structure ressemble à un arbre : un tronc et des branches. Pour le nommer, il faut d'abord identifier le « tronc ».

\begin{definition}[Chaîne principale]
La \cle{chaîne principale} d'un alcane ramifié est la \textbf{plus longue chaîne carbonée continue} (non interrompue) que l'on peut tracer dans la molécule. Elle n'est pas nécessairement écrite horizontalement : elle peut « tourner » dans la formule développée. C'est elle qui donne la racine du nom.
\end{definition}

\begin{attention}[Piège classique]
La chaîne principale n'est pas toujours la chaîne écrite en ligne droite sur ta copie ! Il faut explorer tous les chemins possibles d'un bout à l'autre de la molécule et retenir le plus long. Une chaîne qui « monte » puis « redescend » est parfaitement valable, car les liaisons C--C peuvent tourner librement.
\end{attention}

\courssub{Les groupes alkyles : les ramifications}

Les « branches » greffées sur la chaîne principale s'appellent des \cle{groupes alkyles}.

\begin{definition}[Groupe alkyle]
Un \cle{groupe alkyle} est un fragment de molécule obtenu en retirant un atome d'hydrogène à un alcane. Sa formule générale est $C_nH_{2n+1}-$. Son nom se forme en remplaçant le suffixe \textbf{-ane} de l'alcane par le suffixe \textbf{-yle} : méth\textbf{ane} $\rightarrow$ méth\textbf{yle}, éth\textbf{ane} $\rightarrow$ éth\textbf{yle}.
\end{definition}

\begin{popfigure}
\begin{center}
\begin{tabularx}{\linewidth}{|c|l|l|X|}
\hline
\rowcolor{popgreenL} \textbf{Nombre de C} & \textbf{Racine} & \textbf{Groupe alkyle} & \textbf{Formule semi-développée} \\
\hline
1 & méth- & méthyle & \ce{CH3-} \\
\hline
2 & éth- & éthyle & \ce{CH3-CH2-} ou \ce{C2H5-} \\
\hline
3 & prop- & propyle & \ce{CH3-CH2-CH2-} \\
\hline
4 & but- & butyle & \ce{CH3-CH2-CH2-CH2-} \\
\hline
\end{tabularx}
\end{center}
\legende{Les groupes alkyles les plus courants : on passe du nom de l'alcane au nom du groupe en changeant -ane en -yle.}
\end{popfigure}

\courssub{Les trois règles de la nomenclature}

\begin{methode}[Nommer un alcane ramifié en 3 étapes]
\begin{enumerate}
\item \textbf{Identifier la chaîne principale} : repérer la plus longue chaîne carbonée continue. Elle donne la racine du nom (méth-, éth-, prop-, but-, pent-\ldots) suivie du suffixe -ane.
\item \textbf{Numéroter la chaîne principale} : compter les carbones à partir de l'extrémité la plus proche d'une ramification, de façon à donner aux ramifications les \textbf{indices (numéros) les plus petits possibles} (règle du « premier point de différence »).
\item \textbf{Écrire le nom complet} : citer les ramifications par \textbf{ordre alphabétique} (en ignorant les préfixes di-, tri-, tétra-), chacune précédée de son indice. Si plusieurs groupes identiques sont présents, utiliser les préfixes \textbf{di-}, \textbf{tri-}, \textbf{tétra-} et répéter les indices. Les chiffres sont séparés des lettres par des \textbf{tirets} et entre eux par des \textbf{virgules}.
\end{enumerate}
\end{methode}

Appliquons immédiatement ces règles sur un exemple complet, pas à pas.

\begin{exemplebox}[Nommer le 2,3-diméthylpentane pas à pas]
Considérons la molécule de formule semi-développée :
\[ \ce{CH3-CH(CH3)-CH(CH3)-CH2-CH3} \]
\textbf{Étape 1 — Chaîne principale.} Le plus long chemin continu contient \textbf{5 carbones} : la racine est donc \textbf{pent-}, le nom de base \textbf{pentane}. Il reste deux branches \ce{CH3}, c'est-à-dire deux groupes \textbf{méthyle}.

\textbf{Étape 2 — Numérotation.} En numérotant de gauche à droite, les ramifications portent les indices 2 et 3. De droite à gauche, elles porteraient les indices 3 et 4. On retient donc la numérotation de gauche à droite : $2 < 3$ au premier point de différence.

\textbf{Étape 3 — Nom complet.} Deux groupes méthyle identiques : préfixe \textbf{di-}, indices répétés séparés par une virgule. Le nom est : \textbf{2,3-diméthylpentane}.
\end{exemplebox}

\begin{popfigure}
\begin{center}
\chemfig{CH_3-[0]CH(-[2]CH_3)-[0]CH(-[6]CH_3)-[0]CH_2-[0]CH_3}
\vspace{0.4em}

\begin{tikzpicture}[scale=1.0, every node/.style={font=\small}]
\foreach \x/\n in {0/1, 1.6/2, 3.2/3, 4.8/4, 6.4/5}{
  \fill[popblue] (\x,0) circle (2.2pt);
  \node[below=3pt, popblue] at (\x,0) {\n};
}
\draw[popblue, thick] (0,0) -- (6.4,0);
\fill[poporange] (1.6,0.7) circle (2.2pt);
\fill[poporange] (3.2,-0.7) circle (2.2pt);
\draw[poporange, thick] (1.6,0) -- (1.6,0.7);
\draw[poporange, thick] (3.2,0) -- (3.2,-0.7);
\node[above=2pt, poporange] at (1.6,0.7) {\ce{CH3}};
\node[below=2pt, poporange] at (3.2,-0.7) {\ce{CH3}};
\node[popdark, anchor=west] at (7.0,0) {chaîne principale : 5 C};
\end{tikzpicture}
\end{center}
\legende{Le 2,3-diméthylpentane : la chaîne principale (en bleu) compte 5 carbones numérotés de gauche à droite pour donner les indices 2 et 3 ; les deux groupes méthyle sont en orange.}
\end{popfigure}

\begin{attention}[Erreur fréquente n°1 : numéroter dans le mauvais sens]
Pour la molécule \ce{CH3-CH(CH3)-CH2-CH2-CH3}, deux numérotations sont possibles :
\begin{itemize}
\item de gauche à droite : le méthyle est en position \textbf{2} ;
\item de droite à gauche : le méthyle est en position 4.
\end{itemize}
Le bon nom est \textbf{2-méthylpentane}, et non 4-méthylpentane. Règle d'or : \textbf{toujours choisir la numérotation qui donne les plus petits indices} (règle du premier point de différence). En cas d'égalité sur le premier indice (ex. 2,4 contre 3,5), on compare les indices suivants : $2,4 < 3,5$.
\end{attention}

\begin{attention}[Erreur fréquente n°2 : la ponctuation du nom]
La typographie du nom fait partie de la règle et est sanctionnée au Probatoire :
\begin{itemize}
\item on écrit \textbf{2-méthylpropane} et non « 2-méthyl propane » (pas d'espace) ;
\item les \textbf{chiffres sont séparés des lettres par des tirets} : 2,3-diméthylpentane ;
\item les \textbf{chiffres entre eux sont séparés par des virgules} : 2,2,3-triméthylbutane ;
\item le nom s'écrit en un seul mot, sans majuscule.
\end{itemize}
\end{attention}

\begin{exemplebox}[Attention à la vraie longueur de la chaîne principale]
Nommons la molécule de formule semi-développée :
\[ \ce{CH3-CH(CH3)-CH(C2H5)-CH3} \]
Un élève pressé voit une chaîne horizontale de 4 carbones et propose « 3-éthyl-2-méthylbutane ». \textbf{C'est faux !} Le groupe éthyle \ce{-CH2-CH3} prolonge la chaîne : en passant par lui, on trace une chaîne continue de \textbf{5 carbones}. La chaîne principale est donc un \textbf{pentane}. Le groupe éthyle n'est plus une ramification : il fait partie de la chaîne. Les deux ramifications restantes sont des groupes \textbf{méthyle} (le \ce{CH3} initial sur le C2 et le \ce{CH3} terminal de l'ancienne chaîne butane).

\textbf{Numérotation :} De gauche à droite (par l'extrémité du \ce{CH3} initial), les méthyles sont en 2 et 3. De droite à gauche, ils seraient en 3 et 4. On retient 2 et 3.

\textbf{Nom final :} \textbf{2,3-diméthylpentane}. (Formule brute \ce{C7H16}, isomère de l'heptane).
\end{exemplebox}

\courssub{Exercices résolus d'application}

\begin{exoresolu}[Nommer un alcane ramifié]
Énoncé. Donner le nom de l'alcane de formule semi-développée :
\[ \ce{CH3-C(CH3)2-CH2-CH3} \]
\tcblower
Solution détaillée.

\textbf{Étape 1 — Chaîne principale.} Le plus long chemin continu passe par le carbone central et l'un quelconque des deux méthyles du haut : il contient \textbf{4 carbones}. Racine : \textbf{but-}, nom de base : \textbf{butane}.

\textbf{Étape 2 — Numérotation.} De gauche à droite, les deux groupes méthyle sont sur le carbone n°2 : indices 2 et 2. De droite à gauche, ils seraient en 3 et 3. On retient donc la numérotation de gauche à droite.

\textbf{Étape 3 — Nom complet.} Deux méthyles sur le même carbone : préfixe \textbf{di-} et indice répété deux fois, séparé par une virgule.

Le nom est : \textbf{2,2-diméthylbutane}. Cette molécule, de formule brute \ce{C6H14}, est un isomère de l'hexane.
\end{exoresolu}

\begin{exoresolu}[Passer du nom à la formule]
Énoncé. Écrire la formule semi-développée du \textbf{3-éthyl-2-méthylpentane}, puis donner sa formule brute.
\tcblower
Solution détaillée. On procède à l'envers, en décodant le nom :

\textbf{Étape 1.} « pentane » $\Rightarrow$ la chaîne principale compte 5 carbones. On écrit le squelette numéroté :
\[ \ce{C1-C2-C3-C4-C5} \]

\textbf{Étape 2.} « 2-méthyl » $\Rightarrow$ un groupe \ce{CH3} sur le carbone n°2. « 3-éthyl » $\Rightarrow$ un groupe \ce{C2H5} sur le carbone n°3.

\textbf{Étape 3.} On complète chaque carbone avec des hydrogènes pour qu'il forme 4 liaisons :
\[ \ce{CH3-CH(CH3)-CH(C2H5)-CH2-CH3} \]

\textbf{Vérification.} Comptons les atomes : chaîne de 5 C + 1 C (méthyle) + 2 C (éthyle) = 8 carbones. Formule brute : \ce{C8H18}, qui vérifie bien $C_nH_{2n+2}$ avec $n = 8$ ($2 \times 8 + 2 = 18$). La formule est cohérente.
\end{exoresolu}

\begin{exercice}[À toi de jouer]
\begin{enumerate}
\item Nomme les alcanes suivants : \textbf{a)} \ce{CH3-CH(CH3)-CH3} ; \textbf{b)} \ce{CH3-CH2-CH(CH3)-CH2-CH3} ; \textbf{c)} \ce{CH3-CH(CH3)-CH2-CH(C2H5)-CH3}.
\item Écris la formule semi-développée du 2,2,4-triméthylpentane (un composant de l'essence, référence de l'indice d'octane).
\end{enumerate}
\end{exercice}

\begin{corrige}[Exercice : À toi de jouer]
\textbf{1. a)} Chaîne de 3 C, un méthyle en position 2 : \textbf{2-méthylpropane} (isomère du butane).
\textbf{b)} Chaîne de 5 C, méthyle en 3 (même indice dans les deux sens) : \textbf{3-méthylpentane}.
\textbf{c)} La chaîne la plus longue passe par le groupe éthyle : 6 C $\Rightarrow$ \textbf{hexane}. Numérotation de gauche à droite : indices 2 (méthyle) et 4 (méthyle, l'ancien \ce{CH3} terminal devient ramification). Le nom est \textbf{2,4-diméthylhexane} (ordre alphabétique : un seul type de ramification).

\textbf{2.} Chaîne de 5 C ; deux méthyles en 2, un méthyle en 4 :
\[ \ce{CH3-C(CH3)2-CH2-CH(CH3)-CH3} \]
Formule brute : \ce{C8H18}.
\end{corrige}

\courssub{Nomenclature des cyclanes}

Nous avons vu dans la section précédente que les atomes de carbone peuvent aussi former des \cle{cycles}. La nomenclature des cyclanes est remarquablement simple.

\begin{definition}[Nom d'un cyclane]
Un \cle{cyclane} (ou cycloalcane) de formule générale $C_nH_{2n}$ se nomme en ajoutant le préfixe \cle{cyclo-} devant le nom de l'alcane linéaire possédant le même nombre d'atomes de carbone : cyclopropane ($\ce{C3H6}$, triangle), cyclobutane ($\ce{C4H8}$, carré), cyclopentane ($\ce{C5H10}$, pentagone), cyclohexane ($\ce{C6H12}$, hexagone).
\end{definition}

\begin{exemplebox}[Cyclanes ramifiés]
Si le cycle porte une seule ramification, aucun indice n'est nécessaire : le carbone substitué est le carbone 1 par définition. Ainsi, un cycle à 6 carbones portant un méthyle se nomme le \textbf{méthylcyclohexane}. S'il porte deux ramifications, on numérote le cycle de façon à donner les plus petits indices : \textbf{1,2-diméthylcyclopentane}, par exemple.
\end{exemplebox}

\begin{savaistu}[Un langage universel : l'IUPAC]
Les règles que tu viens d'apprendre ont été établies par l'\textbf{IUPAC} (Union internationale de chimie pure et appliquée), fondée en 1919. Grâce à elles, un chimiste camerounais de l'Université de Yaoundé I, un chimiste japonais à Tokyo et une chimiste brésilienne à São Paulo nomment les molécules \emph{exactement de la même façon}. Le nom « 2,2,4-triméthylpentane » est compris partout sur la planète sans traduction : c'est le même langage que celui des mathématiques. Cette universalité est indispensable : le catalogue mondial recense des dizaines de millions de molécules organiques, et chacune possède un nom unique. Les anciens noms « courants » (comme \emph{isobutane} pour le 2-méthylpropane) survivent parfois dans l'industrie, mais au Probatoire, seule la nomenclature systématique IUPAC est acceptée.
\end{savaistu}

\begin{aretenir}[L'essentiel de la section]
\begin{itemize}
\item Alcane linéaire : \textbf{racine + -ane} (méthane, éthane, propane, butane, pentane, hexane\ldots).
\item Groupe alkyle $C_nH_{2n+1}-$ : suffixe \textbf{-yle} (méthyle, éthyle, propyle).
\item Trois étapes : \textbf{chaîne principale} la plus longue $\rightarrow$ \textbf{numérotation} donnant les plus petits indices (premier point de différence) $\rightarrow$ \textbf{nom} avec ramifications par ordre alphabétique, préfixes di-, tri-, tétra-.
\item Ponctuation : tirets entre chiffres et lettres, virgules entre chiffres ; un seul mot.
\item Cyclanes : préfixe \textbf{cyclo-} devant le nom de l'alcane à $n$ carbones.
\end{itemize}
\end{aretenir}

Tu sais désormais nommer n'importe quel alcane et décoder n'importe quel nom. Dans la section suivante, nous quitterons le papier pour le laboratoire et le quotidien : nous étudierons les \emph{propriétés physiques} des alcanes — et tu comprendras pourquoi le gaz butane se liquéfie si facilement dans les bouteilles de gaz.

\coursec{Propriétés physiques des alcanes}

Dans la section précédente, tu as appris à nommer rigoureusement n'importe quel alcane : méthane, butane, 2-méthylpentane\ldots{} Mais à quoi ressemblent réellement ces composés ? Sont-ils gazeux, liquides ou solides ? Se mélangent-ils à l'eau ? Ces questions ne sont pas anodines : c'est précisément parce que le butane est un gaz facilement liquéfiable que ta famille peut l'acheter en bouteille pour la cuisine, et c'est parce que les alcanes ne se mélangent pas à l'eau qu'une marée noire forme une nappe à la surface de l'océan. Cette section relie la nomenclature que tu viens d'acquérir aux propriétés concrètes qui expliquent les usages quotidiens des alcanes.

\courssub{État physique à température ordinaire}

\textbf{Intuition d'abord.} Les molécules d'alcanes s'attirent les unes les autres par de faibles forces dites de \cle{Van der Waals} (interactions de London). Plus la molécule est grosse (plus $n$ est grand), plus sa surface est grande, plus ces forces s'additionnent, et plus il faut d'énergie pour séparer les molécules. Les petites molécules s'échappent facilement : ce sont des gaz. Les très grosses restent collées : ce sont des solides.

\begin{propriete}[État physique des alcanes à \SI{25}{\celsius}]
À température et pression ordinaires, l'état physique d'un alcane dépend du nombre $n$ d'atomes de carbone :
\begin{itemize}
\item $1 \leq n \leq 4$ : les alcanes sont \cle{gazeux} (méthane, éthane, propane, butane) ;
\item $5 \leq n \leq 17$ : les alcanes sont \cle{liquides} (pentane, hexane, heptane\ldots) ;
\item $n \geq 18$ : les alcanes sont \cle{solides} (paraffines, cires).
\end{itemize}
\end{propriete}

\begin{exemplebox}[Quelques alcanes familiers]
Le \cle{méthane} ($n=1$) est le principal constituant du gaz naturel et du biogaz produit par la fermentation des déchets organiques (marais, biodigesteurs). Le \cle{butane} ($n=4$) est le gaz des bouteilles domestiques. L'\cle{octane} ($n=8$) est un constituant de l'essence : c'est un liquide. La \cle{paraffine} des bougies et la vaseline sont des mélanges d'alcanes solides à longue chaîne ($n \geq 18$).
\end{exemplebox}

\begin{savaistu}[Pourquoi le butane est-il vendu en bouteille ?]
Le butane bout à \SI{-0,5}{\celsius} : à température ordinaire, il suffit d'une pression modérée (environ \SI{2}{bar}) pour le liquéfier. Liquide, il occupe environ 250 fois moins de volume qu'à l'état gazeux : on stocke donc une grande quantité d'énergie dans une petite bouteille ! Quand on ouvre le robinet, la pression chute, le butane se vaporise et alimente la cuisinière. Le propane ($T_{eb} = \SI{-42}{\celsius}$), plus difficile à liquéfier, demande des bouteilles plus résistantes ; le méthane, lui, ne peut pratiquement pas être liquéfié à température ambiante, d'où son transport par gazoduc.
\end{savaistu}

\courssub{Températures de fusion et d'ébullition}

\begin{propriete}[Évolution des températures de changement d'état]
Pour les alcanes \textbf{linéaires}, les températures de fusion et d'ébullition \cle{augmentent régulièrement} avec le nombre $n$ d'atomes de carbone. À nombre de carbone égal, plus la molécule est \cle{ramifiée} (compacte, presque sphérique), plus sa température d'ébullition est \textbf{faible}, car la surface de contact entre molécules diminue.
\end{propriete}

\begin{exemplebox}[Valeurs repères à connaître]
\begin{itemize}
\item Méthane \ce{CH4} : $T_{eb} = \SI{-162}{\celsius}$ (gaz très volatil) ;
\item Butane \ce{C4H10} : $T_{eb} = \SI{-0,5}{\celsius}$ (gaz à température ordinaire, liquéfiable sous pression) ;
\item Pentane \ce{C5H12} : $T_{eb} = \SI{36}{\celsius}$ (liquide très volatil, proche de la température corporelle !).
\end{itemize}
Pour l'influence de la ramification, comparons les trois isomères de \ce{C5H12} : le pentane linéaire bout à \SI{36}{\celsius}, le 2-méthylbutane à \SI{28}{\celsius}, et le 2,2-diméthylpropane (quasiment sphérique) à seulement \SI{9,5}{\celsius}.
\end{exemplebox}

\begin{popfigure}
\begin{center}
\begin{tikzpicture}
\begin{axis}[
width=0.85\linewidth, height=6.5cm,
xlabel={Nombre $n$ d'atomes de carbone},
ylabel={$T_{eb}$ (\si{\celsius})},
xmin=0.5, xmax=10.5, ymin=-200, ymax=200,
xtick={1,2,3,4,5,6,7,8,9,10},
grid=major, grid style={popline!40},
legend style={at={(0.03,0.97)}, anchor=north west, font=\small}
]
\addplot[popcurve, domain=1:10, samples=10, mark=*, mark options={fill=poporange}]
coordinates {(1,-162) (2,-89) (3,-42) (4,-0.5) (5,36) (6,69) (7,98) (8,126) (9,151) (10,174)};
\addlegendentry{Alcanes linéaires}
\addplot[popblue, dashed, domain=0.5:10.5, samples=2] {0};
\end{axis}
\end{tikzpicture}
\legende{Température d'ébullition des alcanes linéaires en fonction du nombre $n$ d'atomes de carbone. La croissance est régulière : chaque groupe \ce{-CH2-} supplémentaire renforce les interactions de Van der Waals.}
\end{center}
\end{popfigure}

\begin{attention}[Ne confonds pas « volatil » et « léger »]
Dire qu'un alcane est volatil signifie qu'il s'évapore facilement (température d'ébullition basse). Cela ne préjuge en rien de sa densité. De plus, ne généralise pas abusivement : la régularité de la courbe concerne les alcanes \textbf{linéaires} ; les isomères ramifiés s'écartent toujours vers le bas.
\end{attention}

\courssub{Solubilité : des composés hydrophobes}

\textbf{Intuition d'abord.} Tu as observé mille fois que l'huile de palme ne se mélange pas à l'eau de cuisson des feuilles : elle forme des gouttelettes qui remontent. L'huile, comme les alcanes, est constituée de longues chaînes carbonées. Pourquoi ce refus de se mélanger ?

\begin{definition}[Composé apolaire]
Une molécule est dite \cle{apolaire} lorsqu'elle ne possède pas de pôle électrique permanent : les charges électriques y sont réparties de façon symétrique. Les alcanes sont des molécules apolaires car les liaisons \ce{C-C} sont parfaitement covalentes (atomes identiques) et les liaisons \ce{C-H} sont quasiment non polarisées (le carbone et l'hydrogène ont des électronégativités très proches).
\end{definition}

L'eau, au contraire, est un solvant \cle{polaire} : ses molécules portent des charges partielles et s'associent fortement entre elles par liaisons hydrogène. Une molécule d'alcane, apolaire, ne peut pas s'insérer dans ce réseau : l'eau la « rejette ». D'où la règle d'expérience bien connue :

\begin{aretenir}[Règle de solubilité]
\textbf{« Qui se ressemble s'assemble » : le semblable dissout le semblable.}
\begin{itemize}
\item Les alcanes sont \cle{insolubles dans l'eau} (solvant polaire) : on dit qu'ils sont \cle{hydrophobes} ;
\item les alcanes sont \cle{solubles dans les solvants organiques apolaires} : cyclohexane, benzène, tétrachlorométhane \ce{CCl4}, éther, essence\ldots
\end{itemize}
\end{aretenir}

\begin{attention}[Erreur fréquente au Probatoire]
Ne dis \textbf{jamais} que les alcanes sont solubles dans l'eau, même « un peu ». Leur solubilité dans l'eau est totalement négligeable. Conséquence pratique : on ne peut pas éteindre un feu d'hydrocarbures (essence, gazole) avec de l'eau — le carburant flotte et continue de brûler en surface, et l'eau répand les flammes ! On utilise un extincteur à poudre ou à mousse.
\end{attention}

\courssub{Densité : les alcanes surnagent sur l'eau}

\begin{propriete}[Densité des alcanes]
Tous les alcanes liquides et solides ont une densité par rapport à l'eau \cle{inférieure à 1} (comprise entre environ 0,6 et 0,8 pour les liquides courants). Moins denses que l'eau et insolubles dans l'eau, ils \cle{surnagent} : ils forment une phase supérieure bien visible.
\end{propriete}

\begin{experience}[Mélange eau + pentane]
\textbf{Protocole.} Dans un tube à essais, verser \SI{5}{mL} d'eau colorée (encre ou jus de bissap) puis \SI{5}{mL} de pentane. Boucher, agiter vigoureusement, puis laisser reposer.

\textbf{Observations.} Juste après agitation, on observe une émulsion trouble. Après repos, deux phases nettes se reforment : la phase aqueuse colorée en \textbf{bas}, la phase organique incolore (le pentane) en \textbf{haut}.

\textbf{Interprétation.} Le pentane, apolaire, est insoluble dans l'eau ; sa densité $d \approx 0,63 < 1$ le fait surnager. C'est le principe de l'extraction liquide-liquide à l'ampoule à décanter.
\end{experience}

\begin{popfigure}
\begin{center}
\begin{tikzpicture}[scale=0.95]
% Tube 1 : eau + pentane
\draw[thick, popdark] (0,0) -- (0,3.2);
\draw[thick, popdark] (1.2,0) -- (1.2,3.2);
\draw[thick, popdark] (0,0) arc[start angle=180, end angle=360, radius=0.6];
\fill[popblue!35] (0.02,0.12) -- (0.02,1.3) -- (1.18,1.3) -- (1.18,0.12) arc[start angle=360, end angle=180, radius=0.58];
\fill[poporange!30] (0.02,1.3) rectangle (1.18,2.9);
\draw[popdark, dashed] (0.02,1.3) -- (1.18,1.3);
\node[font=\small, popdark] at (0.6,2.1) {pentane};
\node[font=\scriptsize, popdark] at (0.6,1.75) {$d \approx 0{,}63$};
\node[font=\small, popdark] at (0.6,0.7) {eau colorée};
\node[font=\scriptsize, popdark] at (0.6,0.35) {$d = 1$};
\node[font=\small\bfseries, popdark] at (0.6,-0.6) {Deux phases};
% Tube 2 : nappe de pétrole (bécher)
\begin{scope}[xshift=4.5cm]
\draw[thick, popdark] (0,0) -- (0,3.2) (1.8,0) -- (1.8,3.2) (0,0) -- (1.8,0);
\fill[popblue!35] (0.03,0.03) rectangle (1.77,1.6);
\fill[popdark!70] (0.03,1.6) rectangle (1.77,2.1);
\node[font=\scriptsize, white] at (0.9,1.85) {nappe de pétrole};
\node[font=\small, popdark] at (0.9,0.8) {eau de mer};
\node[font=\small\bfseries, popdark] at (0.9,-0.6) {Marée noire};
\end{scope}
\end{tikzpicture}
\legende{À gauche : décantation eau/pentane dans un tube à essais. À droite : le pétrole brut (mélange d'alcanes), insoluble et moins dense que l'eau, forme une nappe en surface lors d'une marée noire.}
\end{center}
\end{popfigure}

\courssub{Tableau récapitulatif et conséquences pratiques}

\begin{popfigure}
\begin{center}
\begin{tabularx}{\linewidth}{|l|c|c|X|c|}
\hline
\rowcolor{popblueL} \textbf{Alcane} & \textbf{Formule} & \textbf{État à \SI{25}{\celsius}} & \textbf{$T_{eb}$ (\si{\celsius})} & \textbf{Densité} \\
\hline
Méthane & \ce{CH4} & gaz & $-162$ & --- \\
\hline
Éthane & \ce{C2H6} & gaz & $-89$ & --- \\
\hline
Propane & \ce{C3H8} & gaz & $-42$ & --- \\
\hline
Butane & \ce{C4H10} & gaz & $-0{,}5$ & --- \\
\hline
Pentane & \ce{C5H12} & liquide & $36$ & $0{,}63$ \\
\hline
Hexane & \ce{C6H14} & liquide & $69$ & $0{,}66$ \\
\hline
Octane & \ce{C8H18} & liquide & $126$ & $0{,}70$ \\
\hline
Octadécane & \ce{C18H38} & solide & $317$ & $0{,}78$ (solide) \\
\hline
\end{tabularx}
\legende{État physique, température d'ébullition et densité de quelques alcanes linéaires. Tous sont insolubles dans l'eau et solubles dans les solvants organiques apolaires.}
\end{center}
\end{popfigure}

Ces propriétés physiques ont des conséquences directes dans la vie quotidienne et l'industrie camerounaise :

\begin{itemize}
\item \textbf{Stockage des gaz domestiques.} Le butane et le propane, gazeux à pression atmosphérique mais liquéfiables sous quelques bars, sont stockés à l'état liquide dans les bouteilles : un gain de volume considérable. C'est pourquoi il ne faut jamais exposer une bouteille de gaz à une forte chaleur : la pression du gaz augmente dangereusement.
\item \textbf{Marées noires et pollution.} Le pétrole brut, mélange d'alcanes, insoluble et moins dense que l'eau, s'étale en nappes à la surface des océans, étouffant la faune marine. Le long des côtes (Kribi, Limbé, zone de la SONARA), la prévention des dégazages sauvages est un enjeu écologique majeur.
\item \textbf{Raffinage.} À la SONARA comme dans toutes les raffineries, on sépare les alcanes du pétrole brut par \cle{distillation fractionnée}, en exploitant précisément la régularité des températures d'ébullition : gaz en tête de colonne, puis essence, kérosène, gazole, et résidus lourds (bitume, paraffines) en bas.
\end{itemize}

\begin{exoresolu}[Exploiter les propriétés physiques]
\textbf{Énoncé.} On introduit dans une ampoule à décanter \SI{20}{mL} d'eau et \SI{10}{mL} d'hexane ($d = 0{,}66$). On agite puis on laisse reposer.
\begin{enumerate}
\item Combien de phases observe-t-on ? Justifier.
\item Quelle phase occupe la partie supérieure ? Justifier.
\item Citer un solvant dans lequel l'hexane serait miscible.
\end{enumerate}
\tcblower
\textbf{Solution détaillée.}
\begin{enumerate}
\item On observe \textbf{deux phases}. L'hexane est un alcane, donc une molécule apolaire : il est insoluble dans l'eau, solvant polaire. L'agitation ne fait que disperser temporairement les gouttelettes ; la décantation restaure les deux couches.
\item La phase supérieure est l'\textbf{hexane}, car sa densité $d = 0{,}66$ est inférieure à celle de l'eau ($d = 1$) : le liquide le moins dense surnage. La phase aqueuse (\SI{20}{mL}) est en bas, la phase organique (\SI{10}{mL}) en haut.
\item L'hexane est miscible à tout solvant organique apolaire : cyclohexane, tétrachlorométhane \ce{CCl4}, éther éthylique, essence\ldots{} (règle : le semblable dissout le semblable).
\end{enumerate}
\end{exoresolu}

\begin{aretenir}[L'essentiel de la section]
\begin{itemize}
\item État à \SI{25}{\celsius} : gaz pour $n \leq 4$, liquides pour $5 \leq n \leq 17$, solides pour $n \geq 18$.
\item $T_{eb}$ et $T_{fus}$ augmentent régulièrement avec $n$ ; la ramification abaisse la température d'ébullition.
\item Les alcanes sont apolaires : insolubles dans l'eau (hydrophobes), solubles dans les solvants organiques.
\item Densité $d < 1$ : les alcanes surnagent sur l'eau (nappes de pétrole).
\item Applications : bouteilles de butane/propane liquéfiés, distillation fractionnée du pétrole, extraction liquide-liquide.
\end{itemize}
\end{aretenir}

Maintenant que tu sais comment les alcanes se comportent physiquement, il reste une question essentielle : comment réagissent-ils chimiquement ? C'est l'objet de la section suivante, où tu découvriras pourquoi la combustion du butane dégage tant d'énergie et comment la lumière permet de substituer des atomes d'halogène aux atomes d'hydrogène.

\coursec{Propriétés chimiques : combustion et substitution}

Dans la section précédente, nous avons vu que les alcanes sont des composés peu solubles dans l'eau, de densité inférieure à celle de l'eau, et que leur état physique dépend de la longueur de la chaîne carbonée. Mais ces propriétés physiques ne disent rien de leur \emph{comportement chimique}. Or c'est bien ce comportement qui explique pourquoi la bouteille de gaz butane de nos cuisines, l'essence de nos motos-taxis ou le gaz naturel sont si précieux : les alcanes \emph{brûlent} en dégageant beaucoup d'énergie, et ils peuvent aussi \emph{réagir} avec les halogènes. Deux grandes familles de réactions vont donc retenir notre attention : les \cle{réactions de destruction} (la molécule est détruite, le squelette carboné est brisé) et les \cle{réactions de substitution} (le squelette carboné reste intact, mais des atomes d'hydrogène sont remplacés).

\courssub{Les réactions de destruction}

\textbf{Intuition.} Quand tu allumes le gaz butane de la cuisinière, la flamme bleue qui apparaît est le signe visible d'une réaction chimique : les molécules de butane sont littéralement \emph{détruites} par le dioxygène de l'air. Leurs atomes de carbone et d'hydrogène se recombinent avec l'oxygène pour former de nouvelles molécules. C'est une \cle{combustion}.

\courssub{La combustion complète dans le dioxygène}

\begin{definition}[Combustion complète d'un alcane]
La \cle{combustion complète} d'un alcane est sa réaction avec le dioxygène en \emph{excès}. Elle détruit entièrement la molécule et produit exclusivement du \cle{dioxyde de carbone} \ce{CO2} et de l'\cle{eau} \ce{H2O}, avec un important dégagement de chaleur (réaction \emph{exothermique}).
\end{definition}

Pour un alcane de formule générale $C_nH_{2n+2}$, l'équation générale de la combustion complète s'écrit en mode mathématique standard (indices littéraux) :

\[ C_nH_{2n+2} + \frac{3n+1}{2} O_2 \longrightarrow n CO_2 + (n+1) H_2O \]

Vérifions sur le méthane ($n = 1$) : il faut $\frac{3\times 1+1}{2} = 2$ moles de \ce{O2}, et l'on obtient $1$ mole de \ce{CO2} et $2$ moles d'eau :

\[ \ce{CH4 + 2 O2 -> CO2 + 2 H2O} \]

Cette réaction nécessite une \emph{amorce} (une flamme ou une étincelle) ; ensuite elle s'auto-entretient car la chaleur dégagée entretient la combustion.

\begin{methode}[Équilibrer l'équation de combustion d'un alcane]
Pour équilibrer la combustion complète de $C_nH_{2n+2}$, procède toujours dans l'ordre \textbf{C -- H -- O} :
\begin{enumerate}
\item \textbf{Équilibre le carbone} : place le coefficient $n$ devant \ce{CO2}.
\item \textbf{Équilibre l'hydrogène} : place le coefficient $(n+1)$ devant \ce{H2O} (car l'alcane contient $2n+2$ atomes d'hydrogène).
\item \textbf{Équilibre l'oxygène en dernier} : compte les atomes d'oxygène à droite, soit $2n + (n+1) = 3n+1$, puis place $\frac{3n+1}{2}$ devant \ce{O2}.
\item Si un coefficient fractionnaire te gêne, multiplie \emph{toute} l'équation par 2.
\end{enumerate}
\end{methode}

\begin{exemplebox}[Combustion complète du butane, le gaz de nos cuisines]
Le butane a pour formule brute \ce{C4H10} (ici $n = 4$). Appliquons la méthode :
\begin{enumerate}
\item Carbone : $4$ atomes à gauche $\Rightarrow$ coefficient $4$ devant \ce{CO2}.
\item Hydrogène : $10$ atomes à gauche $\Rightarrow$ coefficient $5$ devant \ce{H2O}.
\item Oxygène : à droite, $4\times 2 + 5 = 13$ atomes $\Rightarrow$ coefficient $\frac{13}{2}$ devant \ce{O2}.
\end{enumerate}
D'où : $\ce{C4H10 + 13/2 O2 -> 4 CO2 + 5 H2O}$. En multipliant tout par 2 pour éliminer la fraction :
\[ \ce{2 C4H10 + 13 O2 -> 8 CO2 + 10 H2O} \]
\end{exemplebox}

\begin{exoresolu}[Combustion du propane et volume de dioxygène]
On brûle complètement \SI{2,0}{mol} de propane \ce{C3H8} dans le dioxygène. Écris l'équation de la réaction, puis détermine la quantité de dioxygène nécessaire et le volume de \ce{CO2} dégagé, mesuré dans les conditions où $V_m = \SI{24,0}{L.mol^{-1}}$.
\tcblower
\textbf{Équation.} Avec $n = 3$ : coefficient de \ce{O2} : $\frac{3\times 3+1}{2} = 5$ ; coefficients $3$ devant \ce{CO2} et $4$ devant \ce{H2O} :
\[ \ce{C3H8 + 5 O2 -> 3 CO2 + 4 H2O} \]
\textbf{Quantité de dioxygène.} D'après les coefficients stœchiométriques, $1$ mole de propane consomme $5$ moles de \ce{O2}. Donc :
\[ n(\ce{O2}) = 5 \times n(\ce{C3H8}) = 5 \times 2,0 = \SI{10}{mol} \]
\textbf{Volume de \ce{CO2}.} Il se forme $3 \times 2,0 = \SI{6,0}{mol}$ de \ce{CO2}, soit :
\[ V(\ce{CO2}) = n \times V_m = 6,0 \times 24,0 = \SI{144}{L} \]
\end{exoresolu}

\textbf{Pourquoi les alcanes sont-ils d'excellents combustibles ?} Parce que leur combustion dégage énormément de chaleur : on dit qu'ils possèdent un \cle{pouvoir calorifique} élevé. Le pouvoir calorifique est l'énergie libérée par la combustion complète d'un kilogramme (ou d'un mètre cube) de combustible. Il vaut environ \SI{50}{MJ.kg^{-1}} pour le méthane (gaz naturel) et \SI{46}{MJ.kg^{-1}} pour le butane. C'est ce qui explique l'usage massif des alcanes comme \emph{combustibles} : gaz butane et propane en bouteille pour la cuisine, essence et gasoil (mélanges d'alcanes) pour les moteurs, gaz naturel pour le chauffage et la production d'électricité (centrale de Kribi au Cameroun).

\begin{popfigure}
\begin{center}
\begin{tikzpicture}[font=\small]
% Axe énergie
\draw[-{Stealth},thick,popdark] (0,0) -- (0,5.2) node[above left=2pt and -18pt] {\textbf{Énergie}};
\draw[-{Stealth},thick,popdark] (0,0) -- (9.6,0) node[below right] {Avancement de la réaction};
% Niveaux
\draw[thick,popblue] (0.6,3.8) -- (3.4,3.8);
\draw[thick,poporange] (6.0,0.9) -- (9.0,0.9);
\node[popblue,anchor=south] at (2.0,3.8) {\ce{CH4} + \ce{2 O2} (réactifs)};
\node[poporange,anchor=north] at (7.5,0.85) {\ce{CO2} + \ce{2 H2O} (produits)};
% Flèche de dégagement
\draw[-{Stealth},very thick,popink] (4.6,3.8) -- (4.6,0.9) node[midway,right,align=left] {énergie dégagée\\ $\approx \SI{890}{kJ.mol^{-1}}$};
% Courbe de réaction
\draw[very thick,popgreen] (3.4,3.8) .. controls (4.4,3.8) and (4.6,4.9) .. (5.2,4.6) .. controls (5.8,4.3) and (5.6,0.9) .. (6.0,0.9);
\node[popgreen,anchor=south west] at (4.7,4.55) {\footnotesize énergie d'activation};
\end{tikzpicture}
\end{center}
\legende{Bilan énergétique de la combustion complète du méthane : les produits sont plus stables que les réactifs ; la différence d'énergie est libérée sous forme de chaleur.}
\end{popfigure}

\courssub{La combustion incomplète : un danger mortel}

Que se passe-t-il si le dioxygène \emph{manque} ? La combustion n'est plus complète : au lieu de \ce{CO2}, on obtient du \cle{monoxyde de carbone} \ce{CO}, ou même du \cle{carbone} pur (la suie qui noircit les casseroles et les marmites au-dessus des feux de bois ou des lampes à pétrole).

\[ \ce{2 CH4 + 3 O2 -> 2 CO + 4 H2O} \qquad \text{(dioxygène insuffisant)} \]
\[ \ce{CH4 + O2 -> C + 2 H2O} \qquad \text{(dioxygène très insuffisant : suie)} \]

\begin{attention}[Le monoxyde de carbone, tueur invisible]
Le monoxyde de carbone \ce{CO} est un gaz \textbf{incolore, inodore et mortel}. Inspiré, il se fixe sur l'hémoglobine du sang à la place du dioxygène et provoque l'asphyxie des organes, sans aucun signe avant-coureur. C'est pourquoi il est \textbf{formellement interdit} d'utiliser un brasero, un réchaud à charbon ou une cuisinière à gaz dans une pièce fermée et mal aérée : la combustion devient incomplète et le \ce{CO} s'accumule. Chaque année, des intoxications mortelles surviennent ainsi pendant la saison fraîche. Aère toujours la pièce où brûle un combustible !
\end{attention}

\courssub{La destruction par le dichlore (combustion dans le dichlore)}

Le dioxygène n'est pas le seul capable de détruire un alcane. Mélangé au \cle{dichlore} et exposé brutalement à une lumière intense (ou à une flamme), le méthane subit une réaction \emph{explosive} qui détruit complètement la molécule : il se forme du \textbf{carbone} (un dépôt noir de suie) et du \textbf{chlorure d'hydrogène} \ce{HCl} :

\[ \ce{CH4 + 2 Cl2 -> C + 4 HCl} \qquad \text{(lumière vive, réaction explosive)} \]

On reconnaît le carbone au dépôt noir qui tapisse le récipient, et le chlorure d'hydrogène à son acidité (il bleuit le papier pH et donne un précipité blanc avec une solution de nitrate d'argent). Cette réaction est parfois appelée \emph{combustion dans le dichlore} dans les programmes.

\courssub{Les réactions de substitution : l'halogénation ménagée}

\textbf{Intuition.} Reprenons le mélange méthane + dichlore, mais cette fois \emph{à la lumière douce} (lumière ultraviolette modérée) et en contrôlant les proportions. Au lieu d'exploser, la réaction devient progressive : les atomes de chlore remplacent \emph{un par un} les atomes d'hydrogène du méthane, \textbf{sans briser le squelette carboné}. C'est une toute autre réaction.

\begin{definition}[Réaction de substitution]
Une \cle{réaction de substitution} est une réaction au cours de laquelle un atome (ou un groupe d'atomes) d'une molécule est \emph{remplacé} par un autre atome (ou groupe d'atomes), le squelette carboné de la molécule restant intact.
\end{definition}

Dans le cas des alcanes, l'halogénation (chloration ou bromation) se déroule \textbf{en présence de lumière ultraviolette} (notée $h\nu$ au-dessus de la flèche) : un atome d'hydrogène est remplacé par un atome d'halogène, et il se forme en même temps une molécule d'halogénure d'hydrogène (\ce{HCl} ou \ce{HBr}).

\textbf{Les quatre substitutions successives du méthane (chloration).} Le méthane possède quatre atomes d'hydrogène : la chloration peut donc se poursuivre en quatre étapes, chaque étape produisant un dérivé chloré différent et du \ce{HCl} :

\begin{popfigure}
\begin{center}
\begin{tikzpicture}[node distance=0.35cm and 0.9cm, font=\small]
\node (a) {\chemfig{CH_4}};
\node[right=of a] (b) {\chemfig{CH_3Cl}};
\node[right=of b] (c) {\chemfig{CH_2Cl_2}};
\node[right=of c] (d) {\chemfig{CHCl_3}};
\node[right=of d] (e) {\chemfig{CCl_4}};
\draw[-{Stealth},thick,popblue] (a) -- (b) node[midway,above,popdark]{\footnotesize \ce{Cl2}} node[midway,below,popdark]{\footnotesize \ce{-HCl}};
\draw[-{Stealth},thick,popblue] (b) -- (c) node[midway,above,popdark]{\footnotesize \ce{Cl2}} node[midway,below,popdark]{\footnotesize \ce{-HCl}};
\draw[-{Stealth},thick,popblue] (c) -- (d) node[midway,above,popdark]{\footnotesize \ce{Cl2}} node[midway,below,popdark]{\footnotesize \ce{-HCl}};
\draw[-{Stealth},thick,popblue] (d) -- (e) node[midway,above,popdark]{\footnotesize \ce{Cl2}} node[midway,below,popdark]{\footnotesize \ce{-HCl}};
\node[below=0.25cm of a, popmuted] {\footnotesize méthane};
\node[below=0.25cm of b, popmuted, align=center] {\footnotesize chloro-\\ \footnotesize méthane};
\node[below=0.25cm of c, popmuted, align=center] {\footnotesize dichloro-\\ \footnotesize méthane};
\node[below=0.25cm of d, popmuted, align=center] {\footnotesize trichloro-\\ \footnotesize méthane};
\node[below=0.25cm of e, popmuted, align=center] {\footnotesize tétrachloro-\\ \footnotesize méthane};
\end{tikzpicture}
\end{center}
\legende{Les quatre substitutions successives du méthane par le dichlore, en présence de lumière UV. À chaque étape, un atome d'hydrogène est remplacé par un atome de chlore et une molécule de \ce{HCl} se forme.}
\end{popfigure}

Les équations correspondantes sont :

\begin{align*}
\ce{CH4 + Cl2 &->[h\nu] CH3Cl + HCl} && \text{chlorométhane}\\
\ce{CH3Cl + Cl2 &->[h\nu] CH2Cl2 + HCl} && \text{dichlorométhane}\\
\ce{CH2Cl2 + Cl2 &->[h\nu] CHCl3 + HCl} && \text{trichlorométhane (chloroforme)}\\
\ce{CHCl3 + Cl2 &->[h\nu] CCl4 + HCl} && \text{tétrachlorométhane}
\end{align*}

La bromation suit le même schéma, par exemple pour la première étape :
\[ \ce{CH4 + Br2 ->[h\nu] CH3Br + HBr} \]

\begin{aretenir}[L'essentiel sur la substitution]
\begin{itemize}
\item Conditions : \textbf{lumière ultraviolette} (ou chaleur modérée) ; le dibrome \ce{Br2} réagit de façon analogue.
\item À chaque étape : \textbf{un H remplacé par un Cl (ou Br)} + formation d'\textbf{une molécule de \ce{HCl} (ou \ce{HBr})}.
\item On obtient en général un \emph{mélange} des quatre produits ; en jouant sur les proportions de \ce{Cl2}, on favorise l'un d'eux.
\item Le \textbf{squelette carboné est conservé} : c'est ce qui distingue la substitution de la destruction.
\end{itemize}
\end{aretenir}

\begin{attention}[Ne confonds jamais substitution et destruction !]
C'est l'erreur classique au Probatoire. Avec le \emph{même} couple de réactifs (méthane + dichlore), deux réactions différentes sont possibles :
\begin{itemize}
\item \textbf{Substitution} (lumière UV douce, progressive) : produits \emph{organiques chlorés} \ce{CH3Cl}, \ce{CH2Cl2}, \ce{CHCl3}, \ce{CCl4} + \ce{HCl}. Le squelette carboné survit.
\item \textbf{Destruction} (lumière vive brutale, explosive) : \emph{carbone} \ce{C} + \ce{HCl}. La molécule est pulvérisée.
\end{itemize}
Retiens le réflexe : \emph{substitution = dérivés chlorés ; destruction = dépôt noir de carbone}.
\end{attention}

\textbf{Comment la lumière déclenche-t-elle la substitution ? (mécanisme simplifié).} La chloration du méthane est une \cle{réaction en chaîne radicalaire} : elle fait intervenir des espèces très réactives appelées \emph{radicaux libres}, atomes ou groupes d'atomes possédant un électron célibataire (noté par un point, ex. \ce{Cl^.}). On distingue deux phases principales :

\begin{enumerate}
\item \textbf{Initiation} : la lumière UV fournit l'énergie nécessaire pour casser la liaison Cl–Cl de la molécule de dichlore en deux atomes de chlore radicaux :
\[ \ce{Cl2 ->[h\nu] 2 Cl^.} \]
\item \textbf{Propagation} : le radical \ce{Cl^.} arrache un atome d'hydrogène au méthane, formant \ce{HCl} et un radical méthyle \ce{CH3^.} ; celui-ci attaque à son tour une molécule de \ce{Cl2} pour donner \ce{CH3Cl} et régénérer un radical \ce{Cl^.}, qui peut recommencer le cycle :
\[ \ce{Cl^. + CH4 -> HCl + CH3^.} \qquad \ce{CH3^. + Cl2 -> CH3Cl + Cl^.} \]
\end{enumerate}

Comme le radical \ce{Cl^.} est régénéré à chaque cycle, un seul photon absorbé peut déclencher des milliers de substitutions : c'est une \emph{réaction en chaîne}. La chaîne s'arrête (phase de \emph{terminaison}) lorsque deux radicaux se rencontrent et se recombinent. Cette présentation est volontairement qualitative : retiens simplement que la lumière \emph{initie} la réaction en créant des radicaux, et que ceux-ci \emph{propagent} la réaction en chaîne.

\courssub{Importance des dérivés halogénés des alcanes}

Les dérivés halogénés des alcanes, obtenus par substitution, ont des applications industrielles considérables.

\begin{popfigure}
\begin{center}
\begin{tikzpicture}[font=\small]
\node[draw=popblue, fill=popblueL, rounded corners=3pt, text width=3.1cm, align=center, minimum height=1.5cm] (a) at (0,0) {\textbf{\ce{CH3Cl}}\\ chlorométhane\\ \footnotesize synthèse des silicones};
\node[draw=popgreen, fill=popgreenL, rounded corners=3pt, text width=3.1cm, align=center, minimum height=1.5cm] (b) at (4.6,0) {\textbf{\ce{CH2Cl2}}\\ dichlorométhane\\ \footnotesize solvant, décapant};
\node[draw=poporange, fill=poporangeL, rounded corners=3pt, text width=3.1cm, align=center, minimum height=1.5cm] (c) at (9.2,0) {\textbf{\ce{CHCl3}}\\ trichlorométhane\\ \footnotesize solvant, ancien anesthésique};
\node[draw=poppurple, fill=poppurpleL, rounded corners=3pt, text width=3.1cm, align=center, minimum height=1.5cm] (d) at (13.8,0) {\textbf{\ce{CCl4}}\\ tétrachlorométhane\\ \footnotesize solvant, extincteurs (jadis)};
\end{tikzpicture}
\end{center}
\legende{Noms et usages des quatre dérivés chlorés du méthane.}
\end{popfigure}

Au-delà des dérivés du méthane, citons :

\begin{itemize}
\item \textbf{Solvants} : le trichlorométhane (chloroforme) et le dichlorométhane dissolvent de nombreuses substances organiques ; ils servent au laboratoire et dans l'industrie.
\item \textbf{Anesthésiques} : le chloroforme fut l'un des premiers anesthésiques chirurgicaux (XIX\textsuperscript{e} siècle), avant d'être abandonné à cause de sa toxicité ; l'halothane \ce{CF3-CHBrCl}, dérivé halogéné de l'éthane, l'a remplacé.
\item \textbf{Réfrigérants} : les fréons (CFC, chlorofluorocarbures) comme \ce{CF2Cl2} ont longtemps équipé réfrigérateurs et climatiseurs.
\item \textbf{Insecticides} : le DDT, dérivé chloré, a été massivement utilisé contre les moustiques vecteurs du paludisme, avant son interdiction pour sa persistance dans l'environnement.
\end{itemize}

\begin{savaistu}[Les fréons, la couche d'ozone et le protocole de Montréal]
Les CFC (chlorofluorocarbures), dits « fréons », semblaient parfaits : ininflammables, non toxiques, excellents fluides réfrigérants. Mais, très stables, ils montent jusqu'à la stratosphère où les UV les décomposent en libérant des atomes de chlore... qui détruisent l'\cle{ozone} \ce{O3}, ce bouclier naturel qui nous protège des rayonnements ultraviolets dangereux. C'est ainsi que s'est formé le « trou » dans la couche d'ozone. En 1987, la communauté internationale a signé le \textbf{protocole de Montréal}, qui interdit progressivement les CFC : c'est l'un des plus grands succès de la diplomatie environnementale, et la couche d'ozone se reconstitue lentement. Les réfrigérateurs modernes utilisent des substituts sans chlore.
\end{savaistu}

\begin{exercice}[Pour t'entraîner]
\begin{enumerate}
\item Écris et équilibre l'équation de la combustion complète du pentane \ce{C5H12}.
\item On expose un mélange d'éthane \ce{C2H6} et de dichlore à la lumière UV. Écris l'équation de la première substitution et nomme le produit organique obtenu.
\item Explique pourquoi il est dangereux de faire fonctionner un groupe électrogène à essence dans une pièce fermée.
\end{enumerate}
\end{exercice}

\begin{corrige}[Exercice]
\begin{enumerate}
\item Avec $n = 5$ : coefficient de \ce{O2} : $\frac{3\times 5+1}{2} = 8$ ; coefficients $5$ devant \ce{CO2} et $6$ devant \ce{H2O} :
\[ \ce{C5H12 + 8 O2 -> 5 CO2 + 6 H2O} \]
\item Première substitution : un hydrogène de l'éthane est remplacé par un chlore :
\[ \ce{C2H6 + Cl2 ->[h\nu] C2H5Cl + HCl} \]
Le produit organique est le \textbf{chloroéthane}.
\item Le moteur à essence réalise la combustion d'alcanes. Dans une pièce fermée, le dioxygène s'épuise et la combustion devient incomplète : elle produit du monoxyde de carbone \ce{CO}, gaz inodore et mortel qui s'accumule et intoxique les occupants.
\end{enumerate}
\end{corrige}

\begin{aretenir}[Bilan de la section]
\begin{itemize}
\item \textbf{Combustion complète} (excès de \ce{O2}) : $\ce{$C_nH_{2n+2}$ + $\frac{3n+1}{2}$ O2 -> $n$ CO2 + $(n+1)$ H2O}$, très exothermique : c'est la base de l'usage des alcanes comme combustibles.
\item \textbf{Combustion incomplète} (manque de \ce{O2}) : formation de \ce{CO} (gaz mortel) ou de carbone (suie).
\item \textbf{Destruction par \ce{Cl2}} (lumière vive) : réaction explosive donnant \ce{C} + \ce{HCl} (combustion dans le dichlore).
\item \textbf{Substitution} (lumière UV) : remplacement successif des H par des Cl (ou Br), squelette carboné conservé, formation de \ce{CH3Cl}, \ce{CH2Cl2}, \ce{CHCl3}, \ce{CCl4} et de \ce{HCl} à chaque étape ; mécanisme radicalaire en chaîne.
\item Les dérivés halogénés sont utiles (solvants, anesthésiques, réfrigérants) mais certains, comme les CFC, sont dangereux pour l'environnement.
\end{itemize}
\end{aretenir}

\coursec{Préparation du méthane au laboratoire}

Dans la section précédente, tu as découvert que les alcanes réagissent avec le dioxygène (combustion) et avec les halogènes (substitution). Mais une question naturelle se pose : d'où viennent ces alcanes ? Comment le chimiste peut-il en obtenir au laboratoire pour les étudier ? Dans cette dernière section, nous allons apprendre à \cle{préparer le méthane}, le plus simple des alcanes, à partir d'un solide grisâtre appelé \cle{carbure d'aluminium}. Cette préparation, classique au programme de Première, illustre parfaitement la démarche expérimentale : choisir les réactifs, monter un dispositif adapté, recueillir le gaz, puis l'identifier.

\courssub{Principe de la préparation : action de l'eau sur le carbure d'aluminium}

\textbf{L'idée intuitive.}\par
Le méthane \ce{CH4} contient du carbone et de l'hydrogène. Pour le fabriquer, il faut donc un réactif qui \emph{apporte le carbone} et un réactif qui \emph{apporte l'hydrogène}. Le candidat idéal pour apporter le carbone est le \cle{carbure d'aluminium}, solide de formule \ce{Al4C3}, dans lequel le carbone est lié à l'aluminium. L'eau \ce{H2O}, elle, apporte l'hydrogène. Lorsqu'on verse de l'eau sur le carbure d'aluminium, les liaisons aluminium--carbone sont rompues : le carbone s'associe à l'hydrogène de l'eau pour former le méthane, tandis que l'aluminium s'associe aux groupes hydroxyde pour former l'\cle{hydroxyde d'aluminium} \ce{Al(OH)3}, un précipité blanc gélatineux.

\begin{definition}[Préparation du méthane]
Au laboratoire, on prépare le méthane par action de l'eau sur le carbure d'aluminium \ce{Al4C3}. C'est une réaction d'\cle{hydrolyse} : l'eau décompose le carbure en dégageant du méthane gazeux et en formant de l'hydroxyde d'aluminium.
\end{definition}

\textbf{Équation-bilan de la réaction.}\par
Écrivons et équilibrons l'équation étape par étape. Les réactifs sont \ce{Al4C3} et \ce{H2O} ; les produits sont \ce{Al(OH)3} et \ce{CH4}.

\begin{itemize}
\item \textbf{Aluminium} : 4 atomes à gauche, il faut donc 4 \ce{Al(OH)3} à droite ;
\item \textbf{Carbone} : 3 atomes à gauche, il faut donc 3 \ce{CH4} à droite ;
\item \textbf{Hydrogène} : à droite, $4 \times 3 + 3 \times 4 = 24$ atomes H, il faut donc 12 \ce{H2O} à gauche ;
\item \textbf{Oxygène} : $12$ atomes O à gauche, et $4 \times 3 = 12$ à droite. L'équation est équilibrée.
\end{itemize}

\[ \ce{Al4C3 + 12 H2O -> 4 Al(OH)3 + 3 CH4 ^} \]

\begin{aretenir}[Caractéristiques de la réaction]
\begin{itemize}
\item La réaction se produit à \cle{température ordinaire}, sans chauffage ni catalyseur ;
\item elle est \cle{lente} mais régulière, ce qui permet de contrôler le débit de gaz en réglant l'arrivée d'eau ;
\item elle est \cle{totale} : tout le carbure finit par être consommé ;
\item le méthane se dégage sous forme de bulles ; l'hydroxyde d'aluminium apparaît comme un précipité blanc dans le ballon.
\end{itemize}
\end{aretenir}

\courssub{Dispositif expérimental}

\textbf{Choix du mode de recueillement.}\par
Le méthane est un gaz \cle{très peu soluble dans l'eau} (environ \SI{30}{mL} de gaz dissous par litre d'eau à température ordinaire) et \cle{moins dense que l'air} (masse molaire $M = \SI{16}{g.mol^{-1}} < \SI{29}{g.mol^{-1}}$). On peut donc le recueillir par \cle{déplacement d'eau} dans une éprouvette (ou un tube à essais) renversée sur une cuve à eau : le gaz, insoluble, chasse l'eau et s'accumule en haut du récipient.

\begin{methode}[Recueillir un gaz par déplacement d'eau]
\begin{enumerate}
\item Vérifier que le gaz à recueillir est \textbf{très peu soluble dans l'eau} (c'est le cas du méthane) ;
\item remplir complètement d'eau une éprouvette graduée et la renverser, ouverture vers le bas, dans un cristallisoir rempli d'eau, sans laisser entrer d'air ;
\item introduire l'extrémité du tube dégagement sous l'ouverture de l'éprouvette ;
\item laisser le gaz chasser l'eau : le niveau d'eau descend dans l'éprouvette ;
\item lorsque l'éprouvette est pleine de gaz, la boucher sous l'eau avant de la sortir.
\end{enumerate}
\end{methode}

\textbf{Description du montage.}\par
Le dispositif complet comporte trois parties :
\begin{itemize}
\item un \cle{ballon} (ou flacon réactionnel) contenant le carbure d'aluminium en morceaux, fermé par un bouchon à deux trous ;
\item une \cle{ampoule de coulée} contenant l'eau, dont le robinet permet de faire tomber l'eau goutte à goutte sur le carbure : c'est elle qui commande la vitesse de la réaction ;
\item un \cle{tube à dégagement} qui conduit le méthane jusqu'à une éprouvette renversée sur l'eau d'un cristallisoir.
\end{itemize}

\begin{popfigure}
\begin{center}
\begin{tikzpicture}[scale=0.9, every node/.style={font=\small}, >=Stealth]
% --- Ballon réactionnel ---
\draw[thick, popblue] (0,0) circle (1.2);
\draw[thick, popblue] (-0.3,1.15) -- (-0.3,2.0) -- (0.3,2.0) -- (0.3,1.15);
% Carbure dans le ballon
\foreach \x/\y in {-0.5/-0.7, 0.2/-0.8, 0.6/-0.5, -0.1/-0.4, 0.4/-0.2, -0.6/-0.2}
  \fill[poppurple!60] (\x,\y) circle (0.12);
\node[popdark, anchor=north] at (0,-1.6) {Carbure d'aluminium \ce{Al4C3}};
\node[popdark, anchor=south] at (0,2.3) {Ballon};
% --- Ampoule de coulée ---
\draw[thick, poporange] (-0.8,2.0) -- (-0.8,4.0) .. controls (-0.8,4.5) and (-0.1,4.5) .. (-0.1,4.0) -- (-0.1,2.0);
\draw[thick, poporange] (-0.55,2.0) -- (-0.55,1.3);
\draw[thick, poporange] (-0.35,2.0) -- (-0.35,1.3);
\draw[thick, poporange] (-0.7,1.75) -- (-0.2,1.75); % robinet
\fill[popblue!20] (-0.85,3.2) rectangle (0.05,3.9);
\node[popdark, align=center, anchor=east] at (-1.5,3.5) {Ampoule de coulée\\ (eau)};
\draw[popdark, ->] (-1.0,3.5) -- (-0.85,3.5);
% Gouttes d'eau
\fill[popblue!70] (-0.45,1.1) circle (0.04);
\fill[popblue!70] (-0.45,0.7) circle (0.04);
\fill[popblue!70] (-0.45,0.3) circle (0.04);
% --- Tube à dégagement ---
\draw[thick, popgreen] (0.3,1.8) -- (1.2,1.8) -- (1.2,0.5) -- (5.0,0.5) -- (5.0,-1.0);
\node[popdark, anchor=south] at (2.5,1.9) {Tube à dégagement};
% --- Cristallisoir ---
\draw[thick, popblue] (3.5,-2.8) -- (3.5,-1.3) -- (6.5,-1.3) -- (6.5,-2.8) -- cycle;
\fill[popblue!10] (3.55,-2.75) rectangle (6.45,-1.35);
\node[popdark, anchor=west] at (6.8,-2.0) {Cristallisoir (eau)};
% --- Éprouvette renversée ---
\draw[thick, poppink] (4.5,-1.3) -- (4.5,1.0) -- (5.5,1.0) -- (5.5,-1.3);
\fill[popgold!30] (4.55,-0.5) rectangle (5.45,0.95);
\node[popdark, align=center, anchor=west] at (5.8,0.5) {Éprouvette renversée\\ pleine de méthane};
\draw[popdark, ->] (5.55,0.5) -- (5.8,0.5);
% Bulles de méthane
\foreach \x/\y/\r in {4.8/-1.0/0.05, 4.7/-0.7/0.06, 4.9/-0.4/0.05, 5.0/-0.1/0.04}
  \fill[popgold!80] (\x,\y) circle (\r);
% Précipité Al(OH)3 dans le ballon (optionnel, visuel)
\fill[popgray!50] (-0.3,-0.9) ellipse (0.4 and 0.15);
\end{tikzpicture}
\end{center}
\legende{Dispositif de préparation du méthane au laboratoire : l'eau tombe goutte à goutte sur le carbure d'aluminium ; le méthane dégagé est recueilli par déplacement d'eau.}
\end{popfigure}

\begin{experience}[Protocole opératoire]
\begin{itemize}
\item \textbf{Dispositif} : celui de la figure ci-dessus.
\item \textbf{Protocole} : introduire quelques morceaux de carbure d'aluminium dans le ballon sec ; adapter le bouchon muni de l'ampoule de coulée et du tube à dégagement ; vérifier l'étanchéité ; ouvrir doucement le robinet de l'ampoule pour laisser couler l'eau goutte à goutte ; laisser d'abord s'échapper le gaz pendant quelques secondes (il chasse l'air du montage), puis recueillir le méthane dans l'éprouvette renversée.
\item \textbf{Observations} : dégagement régulier de bulles ; apparition d'un précipité blanc gélatineux d'hydroxyde d'aluminium ; l'eau descend progressivement dans l'éprouvette.
\item \textbf{Interprétation} : le gaz recueilli est le méthane, produit par hydrolyse du carbure selon $\ce{Al4C3 + 12 H2O -> 4 Al(OH)3 + 3 CH4}$.
\end{itemize}
\end{experience}

\courssub{Propriétés et identification du méthane recueilli}

\textbf{Propriétés physiques.}\par
L'observation directe du gaz recueilli montre que le méthane est :
\begin{itemize}
\item \cle{incolore} : l'éprouvette pleine de gaz paraît vide ;
\item \cle{inodore} : il n'a aucune odeur (l'odeur du gaz domestique vient d'un produit ajouté volontairement, le tétrahydrothiophène, pour détecter les fuites) ;
\item \cle{très peu soluble dans l'eau} : il peut être recueilli par déplacement d'eau ;
\item \cle{moins dense que l'air} : une éprouvette pleine de méthane doit être conservée ouverture vers le bas, sinon le gaz s'échappe vers le haut ;
\item \cle{inflammable} : il brûle facilement en présence de dioxygène.
\end{itemize}

\textbf{Tests d'identification.}\par
Pour prouver que le gaz recueilli est bien un alcane combustible, on réalise sa combustion, déjà étudiée dans la section précédente :

\[ \ce{CH4 + 2 O2 -> CO2 + 2 H2O} \]

On approche une flamme de l'éprouvette : le méthane brûle avec une \cle{flamme bleue} non éclairante (combustion complète). On identifie ensuite les produits :
\begin{itemize}
\item la \textbf{vapeur d'eau} formée se condense en buée sur les parois froides d'un verre de montre placé au-dessus de la flamme ;
\item le \textbf{dioxyde de carbone} formé \cle{trouble l'eau de chaux} (solution d'hydroxyde de calcium) : il se forme un précipité blanc de carbonate de calcium selon $\ce{CO2 + Ca(OH)2 -> CaCO3 v + H2O}$.
\end{itemize}

\begin{attention}[Sécurité : un gaz inflammable]
Le méthane est inflammable et forme avec l'air des mélanges \textbf{explosifs} (entre 5 \% et 15 \% en volume). Avant et pendant l'expérience :
\begin{itemize}
\item éloigner \textbf{toute flamme} du dégagement gazeux (bec Bunsen, bougie) ;
\item vérifier l'\textbf{étanchéité} du montage (bouchons, raccords) avant de commencer ;
\item laisser d'abord s'échapper le gaz quelques secondes avant le recueillement, afin de chasser l'air contenu dans le ballon : un mélange air--méthane s'enflammerait violemment ;
\item ne jamais boucher le tube à dégagement pendant que la réaction se poursuit (risque de surpression).
\end{itemize}
\end{attention}

\begin{attention}[Erreur fréquente : carbure d'aluminium ou carbure de calcium ?]
Ne confonds pas les deux carbures classiques du programme :
\begin{itemize}
\item le \textbf{carbure d'aluminium} \ce{Al4C3} + eau donne le \textbf{méthane} \ce{CH4} (un alcane) ;
\item le \textbf{carbure de calcium} \ce{CaC2} + eau donne l'\textbf{acétylène} (éthyne) \ce{C2H2} (un alcyne) : $\ce{CaC2 + 2 H2O -> Ca(OH)2 + C2H2}$.
\end{itemize}
Astuce mnémotechnique : \emph{Aluminium} commence comme \emph{alcane}... le carbure d'aluminium donne l'alcane le plus simple. Cette confusion coûte cher au Probatoire !
\end{attention}

\courssub{Exemple chiffré : calcul du volume de méthane produit}

\begin{exoresolu}[Volume de méthane dégagé]
On fait agir de l'eau en excès sur \SI{144}{g} de carbure d'aluminium pur. Calculer la quantité de matière de méthane obtenue, puis le volume de gaz recueilli dans les CNTP ($V_m = \SI{22,4}{L.mol^{-1}}$). On donne $M(\ce{Al4C3}) = 4 \times 27 + 3 \times 12 = \SI{144}{g.mol^{-1}}$.
\tcblower
\textbf{Étape 1 : quantité de matière de carbure.}
\[ n(\ce{Al4C3}) = \frac{m}{M} = \frac{144}{144} = \SI{1,00}{mol} \]
\textbf{Étape 2 : tableau d'avancement} (l'eau est en excès).
\begin{center}
\begin{tabular}{lccc}
\toprule
Équation & \ce{Al4C3} & \ce{+ 12 H2O ->} & \ce{4 Al(OH)3 + 3 CH4} \\
\midrule
État initial (mol) & $1,00$ & excès & $0$ \\
État final (mol) & $0$ & excès & $3,00$ \\
\bottomrule
\end{tabular}
\end{center}
D'après les coefficients stœchiométriques, 1 mol de carbure produit 3 mol de méthane :
\[ n(\ce{CH4}) = 3 \times n(\ce{Al4C3}) = \SI{3,00}{mol} \]
\textbf{Étape 3 : volume de gaz dans les CNTP.}
\[ V(\ce{CH4}) = n \times V_m = 3,00 \times 22,4 = \SI{67,2}{L} \]
\textbf{Conclusion} : \SI{144}{g} de carbure d'aluminium produisent \SI{3,00}{mol} de méthane, soit un volume de \SI{67,2}{L} dans les CNTP.
\end{exoresolu}

\begin{savaistu}[Le méthane, gaz des marais et gaz du futur]
Le méthane est le principal constituant du \textbf{gaz naturel} (environ 90 \%), exploité notamment au large de Kribi et de Limbe pour alimenter la centrale thermique de Kribi. On l'appelle aussi \emph{gaz des marais} car il se dégage des eaux stagnantes où la matière organique se décompose sans oxygène (méthanogénèse) : c'est lui qui s'enflamme parfois au-dessus des marécages (les fameux « feux follets »). Il se forme également dans les décharges et les biodigesteurs : en milieu rural camerounais, la \textbf{méthanisation} des déchets de ferme et des résidus de cuisine permet de produire un biogaz utilisable pour la cuisson, une alternative précieuse au bois de chauffe. Attention toutefois : le méthane est un gaz à effet de serre environ 25 fois plus puissant que le dioxyde de carbone sur 100 ans !
\end{savaistu}

\begin{exercice}[À toi de jouer]
On désire recueillir \SI{11,2}{L} de méthane dans les CNTP.
\begin{enumerate}
\item Quelle quantité de matière de méthane cela représente-t-il ?
\item Quelle masse de carbure d'aluminium faut-il utiliser, l'eau étant en excès ?
\item Pourquoi ne peut-on pas recueillir le méthane par déplacement d'air dans une éprouvette placée ouverture vers le haut ?
\end{enumerate}
\end{exercice}

\begin{corrige}[Exercice]
\begin{enumerate}
\item $n(\ce{CH4}) = \dfrac{V}{V_m} = \dfrac{11,2}{22,4} = \SI{0,500}{mol}$.
\item D'après l'équation $\ce{Al4C3 + 12 H2O -> 4 Al(OH)3 + 3 CH4}$, $n(\ce{Al4C3}) = \dfrac{n(\ce{CH4})}{3} = \dfrac{0,500}{3} = \SI{0,167}{mol}$ (soit $\frac{1}{6}$ mol). La masse nécessaire est $m = n \times M = \frac{1}{6} \times 144 = \SI{24,0}{g}$.
\item Le méthane est moins dense que l'air ($M = \SI{16}{g.mol^{-1}} < \SI{29}{g.mol^{-1}}$) : placée ouverture vers le haut, l'éprouvette laisserait le gaz s'échapper vers le haut. Il faudrait la placer ouverture vers le bas pour un recueillement par déplacement d'air ; mais le recueillement par déplacement d'eau reste préférable car il donne un gaz plus pur (sans vapeur d'eau significative) et permet de doser le volume recueilli.
\end{enumerate}
\end{corrige}

\begin{aretenir}[L'essentiel de la section]
\begin{itemize}
\item Le méthane se prépare au laboratoire par action de l'eau sur le carbure d'aluminium : $\ce{Al4C3 + 12 H2O -> 4 Al(OH)3 + 3 CH4}$, réaction totale, lente, à température ordinaire.
\item Le dispositif comprend un ballon contenant le carbure, une ampoule de coulée d'eau et un recueillement par déplacement d'eau, car le méthane est très peu soluble dans l'eau.
\item Le méthane est un gaz incolore, inodore, inflammable, moins dense que l'air ; il brûle avec une flamme bleue en donnant de l'eau et du \ce{CO2} qui trouble l'eau de chaux.
\item Sécurité : éloigner toute flamme et vérifier l'étanchéité du montage ; purger l'air avant recueillement.
\item Ne pas confondre \ce{Al4C3} (donne le méthane) et \ce{CaC2} (donne l'acétylène).
\end{itemize}
\end{aretenir}

\newpage

\coursec{Travaux pratiques}

\courssub{Objectifs du TP}

À l'issue de cette séance de travaux pratiques, tu dois être capable de :

\begin{itemize}
\item construire, à l'aide de \cle{modèles moléculaires}, les molécules de méthane et d'éthane, et visualiser le \cle{carbone tétragonal} (angles de \SI{109,5}{\degree}) ;
\item observer la \cle{rotation libre} autour de la liaison simple C--C de l'éthane et représenter ses conformations (éclipsée et alternée) ;
\item construire \textbf{tous les isomères de chaîne} de formules brutes \ce{C4H10} et \ce{C5H12}, écrire leurs formules semi-développées et les nommer correctement selon la nomenclature IUPAC ;
\item décrire et schématiser le dispositif de \cle{préparation du méthane} à partir du carbure d'aluminium ;
\item identifier le méthane par ses propriétés : combustion à la flamme bleue non fuligineuse, formation d'eau (condensation) et de dioxyde de carbone (trouble de l'eau de chaux).
\end{itemize}

\courssub{Matériel et produits}

\begin{center}
\rowcolors{2}{popblueL}{white}
\begin{tabularx}{\linewidth}{|l|X|}
\hline
\rowcolor{popblue!40} \textbf{Matériel / produit} & \textbf{Quantité et remarques} \\
\hline
Modèles moléculaires (boules et tiges) & 1 boîte par groupe de 4 élèves \\
\hline
Boules de pâte à modeler (2 couleurs) et allumettes & \emph{Alternative locale} : noir pour le carbone, blanc ou rouge pour l'hydrogène ; les allumettes (tête coupée) jouent le rôle des liaisons \\
\hline
Carbure d'aluminium \ce{Al4C3} & environ \SI{5}{g} (poudre grise) \\
\hline
Ballon à fond rond à col rodé & 1 (\SI{250}{mL}) \\
\hline
Ampoule de coulée à robinet & 1, remplie d'eau distillée (\SI{50}{mL}) \\
\hline
Cristallisoir ou grande cuvette & 1, rempli d'eau aux trois quarts \\
\hline
Éprouvettes graduées (\SI{250}{mL}) & 4 à 6, pour recueillir le gaz par déplacement d'eau \\
\hline
Bouchon à deux trous + tubes à dégagement & 1 jeu, avec raccords en caoutchouc \\
\hline
Bec Bunsen ou allumettes & pour l'essai de combustion \\
\hline
Eau de chaux & quelques mL dans un flacon \\
\hline
Verre de montre ou coupelle froide & pour condenser la vapeur d'eau \\
\hline
Lunettes, gants, blouse & obligatoires pour chaque élève \\
\hline
\end{tabularx}
\end{center}

\begin{attention}[Sécurité : les gestes qui sauvent]
\begin{itemize}
\item \textbf{Carbure d'aluminium} : poudre qui réagit violemment avec l'eau en dégageant un gaz \emph{inflammable}. Pictogramme « \emph{matière inflammable} » (flamme noire sur fond rouge) et « \emph{corrosif} » (contact avec l'eau). Manipuler avec une spatule sèche, loin de toute flamme, et ne jamais en verser dans l'évier.
\item \textbf{Méthane} : gaz \emph{extrêmement inflammable} ; son mélange avec l'air est \emph{explosif} (limites d'explosivité : 5--15 \% en volume). Pictogramme « \emph{gaz sous pression / inflammable} ». Aucune flamme ne doit approcher le ballon pendant le dégagement ; on n'allume le gaz \textbf{qu'après avoir vérifié sa pureté} (essai à la petite éprouvette : une légère détonation sourde indique un mélange encore impur, un « pop » discret ou une combustion calme indique un gaz pur).
\item \textbf{Eau de chaux} : solution irritante. Pictogramme « \emph{point d'exclamation} ». En cas de contact avec la peau ou les yeux, rincer abondamment à l'eau claire.
\item Porter en permanence \textbf{blouse, lunettes et gants} ; attacher les cheveux longs ; ne jamais goûter ni sentir directement un gaz (éventer avec la main si nécessaire).
\item L'ampoule de coulée doit être \textbf{fermée} avant toute manipulation du ballon ; ajouter l'eau \emph{goutte à goutte}, jamais en un seul jet.
\end{itemize}
\end{attention}

\courssub{Schéma du dispositif}

\begin{popfigure}
\begin{center}
\begin{tikzpicture}[scale=0.95, every node/.style={font=\small}]
% Cristallisoir
\draw[thick, popblue] (6.2,-2.6) -- (6.2,0.2) -- (11.8,0.2) -- (11.8,-2.6);
\draw[thick, popblue] (6.2,-2.6) -- (11.8,-2.6);
% Eau du cristallisoir
\fill[popblue!20] (6.25,-2.55) rectangle (11.75,-0.5);
\draw[popblue!60] (6.25,-0.5) -- (11.75,-0.5);
\node[popblue!80] at (10.9,-2.2) {eau};
% Éprouvette renversée
\draw[thick, popdark] (7.6,-2.2) -- (7.6,1.4) -- (8.6,1.4) -- (8.6,-2.2);
\fill[popblue!35] (7.65,-2.15) rectangle (8.55,-0.2);
\node[popdark] at (8.1,1.7) {éprouvette};
\node[popdark, align=center] at (8.1,-1.1) {\footnotesize méthane};
% Tube à dégagement coudé
\draw[thick, popdark] (2.6,1.6) -- (4.6,1.6) -- (4.6,-1.6) -- (8.1,-1.6) -- (8.1,-2.05);
% Ballon
\draw[thick, popdark] (0,0) circle (1.15);
\draw[thick, popdark] (-0.35,1.05) -- (-0.35,2.0) -- (0.35,2.0) -- (0.35,1.05);
% Solide dans le ballon
\fill[poppurple!50] (-0.55,-0.75) .. controls (-0.2,-0.5) and (0.2,-0.95) .. (0.55,-0.7) -- (0.55,-1.05) -- (-0.55,-1.05) -- cycle;
\node[poppurple] at (0,-0.85) {\footnotesize \ce{Al4C3}};
% Bulles
\foreach \p in {(-0.3,-0.2),(0.15,0.1),(-0.1,0.45),(0.35,0.6),(0.05,0.8)} \draw[popblue] \p circle (0.05);
% Bouchon
\fill[popgold!70] (-0.38,1.55) rectangle (0.38,1.95);
% Ampoule de coulée
\draw[thick, popdark] (-1.6,2.0) -- (-1.6,3.6) -- (-0.9,3.6) -- (-0.9,2.0);
\draw[thick, popdark] (-1.6,3.6) .. controls (-1.9,4.1) and (-0.6,4.1) .. (-0.9,3.6);
\draw[thick, popdark] (-1.25,2.0) -- (-1.25,1.55);
\fill[popblue!30] (-1.55,2.6) rectangle (-0.95,3.55);
\draw[thick, popdark] (-1.45,2.35) -- (-1.05,2.35);
\node[popdark, align=center] at (-2.9,3.2) {ampoule de\\coulée (eau)};
% Liaison ampoule-bouchon
\draw[thick, popdark] (-1.25,1.55) -- (-1.25,1.75) -- (-0.38,1.75);
% Tube sortie bouchon
\draw[thick, popdark] (0.2,1.95) -- (0.2,2.3) -- (2.6,2.3) -- (2.6,1.6);
% Étiquettes
\node[popdark, align=center] at (0,-1.9) {ballon};
\node[popdark, align=center] at (4.6,2.0) {\footnotesize tube à dégagement};
\node[popdark, align=center] at (9.9,0.9) {cristallisoir};
% Flèches indicatives
\draw[-{Stealth[length=2mm]}, poporange, thick] (0.5,0.5) -- (1.5,1.0) node[right, poporange] {\footnotesize dégagement \ce{CH4}};
\draw[-{Stealth[length=2mm]}, popblue, thick] (8.1,-1.6) -- (8.1,-2.4) node[below, popblue] {\footnotesize recueillement};
\end{tikzpicture}
\end{center}
\legende{Dispositif de préparation du méthane par action de l'eau sur le carbure d'aluminium, avec recueillement du gaz par déplacement d'eau.}
\end{popfigure}

\courssub{Protocole}

\textbf{Partie 1 : construction des modèles moléculaires (par groupes de 4)}

\begin{enumerate}
\item Construire le modèle du \textbf{méthane} \ce{CH4} : une boule noire (carbone) au centre, quatre tiges dirigées vers les sommets d'un tétraèdre régulier, quatre boules claires (hydrogènes). Vérifier l'angle entre deux liaisons : il vaut \SI{109,5}{\degree}.
\item Construire le modèle de l'\textbf{éthane} \ce{C2H6}. Faire tourner doucement un groupe \ce{CH3} par rapport à l'autre autour de la liaison C--C : observer que la rotation est \emph{libre} (faible barrière d'énergie). Arrêter la rotation dans la position où les hydrogènes des deux carbones se font face (\emph{conformation éclipsée}), puis dans celle où ils sont décalés de \SI{60}{\degree} (\emph{conformation alternée} ou décalée, plus stable). Dessiner les deux vues de Newman (ou perspective Cram simplifiée).
\item Construire le \textbf{propane} \ce{C3H8}, puis chercher toutes les façons d'assembler 4 atomes de carbone sans créer de cycles : on obtient \textbf{deux} modèles non superposables de \ce{C4H10} (chaîne linéaire et chaîne ramifiée). Dessiner les formules semi-développées et nommer chaque isomère.
\item Recommencer avec 5 atomes de carbone : trouver les \textbf{trois} isomères de \ce{C5H12}. Pour chacun, écrire la formule semi-développée et le nom officiel.
\item Vérification croisée : un autre groupe doit pouvoir nommer tes modèles sans regarder ta feuille.
\end{enumerate}

\begin{aretenir}[Bilan structural : tétragonalité, rotation et isomérie]
\begin{itemize}
\item Le carbone des alcanes est \textbf{sp$^3$ hybridé} : ses 4 liaisons $\sigma$ pointent vers les sommets d'un tétraèdre régulier (angle \SI{109,5}{\degree}). C'est le \cle{carbone tétragonal}.
\item La liaison simple C--C autorise une \textbf{rotation libre} (barrière $\approx$ \SI{12}{kJ.mol^{-1}} pour l'éthane). L'éthane adopte principalement la \textbf{conformation alternée} (plus stable, hydrogènes écartés) ; la conformation éclipsée est un maximum d'énergie.
\item À partir de \textbf{4 atomes de carbone}, une même formule brute \ce{C_nH_{2n+2}} peut correspondre à plusieurs squelettes carbonés distincts : ce sont les \cle{isomères de chaîne} (aussi appelés isomères de constitution). Ils ont des propriétés physiques différentes (points d'ébullition, densités).
\end{itemize}
\end{aretenir}

\textbf{Partie 2 : préparation et identification du méthane}

\begin{enumerate}
\item Introduire environ \SI{5}{g} de carbure d'aluminium en poudre dans le ballon sec. Boucher le ballon avec le bouchon à deux trous muni de l'ampoule de coulée et du tube à dégagement. Vérifier l'étanchéité.
\item Remplir l'ampoule de coulée d'environ \SI{50}{mL} d'eau distillée et vérifier que le robinet est fermé.
\item Remplir les éprouvettes d'eau et les retourner dans le cristallisoir, sans laisser entrer d'air (les boucher sous l'eau avec le pouce avant de les retourner). Placer l'extrémité du tube à dégagement sous l'ouverture de la première éprouvette.
\item Ouvrir délicatement le robinet pour laisser couler l'eau \textbf{goutte à goutte} (environ 1 goutte par seconde). Observer le dégagement gazeux. \emph{Laisser s'échapper les premières bulles} (air du ballon et des tuyaux) pendant \SIrange{1}{2}{\minute} avant de recueillir le gaz.
\item Recueillir 2 à 3 éprouvettes pleines de gaz. Les boucher sous l'eau, les sortir et les placer debout sur la paillasse.
\item \textbf{Essai de combustion} : approcher une flamme (allumette ou bec Bunsen) de l'ouverture de la première éprouvette. Observer la couleur de la flamme et l'éventuel bruit.
\item Placer une coupelle froide (ou un verre de montre) au-dessus de la flamme de la deuxième éprouvette : observer le dépôt sur la paroi froide.
\item Après extinction, ajouter quelques millilitres d'eau de chaux dans la deuxième éprouvette, boucher avec le pouce (ganté) et agiter : observer.
\item Recommencer les essais 6 à 8 avec les éprouvettes restantes si nécessaire.
\end{enumerate}

\courssub{Observations et résultats attendus}

\begin{center}
\rowcolors{2}{popgreenL}{white}
\begin{tabularx}{\linewidth}{|X|X|}
\hline
\rowcolor{popgreen!45} \textbf{Étape} & \textbf{Observation attendue} \\
\hline
Addition d'eau sur \ce{Al4C3} & Dégagement gazeux régulier, effervescence, léger échauffement du ballon (réaction exothermique) ; il reste un résidu blanchâtre gélatineux (hydroxyde d'aluminium \ce{Al(OH)3}) \\
\hline
Recueillement & Gaz incolore, inodore, très peu soluble dans l'eau (le niveau d'eau descend dans l'éprouvette) ; volume total théorique pour \SI{5}{g} \ce{Al4C3} : $\approx$ \SI{2,5}{L} (débit contrôlé par l'ampoule) \\
\hline
Essai à la flamme & Le gaz brûle avec une \textbf{flamme bleue} pâle, non fuligineuse ; « pop » discret si le gaz est pur (combustion calme), détonation sourde si mélange air/méthane \\
\hline
Coupelle froide au-dessus de la flamme & Formation de \textbf{buée} (fines gouttelettes d'eau liquide) qui trouble le verre : preuve de formation de \ce{H2O} \\
\hline
Eau de chaux après combustion & L'eau de chaux se \textbf{trouble} (précipité blanc laiteux de carbonate de calcium \ce{CaCO3}) : preuve de formation de \ce{CO2} \\
\hline
Modèles moléculaires & 1 modèle pour \ce{CH4}, \ce{C2H6}, \ce{C3H8} ; \textbf{2 isomères} pour \ce{C4H10} ; \textbf{3 isomères} pour \ce{C5H12} \\
\hline
\end{tabularx}
\end{center}

\courssub{Exploitation}

\begin{exoresolu}[Questions guidées et réponses]
\textbf{Question 1.} Écrire l'équation-bilan de la réaction entre le carbure d'aluminium et l'eau, sachant qu'il se forme aussi de l'hydroxyde d'aluminium \ce{Al(OH)3}.

\textbf{Question 2.} Pourquoi recueille-t-on le méthane par déplacement d'eau ?

\textbf{Question 3.} Écrire l'équation-bilan de la combustion complète du méthane dans le dioxygène et justifier les deux tests d'identification des produits.

\textbf{Question 4.} Nommer les deux isomères de \ce{C4H10} et les trois isomères de \ce{C5H12} en donnant leurs formules semi-développées.

\textbf{Question 5.} Pourquoi faut-il laisser s'échapper les premières bulles avant de recueillir le gaz ?
\tcblower
\textbf{Réponse 1.} L'eau attaque le carbure d'aluminium par hydrolyse. L'équation-bilan équilibrée est :
\[ \ce{Al4C3 + 12 H2O -> 4 Al(OH)3 + 3 CH4} \]
Le méthane est le seul produit gazeux : c'est lui qui se dégage. Le résidu solide est l'hydroxyde d'aluminium.

\textbf{Réponse 2.} Le méthane est \emph{très peu soluble dans l'eau} (\SI{\approx 22}{mg.L^{-1}} à \SI{20}{\celsius}) et sa densité (\SI{0,72}{g.L^{-1}}) est très inférieure à celle de l'air : le recueillement par déplacement d'eau est donc la méthode la plus simple et la plus sûre au laboratoire pour obtenir un gaz pur.

\textbf{Réponse 3.} Combustion complète dans le dioxygène (excès d'air) :
\[ \ce{CH4 + 2 O2 -> CO2 + 2 H2O} \]
\begin{itemize}
\item L'\textbf{eau} formée à l'état gazeux se condense au contact de la paroi froide : apparition de buée (test physique).
\item Le \textbf{dioxyde de carbone} réagit avec l'eau de chaux (solution de \ce{Ca(OH)2}) :
\[ \ce{CO2 + Ca(OH)2 -> CaCO3 v + H2O} \]
Le précipité blanc de carbonate de calcium \ce{CaCO3} est caractéristique (test chimique). Excès de \ce{CO2} : le trouble disparaît (formation de \ce{Ca(HCO3)2} soluble).
\end{itemize}

\textbf{Réponse 4.}
\begin{itemize}
\item \ce{C4H10} :
      \begin{enumerate}
      \item \ce{CH3-CH2-CH2-CH3} : \textbf{butane} (ou n-butane).
      \item \ce{CH3-CH(CH3)-CH3} : \textbf{2-méthylpropane} (ou isobutane).
      \end{enumerate}
\item \ce{C5H12} :
      \begin{enumerate}
      \item \ce{CH3-CH2-CH2-CH2-CH3} : \textbf{pentane} (ou n-pentane).
      \item \ce{CH3-CH(CH3)-CH2-CH3} : \textbf{2-méthylbutane}.
      \item \ce{CH3-C(CH3)2-CH3} : \textbf{2,2-diméthylpropane} (ou néopentane).
      \end{enumerate}
\end{itemize}

\textbf{Réponse 5.} Les premières bulles chassent l'\emph{air} contenu dans le ballon et les tubes. Recueillir trop tôt donnerait un mélange méthane--air, \emph{explosif} à la flamme (risque d'éclatement de l'éprouvette), et fausserait l'essai de combustion (flamme jaune, suie, formation de \ce{CO}).
\end{exoresolu}

\begin{attention}[Erreurs fréquentes à éviter dans le compte rendu]
\begin{itemize}
\item Ne pas écrire \ce{Al4C3 + H2O -> CH4} sans équilibrer : toute équation-bilan doit être \textbf{équilibrée} en atomes (et en charges si ions).
\item Ne pas confondre \emph{isomères de chaîne} (même formule brute, squelettes carbonés différents) avec de simples rotations du modèle ou conformations : deux modèles qui se superposent après rotation autour de liaisons simples représentent la \textbf{même} molécule.
\item La flamme du méthane pur en excès d'air est \emph{bleue} ; une flamme jaune et fuligineuse (dépôt de carbone) trahirait une combustion \emph{incomplète} (manque d'\ce{O2}), formant du monoxyde de carbone \ce{CO} (gaz toxique, incolore, inodore) ou du carbone \ce{C}.
\item Ne pas oublier l'état physique de l'eau formée : à la flamme, c'est de la \emph{vapeur d'eau} (gaz) ; elle se condense en \emph{eau liquide} (buée) sur la coupelle froide.
\end{itemize}
\end{attention}

\courssub{Conclusion}

Cette séance t'a permis de concrétiser deux idées majeures de la leçon sur les alcanes. D'une part, le carbone des alcanes est \cle{tétragonal} (\emph{sp$^3$}) : ses quatre liaisons $\sigma$ pointent vers les sommets d'un tétraèdre régulier (angle \SI{109,5}{\degree}), et la rotation libre autour des liaisons C--C simples explique l'existence de conformations (éclipsée, alternée) pour l'éthane et les alcanes supérieurs. D'autre part, dès quatre atomes de carbone, une même formule brute \ce{C_nH_{2n+2}} peut correspondre à plusieurs composés distincts : ce sont les \cle{isomères de chaîne} (2 pour \ce{C4H10}, 3 pour \ce{C5H12}, 5 pour \ce{C6H14}...), que seule une nomenclature rigoureuse (IUPAC) permet de distinguer et de nommer sans ambiguïté. Enfin, tu as préparé le méthane par action de l'eau sur le carbure d'aluminium (hydrolyse) et tu l'as identifié par sa combustion complète, qui produit de l'eau et du dioxyde de carbone. Ce gaz, principal constituant du gaz naturel (exploité au Cameroun par la SONARA/Gaz du Cameroun), est de la même famille chimique que le butane des bouteilles de gaz utilisées dans nos cuisines : les alcanes sont bien des composés de la vie quotidienne et industrielle, dont la maîtrise exige rigueur, respect des règles de sécurité et compréhension fine de leur structure tridimensionnelle.

\newpage

\coursec{Méthodes et exercices résolus}

\courssub{Méthode 1 : Écrire et nommer un alcane (nomenclature)}

\begin{methode}[Nommer un alcane à partir de sa formule semi-développée]
\begin{enumerate}
\item \textbf{Repérer la chaîne la plus longue} contenant le maximum d'atomes de carbone : c'est la chaîne principale. Son nombre de carbones donne la racine : méth-, éth-, prop-, but-, pent-, hex-...
\item \textbf{Numéroter} cette chaîne dans le sens qui donne aux ramifications les \cle{indices de position les plus petits} (comparer les suites d'indices chiffre par chiffre).
\item \textbf{Identifier les ramifications} (alkyles) : \ce{-CH3} méthyle, \ce{-C2H5} éthyle, \ce{-C3H7} propyle.
\item \textbf{Écrire le nom} : indice(s)-ramification(s) + nom de la chaîne principale + suffixe \og ane \fg{}. Si une même ramification apparaît plusieurs fois, utiliser les préfixes di-, tri-, tétra- et répéter les indices séparés par des virgules.
\item \textbf{Vérifier} : les indices sont séparés des mots par des tirets ; les ramifications sont citées par ordre alphabétique (éthyle avant méthyle).
\end{enumerate}
\textbf{Réflexe} : si deux numérotations donnent les mêmes indices, choisir celle qui donne l'indice le plus petit à la ramification citée en premier dans l'ordre alphabétique.
\end{methode}

\begin{exoresolu}[Nommer un alcane ramifié]
On considère l'alcane de formule semi-développée :
\[ \chemfig{CH_3-CH(-[2]CH_3)-CH_2-CH(-[6]CH_3)-CH_3} \]
\begin{enumerate}
\item Déterminer sa formule brute.
\item Donner son nom systématique.
\item Écrire la formule semi-développée et donner le nom de l'isomère de chaîne possédant une ramification éthyle sur une chaîne de 4 carbones.
\end{enumerate}
\tcblower
\textbf{1. Formule brute.} Comptons les atomes de carbone : la chaîne horizontale en contient 5, plus 2 groupes méthyle, soit $n = 7$ carbones. Pour un alcane, $C_nH_{2n+2}$ donne $2 \times 7 + 2 = 16$ hydrogènes.
\[ \boxed{\ce{C7H16}} \]

\textbf{2. Nom systématique.} La chaîne la plus longue compte 5 carbones : racine \og pent \fg{}, suffixe \og ane \fg{} $\rightarrow$ pentane. Numérotation :
\begin{itemize}
\item de gauche à droite : ramifications en positions 2 et 4 ;
\item de droite à gauche : ramifications en positions 2 et 4.
\end{itemize}
Les deux sens donnent la même suite d'indices $(2,4)$. Les deux ramifications sont des méthyles : on écrit \og 2,4-diméthyl \fg{}.
\[ \boxed{\text{2,4-diméthylpentane}} \]

\textbf{3. Isomère à ramification éthyle.} Une chaîne de 4 carbones (butane) portant un éthyle \ce{-C2H5} : l'éthyle ne peut se fixer qu'en position 2 (la position 3 est équivalente) :
\[ \chemfig{CH_3-CH(-[2]CH_2-CH_3)-CH_2-CH_3} \]
En réalité, la chaîne la plus longue de cette molécule compte 5 carbones : le nom correct est donc \textbf{3-méthylpentane} (et non 2-éthylbutane, nom interdit car la chaîne principale choisie n'est pas la plus longue). Vérification : \ce{C6H14}... Attention : comptons : 5 C dans la chaîne + 1 C du méthyle = 6 C, soit \ce{C6H14}. Ce n'est pas un isomère de \ce{C7H16} ! Pour obtenir un isomère de \ce{C7H16} avec une ramification éthyle, il faut une chaîne de 5 carbones + un éthyle :
\[ \chemfig{CH_3-CH_2-CH(-[2]CH_2-CH_3)-CH_2-CH_3} \]
soit le \textbf{3-éthylpentane}, de formule brute \ce{C7H16}. C'est bien un isomère de constitution du 2,4-diméthylpentane.

\textbf{Conclusion} : l'alcane étudié est le 2,4-diméthylpentane, \ce{C7H16} ; un de ses isomères est le 3-éthylpentane.
\end{exoresolu}

\courssub{Méthode 2 : Trouver tous les isomères d'un alcane ($n \leq 5$)}

\begin{methode}[Dénombrer et écrire les isomères de constitution]
\begin{enumerate}
\item \textbf{Écrire la chaîne linéaire} à $n$ carbones (alcane \og normal \fg{}, préfixe n-).
\item \textbf{Raccourcir la chaîne principale d'un carbone} et placer le carbone retiré comme méthyle sur tous les carbones \emph{non équivalents} (jamais en bout de chaîne : cela redonnerait la chaîne linéaire !).
\item \textbf{Raccourcir encore} et placer deux méthyles (ou un éthyle) en explorant toutes les positions non équivalentes.
\item \textbf{Éliminer les doublons} en nommant chaque formule obtenue : deux formules portant le même nom sont le même composé.
\item \textbf{Vérifier} que chaque formule répond bien à $C_nH_{2n+2}$ et que chaque carbone forme exactement 4 liaisons.
\end{enumerate}
\textbf{Résultats à connaître} : \ce{C4H10} : 2 isomères ; \ce{C5H12} : 3 isomères ; \ce{C6H14} : 5 isomères.
\end{methode}

\begin{exoresolu}[Isomères du pentane]
Un hydrocarbure gazeux A a pour formule brute \ce{C5H12}.
\begin{enumerate}
\item Montrer que A est un alcane.
\item Écrire les formules semi-développées de tous les isomères de A et les nommer.
\item Réaliser à l'aide de modèles moléculaires (ou décrire) la représentation spatiale de l'isomère le plus ramifié : quelle est la géométrie autour du carbone central ?
\end{enumerate}
\tcblower
\textbf{1. Nature de A.} Pour $n = 5$ : $2n + 2 = 12$. La formule \ce{C5H12} répond à la formule générale des alcanes $C_nH_{2n+2}$ : A est bien un alcane (le pentane et ses isomères).

\textbf{2. Recherche systématique des isomères.}

\emph{Étape 1 — chaîne à 5 carbones :}
\[ \chemfig{CH_3-CH_2-CH_2-CH_2-CH_3} \qquad \text{pentane (n-pentane)} \]

\emph{Étape 2 — chaîne à 4 carbones + 1 méthyle :} le méthyle ne peut se placer qu'en position 2 (les positions 1 et 4 redonneraient le pentane ; les positions 2 et 3 sont équivalentes) :
\[ \chemfig{CH_3-CH(-[2]CH_3)-CH_2-CH_3} \qquad \text{2-méthylbutane} \]

\emph{Étape 3 — chaîne à 3 carbones + 2 méthyles :} les deux méthyles se placent obligatoirement sur le carbone central :
\[ \chemfig{C(-[2]CH_3)(-[4]CH_3)(-[6]CH_3)-CH_3} \qquad \text{2,2-diméthylpropane (néopentane)} \]

Vérification : chaque formule contient 5 C et 12 H, et chaque carbone est tétravalent. Il n'existe pas d'autre possibilité.
\[ \boxed{\text{\ce{C5H12} possède exactement 3 isomères de constitution.}} \]

\textbf{3. Géométrie du néopentane.} Le carbone central est lié à 4 atomes de carbone. Comme tout carbone tétragonal, ses 4 liaisons pointent vers les sommets d'un \cle{tétraèdre régulier}, avec des angles de \SI{109,5}{\degree}. Le modèle moléculaire montre une molécule quasi sphérique, très compacte : c'est pourquoi le néopentane bout à \SI{9,5}{\celsius} seulement, contre \SI{36}{\celsius} pour le n-pentane (molécule allongée, contacts intermoléculaires plus étendus).

\textbf{Conclusion} : \ce{C5H12} admet trois isomères : le pentane, le 2-méthylbutane et le 2,2-diméthylpropane, dont la géométrie tétraédrique illustre le caractère tétragonal du carbone.
\end{exoresolu}

\courssub{Méthode 3 : Exploiter la combustion d'un alcane}

\begin{methode}[Équilibrer et exploiter une combustion complète]
\begin{enumerate}
\item \textbf{Écrire l'équation générale} :
\[ \ce{C_nH_{2n+2} + \frac{3n+1}{2} O2 -> n CO2 + (n+1) H2O} \]
Si le coefficient de \ce{O2} est fractionnaire, multiplier toute l'équation par 2.
\item \textbf{Vérifier la conservation} des atomes de C, H et O de part et d'autre.
\item \textbf{Construire un tableau d'avancement} si les quantités des deux réactifs sont données, afin d'identifier le réactif limitant : la combustion complète exige du dioxygène en excès.
\item \textbf{Calculer} avec $n = \dfrac{m}{M}$ pour les masses, $n = \dfrac{V}{V_m}$ pour les gaz (préciser les conditions : à \SI{0}{\celsius} sous \SI{1}{atm}, $V_m = \SI{22,4}{L.mol^{-1}}$).
\item \textbf{Conclure} par une phrase répondant exactement à la question posée, avec unités et chiffres significatifs cohérents.
\end{enumerate}
\textbf{Réflexe} : combustion \emph{incomplète} (manque de \ce{O2}) $\rightarrow$ formation de \ce{CO} (toxique !) ou de carbone \ce{C} (suie). C'est le danger des braseros et des groupes électrogènes dans les pièces mal aérées.
\end{methode}

\begin{exoresolu}[Combustion du butane d'une bouteille de gaz domestique]
Le gaz domestique vendu à Douala et Yaoundé est essentiellement du butane \ce{C4H10}. On brûle complètement une masse $m = \SI{11,6}{g}$ de butane dans le dioxygène de l'air.
\begin{enumerate}
\item Écrire l'équation équilibrée de la combustion complète du butane.
\item Calculer la quantité de butane brûlé.
\item Déterminer le volume de dioxygène nécessaire et le volume de dioxyde de carbone dégagé, mesurés dans les CNTP ($V_m = \SI{22,4}{L.mol^{-1}}$).
\item Quelle masse d'eau obtient-on ?
\end{enumerate}
On donne : $M(\text{C}) = \SI{12}{g.mol^{-1}}$, $M(\text{H}) = \SI{1}{g.mol^{-1}}$, $M(\text{O}) = \SI{16}{g.mol^{-1}}$.
\tcblower
\textbf{1. Équation de combustion.} Avec $n = 4$ : coefficient de \ce{O2} : $\dfrac{3 \times 4 + 1}{2} = \dfrac{13}{2}$. On multiplie par 2 :
\[ \boxed{\ce{2 C4H10 + 13 O2 -> 8 CO2 + 10 H2O}} \]
Vérification : C : $2 \times 4 = 8$ à gauche, 8 à droite ; H : 20 de chaque côté ; O : 26 de chaque côté. L'équation est équilibrée.

\textbf{2. Quantité de butane.} $M(\ce{C4H10}) = 4 \times 12 + 10 \times 1 = \SI{58}{g.mol^{-1}}$.
\[ n(\ce{C4H10}) = \frac{m}{M} = \frac{11,6}{58} = \SI{0,200}{mol} \]

\textbf{3. Volumes gazeux.} Tableau d'avancement (en mol), le dioxygène étant en excès :
\begin{center}
\begin{tabular}{lcccc}
\toprule
 & \ce{2 C4H10} & \ce{+ 13 O2} & \ce{-> 8 CO2} & \ce{+ 10 H2O} \\
\midrule
État initial & $0,200$ & $n_0$ & $0$ & $0$ \\
En cours & $0,200 - 2x$ & $n_0 - 13x$ & $8x$ & $10x$ \\
État final ($x_{max} = 0,100$) & $0$ & $n_0 - 1,30$ & $0,800$ & $1,00$ \\
\bottomrule
\end{tabular}
\end{center}
Le réactif limitant est le butane : $x_{max} = \dfrac{0,200}{2} = \SI{0,100}{mol}$.

Dioxygène consommé : $n(\ce{O2}) = 13 \times 0,100 = \SI{1,30}{mol}$, soit :
\[ V(\ce{O2}) = 1,30 \times 22,4 = \SI{29,1}{L} \]
Dioxyde de carbone formé : $n(\ce{CO2}) = 8 \times 0,100 = \SI{0,800}{mol}$, soit :
\[ V(\ce{CO2}) = 0,800 \times 22,4 = \SI{17,9}{L} \]

\textbf{4. Masse d'eau.} $n(\ce{H2O}) = 10 \times 0,100 = \SI{1,00}{mol}$ ; $M(\ce{H2O}) = \SI{18}{g.mol^{-1}}$ :
\[ m(\ce{H2O}) = 1,00 \times 18 = \SI{18,0}{g} \]

\textbf{Conclusion} : la combustion complète de \SI{11,6}{g} de butane consomme \SI{29,1}{L} de dioxygène et dégage \SI{17,9}{L} de \ce{CO2} et \SI{18,0}{g} d'eau dans les CNTP.
\end{exoresolu}

\courssub{Méthode 4 : Décrire une réaction de substitution (halogénation)}

\begin{methode}[Écrire une halogénation d'alcane]
\begin{enumerate}
\item \textbf{Identifier les conditions} : la substitution d'un alcane par \ce{Cl2} ou \ce{Br2} exige la \cle{lumière} (rayonnement UV) ou la chaleur. Sans lumière, pas de réaction.
\item \textbf{Remplacer un H par un halogène} : un atome d'hydrogène de l'alcane est remplacé par un atome de chlore ou de brome ; l'autre atome d'halogène forme \ce{HCl} ou \ce{HBr} :
\[ C_nH_{2n+2} + \ce{Cl2} \xrightarrow{\text{lumière}} C_nH_{2n+1}Cl + \ce{HCl} \]
\item \textbf{Prévoir la poursuite de la substitution} : le produit obtenu possède encore des atomes d'hydrogène ; la réaction continue (di-, tri-, tétra-substitution) tant qu'il reste des H et de l'halogène. On obtient un \textbf{mélange} de dérivés halogénés.
\item \textbf{Nommer les produits} avec le préfixe chloro- ou bromo- et l'indice de position (ex. 1-chloropropane, 2-chloropropane).
\item \textbf{Dénombrer les isomères} monosubstitués en repérant les atomes d'hydrogène non équivalents de l'alcane de départ.
\end{enumerate}
\textbf{Réflexe} : ne jamais écrire \ce{CH4 + Cl2 -> CH3Cl} sans \ce{HCl} : la conservation des atomes l'interdit.
\end{methode}

\begin{exoresolu}[Chloration du méthane et du propane]
\textbf{Partie A.} On expose un mélange de méthane et de dichlore à la lumière.
\begin{enumerate}
\item Écrire les équations des quatre substitutions successives possibles et nommer chaque produit organique.
\item Citer une utilisation de l'un de ces dérivés.
\end{enumerate}
\textbf{Partie B.} On fait agir le dibrome sur le propane en présence de lumière, en limitant la réaction à la monosubstitution.
\begin{enumerate}
\setcounter{enumi}{2}
\item Combien de monobromopropanes différents peut-on obtenir ? Écrire leurs formules semi-développées et les nommer.
\end{enumerate}
\tcblower
\textbf{Partie A. 1. Équations des substitutions successives :}
\begin{align*}
\ce{CH4 + Cl2 ->[h\nu] CH3Cl + HCl} & \quad \text{chlorométhane} \\
\ce{CH3Cl + Cl2 ->[h\nu] CH2Cl2 + HCl} & \quad \text{dichlorométhane} \\
\ce{CH2Cl2 + Cl2 ->[h\nu] CHCl3 + HCl} & \quad \text{trichlorométhane (chloroforme)} \\
\ce{CHCl3 + Cl2 ->[h\nu] CCl4 + HCl} & \quad \text{tétrachlorométhane}
\end{align*}
À chaque étape, un H est remplacé par un Cl et une molécule de \ce{HCl} se forme. Le mélange final contient les quatre dérivés, qu'il faut séparer par distillation.

\textbf{2. Utilisation.} Le dichlorométhane \ce{CH2Cl2} est un solvant très utilisé (décaféination, extraction en laboratoire) ; le chloroforme \ce{CHCl3} fut le premier anesthésique chirurgical. Le tétrachlorométhane, autrefois solvant de détachage, est désormais proscrit car toxique et destructeur de la couche d'ozone.

\textbf{Partie B. 3. Monobromopropanes.} Le propane \chemfig{CH_3-CH_2-CH_3} possède deux types d'hydrogènes non équivalents : les 6 H des deux groupes \ce{CH3} (équivalents) et les 2 H du groupe \ce{CH2} central. On obtient donc \textbf{deux} monobromopropanes :
\begin{center}
\chemfig{CH_3-CH_2-CH_2-Br} \qquad \text{1-bromopropane} \\
\vspace{0,5em}
\chemfig{CH_3-CH(-[2]Br)-CH_3} \qquad \text{2-bromopropane}
\end{center}
Équation (pour le 2-bromopropane) :
\[ \ce{C3H8 + Br2 ->[\text{lumière}] C3H7Br + HBr} \]

\textbf{Conclusion} : la monobromation du propane donne un mélange de deux isomères de position, le 1-bromopropane et le 2-bromopropane, accompagnés de \ce{HBr}.
\end{exoresolu}

\courssub{Méthode 5 : Décrire la préparation du méthane au laboratoire}

\begin{methode}[Préparer et recueillir le méthane à partir du carbure d'aluminium]
\begin{enumerate}
\item \textbf{Retenir le principe} : l'hydrolyse du carbure d'aluminium \ce{Al4C3} par l'eau dégage du méthane :
\[ \ce{Al4C3 + 12 H2O -> 4 Al(OH)3 + 3 CH4} \]
\item \textbf{Schématiser le dispositif} : un flacon (ou tube à dégagement) contenant le carbure d'aluminium, surmonté d'un entonnoir à robinet (ampoule de coulée) pour verser l'eau goutte à goutte ; le gaz est conduit par un tube à dégagement.
\item \textbf{Recueillir le gaz} : le méthane est très peu soluble dans l'eau et moins dense que l'air ; on le recueille par \cle{déplacement d'eau} (éprouvette renversée sur une cuve à eau).
\item \textbf{Caractériser le produit} : le méthane brûle avec une flamme bleue (combustion) ; le \ce{CO2} formé trouble l'eau de chaux.
\item \textbf{Sécurité} : éloigner toute flamme du dispositif pendant le dégagement ; le mélange méthane-air est explosif (grisou des mines).
\end{enumerate}
\end{methode}

\begin{exoresolu}[Hydrolyse du carbure d'aluminium]
Au laboratoire du lycée, on désire préparer un volume $V = \SI{1,12}{L}$ de méthane (CNTP) par action de l'eau sur le carbure d'aluminium \ce{Al4C3}.
\begin{enumerate}
\item Écrire l'équation de la réaction.
\item Schématiser le dispositif expérimental et justifier le mode de recueillement du gaz.
\item Calculer la masse minimale de carbure d'aluminium à utiliser.
\end{enumerate}
On donne : $M(\text{Al}) = \SI{27}{g.mol^{-1}}$, $M(\text{C}) = \SI{12}{g.mol^{-1}}$, $V_m = \SI{22,4}{L.mol^{-1}}$ (CNTP).
\tcblower
\textbf{1. Équation de la réaction :}
\[ \ce{Al4C3 + 12 H2O -> 4 Al(OH)3 + 3 CH4} \]
Vérification : Al : 4 de chaque côté ; C : 3 de chaque côté ; H : 24 de chaque côté ; O : 12 de chaque côté.

\textbf{2. Dispositif.} On place le carbure d'aluminium solide dans un ballon ; une ampoule de coulée permet d'ajouter l'eau lentement ; le gaz dégagé est recueilli dans une éprouvette graduée renversée sur une cuve à eau.
\begin{popfigure}
\begin{tikzpicture}[scale=0.9, every node/.style={font=\small}]
% Ballon
\draw[thick, popblue] (0,0) circle (1);
\draw[thick, popblue] (-0.35,0.95) -- (-0.35,1.9) -- (0.35,1.9) -- (0.35,0.95);
% Solide carbure
\fill[poppurple!40] (-0.6,-0.5) circle (0.12); \fill[poppurple!40] (0.1,-0.65) circle (0.12);
\fill[poppurple!40] (0.5,-0.35) circle (0.12); \fill[poppurple!40] (-0.15,-0.3) circle (0.12);
\node[popdark] at (0,-1.45) {\ce{Al4C3} (solide)};
% Ampoule de coulée
\draw[thick, poporange] (-0.12,1.9) -- (-0.12,3.1) -- (0.12,3.1) -- (0.12,1.9);
\draw[thick, poporange] (-0.5,3.1) -- (-0.12,3.9) -- (0.12,3.9) -- (0.5,3.1);
\fill[popblue!30] (-0.1,3.15) rectangle (0.1,3.8);
\node[popdark, right] at (0.6,3.5) {eau};
% Tube à dégagement
\draw[thick, popgreen] (0.35,1.3) .. controls (1.6,1.3) and (2.4,0.6) .. (3.2,0.2) -- (3.2,-1.3);
% Cuve à eau
\draw[thick, popblue] (2.1,-1.6) -- (2.1,-2.6) -- (5.3,-2.6) -- (5.3,-1.6);
\fill[popblue!20] (2.15,-2.55) rectangle (5.25,-1.9);
\node[popdark] at (5.9,-2.2) {eau};
% Éprouvette renversée
\draw[thick, popdark] (2.85,-0.3) -- (2.85,-2.3) (3.55,-0.3) -- (3.55,-2.3);
\draw[thick, popdark] (2.85,-0.3) -- (3.55,-0.3);
\fill[popcream] (2.87,-0.32) rectangle (3.53,-1.5);
\node[popdark, above] at (3.2,-0.25) {\ce{CH4}};
\end{tikzpicture}
\legende{Préparation du méthane par hydrolyse du carbure d'aluminium et recueillement par déplacement d'eau.}
\end{popfigure}
Le méthane est \textbf{très peu soluble dans l'eau} : on peut donc le recueillir par déplacement d'eau (il chasse l'eau de l'éprouvette au fur et à mesure de son arrivée).

\textbf{3. Masse de carbure d'aluminium.}
\[ n(\ce{CH4}) = \frac{V}{V_m} = \frac{1,12}{22,4} = \SI{5,00e-2}{mol} \]
D'après l'équation, 1 mol de \ce{Al4C3} donne 3 mol de \ce{CH4} :
\[ n(\ce{Al4C3}) = \frac{n(\ce{CH4})}{3} = \frac{5,00 \times 10^{-2}}{3} = \SI{1,67e-2}{mol} \]
Masse molaire : $M(\ce{Al4C3}) = 4 \times 27 + 3 \times 12 = 108 + 36 = \SI{144}{g.mol^{-1}}$.
\[ m(\ce{Al4C3}) = 1,67 \times 10^{-2} \times 144 = \SI{2,40}{g} \]

\textbf{Conclusion} : il faut hydrolyser au minimum \SI{2,40}{g} de carbure d'aluminium pour recueillir \SI{1,12}{L} de méthane dans les CNTP (en supposant un rendement de \SI{100}{\percent}).
\end{exoresolu}

\courssub{Méthode 6 : Représenter la géométrie et les conformations des alcanes}

\begin{methode}[Représenter un alcane dans l'espace]
\begin{enumerate}
\item \textbf{Partir du carbone tétragonal} : chaque carbone d'un alcane est au centre d'un tétraèdre régulier ; les 4 liaisons pointent vers les sommets, avec des angles de \SI{109,5}{\degree}.
\item \textbf{Utiliser la représentation de Cram} : liaison dans le plan (trait plein fin), liaison en avant du plan (trait plein élargi, triangle), liaison en arrière du plan (trait hachuré).
\item \textbf{Pour les modèles moléculaires} : boules noires (C), blanches (H) ; bâtonnets = liaisons covalentes. Vérifier toujours : 4 bâtonnets par carbone, 1 par hydrogène.
\item \textbf{Conformations} : autour d'une liaison simple \ce{C-C}, la rotation est libre ; dans l'éthane, on distingue la conformation \cle{éclipsée} (H en vis-à-vis, énergie maximale) et la conformation \cle{décalée} (énergie minimale, plus stable), représentées en projection de Newman.
\item \textbf{Pour les cyclanes} : la chaîne se referme en cycle ; la formule générale devient $C_nH_{2n}$ (ex. cyclopropane \ce{C3H6}, cyclohexane \ce{C6H12}).
\end{enumerate}
\end{methode}

\begin{exoresolu}[Modèle moléculaire et représentation de l'éthane]
Un élève construit le modèle moléculaire de l'éthane \ce{C2H6}.
\begin{enumerate}
\item Combien de boules et de bâtonnets lui faut-il ?
\item Représenter un carbone de l'éthane en convention de Cram.
\item Expliquer pourquoi l'éthane peut adopter plusieurs conformations et préciser laquelle est la plus stable.
\end{enumerate}
\tcblower
\textbf{1. Modèle moléculaire.} \ce{C2H6} contient 2 atomes de carbone et 6 atomes d'hydrogène : il faut \textbf{2 boules noires et 6 boules blanches}. Nombre de liaisons (bâtonnets) : 1 liaison \ce{C-C} et 6 liaisons \ce{C-H}, soit \textbf{7 bâtonnets}. Vérification : chaque carbone porte bien 4 bâtonnets (1 vers l'autre C, 3 vers des H), chaque hydrogène 1 seul.

\textbf{2. Représentation de Cram d'un carbone de l'éthane :}
\[ \chemfig{C(-[2]H)(-[6]H)(<[4]H)(<:[8]CH_3)} \]
Le carbone est au centre d'un tétraèdre : deux liaisons dans le plan (vers H et \ce{CH3} représentées ici dans le plan), une liaison en avant (triangle plein vers H), une liaison en arrière (hachures vers \ce{CH3} ou H selon la vue choisie). Les angles valent \SI{109,5}{\degree}.

\textbf{3. Conformations de l'éthane.} La liaison \ce{C-C} est une liaison simple ($\sigma$) : les deux groupes \ce{CH3} peuvent \textbf{tourner librement} l'un par rapport à l'autre autour de cet axe. La molécule adopte alors une infinité de conformations, dont deux extrêmes :
\begin{itemize}
\item la conformation \textbf{éclipsée} : les hydrogènes des deux carbones se font face ; les répulsions entre doublets liants sont maximales $\rightarrow$ conformation la moins stable ;
\item la conformation \textbf{décalée} : les hydrogènes sont intercalés (rotation de \SI{60}{\degree}) ; les répulsions sont minimales $\rightarrow$ \textbf{conformation la plus stable}.
\end{itemize}
À température ordinaire, la rotation est si rapide qu'on ne peut pas isoler ces conformations : ce ne sont \emph{pas} des isomères.

\textbf{Conclusion} : l'éthane est une molécule flexible autour de sa liaison \ce{C-C} ; sa conformation décalée, d'énergie minimale, est la plus stable.
\end{exoresolu}

\begin{attention}[Les 5 erreurs qui coûtent des points au Probatoire]
\begin{enumerate}
\item \textbf{Oublier \ce{HCl} (ou \ce{HBr}) dans une substitution.} Écrire \ce{CH4 + Cl2 -> CH3Cl} est faux : l'équation n'est pas équilibrée. Toute substitution d'un alcane produit un dérivé halogéné \emph{et} un halogénure d'hydrogène.
\item \textbf{Choisir une chaîne principale qui n'est pas la plus longue.} Le nom \og 2-éthylbutane \fg{} est interdit : la vraie chaîne principale compte 5 carbones (3-méthylpentane). Cherche toujours la chaîne la plus longue, même si elle \og monte \fg{} dans une ramification.
\item \textbf{Mal numéroter la chaîne.} Compare les suites d'indices chiffre par chiffre : $(2,3)$ l'emporte sur $(2,4)$, et non sur la somme des indices. Et n'oublie pas les préfixes di-, tri- avec répétition des indices : 2,2-diméthylpropane.
\item \textbf{Confondre isomères et conformations.} Les conformations éclipsée et décalée de l'éthane ne sont \emph{pas} des isomères : ce sont des états de rotation autour d'une liaison simple. Les isomères ont des formules développées différentes (enchaînement des atomes différent).
\item \textbf{Négliger les conditions opératoires et les unités dans les calculs.} Précise toujours les conditions (lumière pour la substitution, CNTP pour $V_m = \SI{22,4}{L.mol^{-1}}$), écris le tableau d'avancement pour identifier le réactif limitant, et garde les unités à chaque ligne de calcul. Une combustion \og complète \fg{} suppose le dioxygène en excès : dis-le explicitement.
\end{enumerate}
\end{attention}

\newpage

\coursec{Exercices}

\coursec{Les alcanes}

\courssub{Structure des alcanes : le carbone tétragonal}

\begin{definition}[Alcane]
Les \cle{alcanes} sont des hydrocarbures saturés acycliques (chaîne ouverte) ne comportant que des liaisons covalentes simples \ce{C-C} et \ce{C-H}. Leur formule générale est $C_nH_{2n+2}$.
\end{definition}

\begin{propriete}[Géométrie du carbone tétragonal]
Dans les alcanes, chaque atome de carbone est \cle{hybridé $sp^3$}. Il forme quatre liaisons $\sigma$ dirigées vers les sommets d'un \cle{tétraèdre} régulier.
\begin{itemize}
    \item Longueur liaison \ce{C-H} $\approx \SI{1,09}{\angstrom}$ ; \ce{C-C} $\approx \SI{1,54}{\angstrom}$.
    \item Angles de liaison : $\widehat{HCH} = \widehat{CCH} = \widehat{CCC} \approx \SI{109,5}{\degree}$.
\end{itemize}
\end{propriete}

\begin{exemplebox}[Méthane \ce{CH4} et Éthane \ce{C2H6}]
\textbf{Méthane :} Carbone central, 4 H aux sommets du tétraèdre.
\[
\chemfig{
    C(-[0]H)(-[2]H)(<:[4]H)(<:[8]H)
}
\quad \rightarrow \quad
\text{Géométrie tétraédrique parfaite.}
\]

\textbf{Éthane :} Deux tétraèdres joints par un sommet commun (liaison \ce{C-C}).
\[
\chemfig{
    CH_3(-[0]H)(-[2]H)(<:[4]H)-CH_3(-[0]H)(-[2]H)(<:[4]H)
}
\]
\emph{Représentation de Cram (trait plein = plan, triangle plein = hors plan vers l'avant, trait pointillé = hors plan vers l'arrière).}
\end{exemplebox}

\begin{aretenir}[Conformations de l'éthane]
La rotation autour de la liaison \ce{C-C} simple est libre (faible barrière $\approx \SI{12}{kJ.mol^{-1}}$). Deux conformations limites :
\begin{itemize}
    \item \cle{Conformation éclipsée} : atomes d'hydrogène face à face (énergie maximale, instable).
    \item \cle{Conformation décalée (staggered)} : atomes d'hydrogène décalés de \SI{60}{\degree} (énergie minimale, stable).
\end{itemize}
À température ambiante, les molécules passent rapidement d'une conformation à l'autre ; on représente souvent la conformation décalée.
\end{aretenir}

\begin{savaistu}[Le carbone tétragonal au Cameroun]
La géométrie tétraédrique du carbone est à la base de la stéréochimie. Elle explique pourquoi le \emph{bili-bili} (bière de mil) ou le \emph{vin de palme} contiennent de l'éthanol (\ce{CH3-CH2-OH}) dont le carbone \ce{CH2} est tétragonal, et pourquoi les médicaments chiraux (ex. antipaludéens) ont des énantiomères aux activités différentes. Le Mont Cameroun, par son activité volcanique, émet naturellement du méthane (\ce{CH4}), premier alcane, dont la structure tétraédrique a été élucidée par Van't Hoff et Le Bel (1874).
\end{savaistu}

\courssub{Formule générale, formules développées et isomérie}

\begin{propriete}[Formules des alcanes acycliques]
\begin{itemize}
    \item \textbf{Formule brute (générale) :} $C_nH_{2n+2}$ ($n \ge 1$).
    \item \textbf{Formule semi-développée :} Liaisons \ce{C-C} explicites, groupements \ce{CH3}, \ce{CH2}, \ce{CH} groupés. Ex: \ce{CH3-CH2-CH3} (propane).
    \item \textbf{Formule développée :} Toutes les liaisons \ce{C-H} et \ce{C-C} tracées.
\end{itemize}
\end{propriete}

\begin{definition}[Isomérie de constitution (squelette)]
Des \cle{isomères de constitution} ont la même formule brute mais des formules développées différentes (enchaînement des atomes de carbone distinct). Pour les alcanes, c'est de l'isomérie de chaîne (linéaire vs ramifiée).
\begin{center}
\begin{tabularx}{\linewidth}{|c|X|X|}\hline
$n$ & Formule brute & Isomères (semi-développée + nom) \\\hline
1 & \ce{CH4} & Méthane \\\hline
2 & \ce{C2H6} & Éthane \\\hline
3 & \ce{C3H8} & Propane \\\hline
4 & \ce{C4H10} & \ce{CH3-CH2-CH2-CH3} (Butane) ; \ce{CH3-CH(CH3)-CH3} (2-Méthylpropane) \\\hline
5 & \ce{C5H12} & Pentane ; 2-Méthylbutane ; 2,2-Diméthylpropane \\\hline
\end{tabularx}
\end{center}
\end{definition}

\begin{propriete}[Cycloalcanes (Cyclanes)]
Hydrocarbures saturés \textbf{cycliques}. Formule générale : \cle{$C_nH_{2n}$} ($n \ge 3$). Ex: Cyclopropane, Cyclobutane, Cyclopentane, Cyclohexane (conformation chaise).
\end{propriete}

\courssub{Nomenclature des alcanes (Règles UICPA)}

\begin{methode}[Nommer un alcane (UICPA)]
\begin{enumerate}
    \item \textbf{Chaîne principale :} Plus longue chaîne continue d'atomes de carbone. Donne le radical de base (\emph{meth-, eth-, prop-, but-, pent-, hex-, hept-, oct-, non-, dec-}) + suffixe \emph{-ane}.
    \item \textbf{Numérotation :} Sens donnant les plus bas locants aux substituants (règle du \emph{premier point de différence}).
    \item \textbf{Substituants (groupes alkyl) :} Nommés par le radical correspondant (\emph{méthyl} \ce{-CH3}, \emph{éthyl} \ce{-CH2CH3}, \emph{propyl} \ce{-CH2CH2CH3}, \emph{isopropyl} \ce{-CH(CH3)2}, \emph{tert-butyle} \ce{-C(CH3)3}...). Préfixes multiplicatifs : \emph{di-, tri-, tétra-...}
    \item \textbf{Écriture finale :} \emph{locants-suffixes-substituants-radical principal-ane}. Trait d'union entre chiffre et lettre ; virgule entre chiffres. Ordre alphabétique des substituants (ignorer di-, tri-).
\end{enumerate}
\end{methode}

\begin{exemplebox}[Exemples de nomenclature]
\begin{itemize}
    \item \ce{CH3-CH2-CH(CH3)-CH3} $\rightarrow$ Chaîne 4C (butane), méthyle en C2 $\rightarrow$ \textbf{2-méthylbutane}.
    \item \ce{CH3-C(CH3)2-CH3} $\rightarrow$ Chaîne 3C (propane), 2 méthyles en C2 $\rightarrow$ \textbf{2,2-diméthylpropane}.
    \item \ce{CH3-CH2-CH(CH2CH3)-CH2-CH3} $\rightarrow$ Chaîne 5C (pentane), éthyle en C3 $\rightarrow$ \textbf{3-éthylpentane}.
\end{itemize}
\end{exemplebox}

\begin{attention}[Pièges fréquents]
\begin{itemize}
    \item Ne pas confondre \emph{isopropyl} (\ce{-CH(CH3)2}, substituant) et \emph{propane} (alcane).
    \item Choisir la \textbf{plus longue chaîne}, même si elle n'est pas horizontale sur le papier.
    \item \ce{C(CH3)4} s'appelle \textbf{2,2-diméthylpropane} (chaîne principale = 3C), pas "néopentane" (nom trivial toléré mais non UICPA).
\end{itemize}
\end{attention}

\courssub{Propriétés physiques des alcanes}

\begin{propriete}[Propriétés physiques]
\begin{itemize}
    \item \textbf{Polarité :} Molécules \cle{apolaires} (liaisons \ce{C-H} et \ce{C-C} faiblement polaires, symétrie tétraédrique).
    \item \textbf{Solubilité :} \cle{Insolubles dans l'eau} (solvant polaire) ; \cle{solubles dans les solvants apolaires} (éther, chloroforme, benzène, hexane). \emph{Miscibles entre eux}.
    \item \textbf{Densité :} $d < 1$ (flottent sur l'eau). $d_{\text{gaz}} < d_{\text{air}}$ pour $n \le 2$ ; $d_{\text{gaz}} > d_{\text{air}}$ pour $n \ge 3$.
    \item \textbf{État physique (CNTP) :}
    \begin{itemize}
        \item $n = 1 \to 4$ : Gaz (Méthane, Éthane, Propane, Butane).
        \item $n = 5 \to 17$ : Liquides (Essence, Kérosène, Gazole).
        \item $n \ge 18$ : Solides (Paraffines, Cires).
    \end{itemize}
    \item \textbf{Point d'ébullition :} Croît avec la masse molaire (forces de Van der Waals) ; \textbf{diminue avec le ramification} (surface de contact moindre). Ex: $T_{\text{éb}}(\text{n-butane}) = \SI{-0,5}{\celsius} > T_{\text{éb}}(\text{2-méthylpropane}) = \SI{-11,7}{\celsius}$.
\end{itemize}
\end{propriete}

\begin{savaistu}[Gaz domestique au Cameroun]
Les bouteilles de \emph{gaz domestique} (CIMENCAM, SONARA, TotalEnergies, Tradex) contiennent du \textbf{GPL (Gaz de Pétrole Liquéfié)} : mélange majoritairement de \textbf{butane} ($n$-butane et isobutane) et de \textbf{propane}. Sous pression (environ \SI{8}{bar}), ils sont liquides ; à l'ouverture du robinet, ils se vaporisent. L'additif odorant (mercaptan, odeur « chou/œuf pourri ») permet de détecter les fuites (gaz inodore et incolore à l'origine).
\end{savaistu}

\courssub{Propriétés chimiques : combustion et substitution}

\begin{propriete}[Inertie relative]
Les alcanes sont peu réactifs à température ambiante (liaisons \ce{C-C} et \ce{C-H} fortes, apolarité). Pas de réaction avec acides, bases, agents oxydants doux, halogènes (sans énergie).
\end{propriete}

\begin{aretenir}[Combustion complète (Réaction de destruction)]
Exothermique, lente à froid, violente à chaud (flamme). Produits : \ce{CO2} + \ce{H2O}.
\[
\ce{C_nH_{2n+2} + \frac{3n+1}{2} O2 -> n CO2 + (n+1) H2O}
\]
\textbf{Méthode d'équilibrage :} 1. Carbone $\rightarrow n\,\ce{CO2}$. 2. Hydrogène $\rightarrow (n+1)\,\ce{H2O}$. 3. Oxygène $\rightarrow \frac{2n + (n+1)}{2} = \frac{3n+1}{2}\,\ce{O2}$. Multiplier par 2 pour coefficients entiers.
\end{aretenir}

\begin{attention}[Combustion incomplète : DANGER MORTEL]
Manque de dioxygène (pièce close, brûleur mal réglé). Produits toxiques :
\begin{itemize}
    \item \cle{Monoxyde de carbone \ce{CO}} : Gaz incolore, inodore, \textbf{toxique} (fixe l'hémoglobine $\to$ asphyxie). $\ce{2 C_nH_{2n+2} + (2n+1) O2 -> 2n CO + (2n+2) H2O}$.
    \item \cle{Carbone (suie) \ce{C}} : Dépôts noirs, encrassement. $\ce{C_nH_{2n+2} + n O2 -> n C + (n+1) H2O}$.
\end{itemize}
\textbf{Ventilation obligatoire !}
\end{attention}

\begin{aretenir}[Halogénation : Substitution radicalaire]
Réaction avec \ce{Cl2} ou \ce{Br2} sous \textbf{rayonnement UV} (lumière) ou \textbf{chaleur} (\SI{> 250}{\celsius}).
\[
\ce{C_nH_{2n+2} + X2 ->[h\nu/\Delta] C_nH_{2n+1}X + HX} \quad (X = \ce{Cl}, \ce{Br})
\]
\textbf{Mécanisme en chaîne (type radicalaire) :}
\begin{enumerate}
    \item \textbf{Initiation :} $\ce{Cl-Cl ->[h\nu] 2 Cl^.}$ (homolyse).
    \item \textbf{Propagation :} $\ce{Cl^. + RH -> HCl + R^.}$ ; $\ce{R^. + Cl2 -> RCl + Cl^.}$.
    \item \textbf{Terminaison :} $\ce{2 Cl^. -> Cl2}$ ; $\ce{R^. + Cl^. -> RCl}$ ; $\ce{2 R^. -> R-R}$.
\end{enumerate}
\emph{Polysubstitution possible} (mélange \ce{RCl}, \ce{RCl2}, \ce{RCl3}...). Le brome est moins réactif, plus sélectif (caractère radicalaire).
\end{aretenir}

\begin{savaistu}[Importance des dérivés halogénés]
\begin{itemize}
    \item \textbf{Chlorofluorocarbures (CFC) :} Ex \ce{CCl3F} (Fréon). Anciens frigorigènes, propulseurs. Interdits (Protocole de Montréal) car détruisent la couche d'ozone.
    \item \textbf{Dichlorométhane \ce{CH2Cl2}, Trichloréthylène} : Solvants industriels (dégraissage, extraction caféine).
    \item \textbf{Chloroforme \ce{CHCl3}} : Ancien anesthésique, toxique (foie, reins).
    \item \textbf{Chlorure de vinyle \ce{CH2=CHCl}} (alcène halogéné) : Monomère du \textbf{PVC} (tuyaux CIMENCAM, câbles, fenêtres).
\end{itemize}
\end{savaistu}

\courssub{Préparation du méthane au laboratoire}

\begin{experience}[Préparation du méthane par hydrolyse du carbure d'aluminium]
\textbf{Principe :} \ce{Al4C3(s) + 12 H2O(l) -> 4 Al(OH)3(s) + 3 CH4(g)}.
\textbf{Dispositif :} Ballon (ou tube à essai) contenant \ce{Al4C3} + entonnoir à robinet (eau) $\to$ tube de sortie $\to$ \textbf{recueil par déplacement d'eau} (bac pneumatique + burette/éprouvette graduée inversée).
\textbf{Observations :} Ébullition locale, dégagement gazeux bouillant dans l'eau, gaz incolore, inodore, inflammable (flamme bleue non lumineuse).
\textbf{Interprétation :} \ce{Al4C3} est un carbure ionique (\ce{Al^3+}, \ce{C^4-}). L'ion \ce{C^4-} est une base très forte qui arrache \ce{H+} à l'eau $\to$ \ce{CH4}. L'hydroxyde d'aluminium \ce{Al(OH)3} précipite (gel blanc).
\end{experience}

\begin{propriete}[Justification du recueil par déplacement d'eau]
Le méthane est \cle{peu soluble dans l'eau} ($\approx \SI{22}{mg.L^{-1}}$ à \SI{20}{\celsius}) et \cle{moins dense que l'air} ($d \approx 0,55$). Il ne se dissout pas et chasse l'eau de la burette.
\end{propriete}

\begin{methode}[Calcul stœchiométrique avec tableau d'avancement]
\begin{enumerate}
    \item Écrire l'équation bilan équilibrée.
    \item Calculer $n_{\text{initial}}$ des réactifs (données : masse, volume gaz $\to n = V/V_m$).
    \item Tableau d'avancement ($x$ en mol) : Initial / Intermédiaire / Final (réactif limitant $\to x_{\text{max}}$).
    \item Calculer $n_{\text{final}}$ du produit cherché $\to$ Masse ou Volume ($V = n \times V_m$).
\end{enumerate}
\end{methode}

\courssub{Je vérifie mes connaissances}

\begin{exercice}[QCM : les alcanes]
Pour chaque affirmation, indique la (ou les) bonne(s) réponse(s).
\begin{enumerate}
\item La molécule de méthane \ce{CH4} est :
\begin{enumerate}
\item plane ; \item tétraédrique ; \item linéaire.
\end{enumerate}
\item L'angle entre deux liaisons \ce{C-H} dans le méthane vaut environ :
\begin{enumerate}
\item \SI{90}{\degree} ; \item \SI{109,5}{\degree} ; \item \SI{120}{\degree}.
\end{enumerate}
\item Les alcanes sont :
\begin{enumerate}
\item solubles dans l'eau ; \item insolubles dans l'eau ; \item solubles dans les solvants organiques.
\end{enumerate}
\item La combustion incomplète d'un alcane peut produire :
\begin{enumerate}
\item du monoxyde de carbone \ce{CO} ; \item du carbone (suie) ; \item uniquement \ce{CO2} et \ce{H2O}.
\end{enumerate}
\item La réaction d'un alcane avec le dichlore en présence de lumière est une réaction de :
\begin{enumerate}
\item addition ; \item substitution ; \item combustion.
\end{enumerate}
\end{enumerate}
\end{exercice}

\begin{exercice}[Vrai ou faux]
Réponds par \textbf{vrai} ou \textbf{faux} en justifiant chaque réponse.
\begin{enumerate}
\item La formule générale des alcanes est $C_nH_{2n+2}$.
\item Deux alcanes de même formule brute ont toujours les mêmes propriétés physiques.
\item Le butane et le 2-méthylpropane sont des isomères de constitution.
\item La combustion complète du méthane produit du dioxyde de carbone et de l'eau.
\item Les alcanes réagissent spontanément avec le dichlore dans l'obscurité.
\end{enumerate}
\end{exercice}

\begin{exercice}[Texte à trous]
Complète le texte suivant avec les mots ou expressions convenables : \emph{tétragonal ; substitution ; isomères ; apolaires ; radicalaire ; $C_nH_{2n+2}$}.

Dans une molécule d'alcane, chaque atome de carbone est \ldots\ldots\ldots\ldots : il forme quatre liaisons covalentes dirigées vers les sommets d'un tétraèdre. Les alcanes ont pour formule générale \ldots\ldots\ldots\ldots. Ce sont des molécules \ldots\ldots\ldots\ldots, ce qui explique leur insolubilité dans l'eau. Des composés de même formule brute mais de formules développées différentes sont des \ldots\ldots\ldots\ldots. En présence de lumière, le dichlore réagit avec un alcane par une réaction de \ldots\ldots\ldots\ldots, dont le mécanisme est de type \ldots\ldots\ldots\ldots.
\end{exercice}

\begin{exercice}[Définitions]
Définis brièvement les termes suivants :
\begin{enumerate}
\item alcane ; \item carbone tétragonal ; \item isomères de constitution ; \item réaction de substitution ; \item densité d'un gaz par rapport à l'air.
\end{enumerate}
\end{exercice}

\courssub{Je m'entraîne}

\begin{exercice}[Formules et nomenclature]
\begin{enumerate}
\item Écris la formule brute des alcanes comportant $n = 3$, $n = 6$ et $n = 8$ atomes de carbone, puis donne leur nom.
\item Écris la formule semi-développée des composés suivants : 2-méthylbutane ; 2,2-diméthylpropane ; 3-éthylpentane.
\item Nomme selon les règles de l'UICPA les composés de formules semi-développées :
\[\ce{CH3-CH(CH3)-CH2-CH3} \qquad \ce{CH3-CH2-CH(CH3)-CH2-CH2-CH3}\]
\end{enumerate}
\end{exercice}

\begin{exercice}[Combustion du butane]
Le gaz vendu en bouteilles au Cameroun (gaz domestique) est essentiellement du butane \ce{C4H10}.
\begin{enumerate}
\item Écris et équilibre l'équation de la combustion complète du butane dans le dioxygène.
\item Calcule le volume de dioxygène nécessaire à la combustion complète de \SI{29}{g} de butane, mesuré dans les conditions normales de température et de pression (CNTP).
\item Cite deux produits toxiques ou dangereux formés lors d'une combustion incomplète du butane dans une cuisine mal aérée.
\end{enumerate}
\emph{Données :} $M(\ce{C}) = \SI{12}{g.mol^{-1}}$ ; $M(\ce{H}) = \SI{1}{g.mol^{-1}}$ ; volume molaire dans les CNTP : $V_m = \SI{22,4}{L.mol^{-1}}$.
\end{exercice}

\begin{exercice}[Chloration de l'éthane]
On fait réagir du dichlore sur l'éthane \ce{C2H6} en présence de lumière.
\begin{enumerate}
\item Écris l'équation de la réaction de monochloration de l'éthane et nomme le produit organique obtenu.
\item Écris l'équation de la réaction de dichloration de l'éthane. Combien de dérivés dichlorés différents peut-on obtenir ? Écris leurs formules semi-développées et nomme-les.
\item Quel est le rôle de la lumière dans ces réactions ?
\end{enumerate}
\end{exercice}

\begin{exercice}[Un alcane gazeux]
Un alcane gazeux $A$ a une densité par rapport à l'air $d \approx 2$.
\begin{enumerate}
\item Détermine la masse molaire de $A$, puis sa formule brute.
\item Écris les formules semi-développées de tous les isomères de $A$ et nomme-les.
\end{enumerate}
\emph{Données :} masse molaire moyenne de l'air : $M_{\text{air}} = \SI{29}{g.mol^{-1}}$ ; $M(\ce{C}) = \SI{12}{g.mol^{-1}}$ ; $M(\ce{H}) = \SI{1}{g.mol^{-1}}$.
\end{exercice}

\courssub{Je me prépare au Probatoire}

\begin{exercice}[Les isomères du pentane et combustion]
\textbf{Partie A : isomérie.} On considère les alcanes de formule brute \ce{C5H12}.
\begin{enumerate}
\item Écris les formules semi-développées de tous les isomères de formule brute \ce{C5H12}.
\item Nomme chacun de ces isomères selon la nomenclature UICPA.
\item Pour l'isomère à chaîne linéaire, écris la formule développée complète en précisant la géométrie autour de chaque atome de carbone.
\end{enumerate}
\textbf{Partie B : combustion.} On réalise la combustion complète d'une masse $m = \SI{10}{g}$ de pentane \ce{C5H12} dans le dioxygène.
\begin{enumerate}
\item Écris et équilibre l'équation de la réaction de combustion complète du pentane.
\item Calcule la quantité de matière de pentane brûlé.
\item À l'aide d'un tableau d'avancement, détermine le volume de dioxygène nécessaire à cette combustion, mesuré dans les CNTP.
\item Calcule le volume de dioxyde de carbone dégagé dans les mêmes conditions.
\end{enumerate}
\emph{Données :} $M(\ce{C}) = \SI{12}{g.mol^{-1}}$ ; $M(\ce{H}) = \SI{1}{g.mol^{-1}}$ ; $M(\ce{O}) = \SI{16}{g.mol^{-1}}$ ; $V_m = \SI{22,4}{L.mol^{-1}}$ dans les CNTP.
\end{exercice}

\begin{exercice}[Identification d'un alcane par combustion]
La combustion complète d'une masse $m = \SI{2,9}{g}$ d'un alcane $A$ produit une masse $m' = \SI{8,8}{g}$ de dioxyde de carbone.
\begin{enumerate}
\item Écris l'équation générale de la combustion complète d'un alcane de formule $C_nH_{2n+2}$ dans le dioxygène.
\item Calcule la quantité de matière de \ce{CO2} formé.
\item Exprime la quantité de matière d'alcane brûlé en fonction de $n$, puis montre que la formule brute de $A$ est \ce{C4H10}.
\item Écris les formules semi-développées des isomères de $A$ et nomme-les.
\item L'un de ces isomères est le principal constituant du gaz domestique utilisé dans les cuisines camerounaises. Lequel ? Pourquoi est-il dangereux de faire brûler ce gaz dans une pièce mal ventilée ? Appuie ta réponse par une équation de réaction.
\end{enumerate}
\emph{Données :} $M(\ce{C}) = \SI{12}{g.mol^{-1}}$ ; $M(\ce{H}) = \SI{1}{g.mol^{-1}}$ ; $M(\ce{O}) = \SI{16}{g.mol^{-1}}$.
\end{exercice}

\begin{exercice}[Préparation du méthane au laboratoire]
Au laboratoire, on prépare le méthane par action de l'eau sur le carbure d'aluminium \ce{Al4C3}. Il se forme du méthane \ce{CH4} et de l'hydroxyde d'aluminium \ce{Al(OH)3}.
\begin{enumerate}
\item Écris et équilibre l'équation-bilan de la réaction.
\item Schématise et légende soigneusement le dispositif expérimental permettant de préparer le méthane et de le recueillir par déplacement d'eau (précise le nom de chaque élément du montage).
\item Pourquoi peut-on recueillir le méthane par déplacement d'eau ? Justifie à l'aide d'une propriété physique des alcanes.
\item On souhaite recueillir un volume $V = \SI{6,72}{L}$ de méthane, mesuré dans les CNTP. Calcule la quantité de matière de méthane correspondante.
\item À l'aide d'un tableau d'avancement, détermine la masse minimale de carbure d'aluminium nécessaire, l'eau étant en excès.
\item Le méthane est le principal constituant du gaz naturel. Écris l'équation de sa combustion complète et calcule le volume de dioxygène nécessaire à la combustion des \SI{6,72}{L} de méthane recueillis.
\end{enumerate}
\emph{Données :} $M(\ce{Al}) = \SI{27}{g.mol^{-1}}$ ; $M(\ce{C}) = \SI{12}{g.mol^{-1}}$ ; $V_m = \SI{22,4}{L.mol^{-1}}$ dans les CNTP.
\end{exercice}

\newpage

\coursec{Corrigés des exercices}

\begin{corrige}[Exercice 1 — QCM : les alcanes]
\begin{enumerate}
\item \textbf{Réponse b.} La molécule de méthane \ce{CH4} est \textbf{tétraédrique} : le carbone, hybridé $sp^3$, forme quatre liaisons dirigées vers les sommets d'un tétraèdre régulier.
\item \textbf{Réponse b.} L'angle $\widehat{HCH}$ vaut environ \SI{109,5}{\degree} (angle tétraédrique).
\item \textbf{Réponses b et c.} Les alcanes sont des molécules apolaires : ils sont \textbf{insolubles dans l'eau} (solvant polaire) mais \textbf{solubles dans les solvants organiques} apolaires (éther, hexane...).
\item \textbf{Réponses a et b.} En combustion incomplète (manque de dioxygène), il se forme du \textbf{monoxyde de carbone \ce{CO}} (gaz toxique) et/ou du \textbf{carbone \ce{C}} (suie). La réponse c correspond à la combustion \emph{complète}.
\item \textbf{Réponse b.} Il s'agit d'une réaction de \textbf{substitution} : un atome d'hydrogène de l'alcane est remplacé par un atome de chlore, selon un mécanisme radicalaire déclenché par la lumière.
\end{enumerate}
\end{corrige}

\begin{corrige}[Exercice 2 — Vrai ou faux]
\begin{enumerate}
\item \textbf{Vrai.} Les alcanes (hydrocarbures saturés acycliques) ont pour formule générale $C_nH_{2n+2}$ avec $n \ge 1$.
\item \textbf{Faux.} Deux alcanes de même formule brute peuvent être des isomères de constitution (ex. butane et 2-méthylpropane, \ce{C4H10}) ; ils ont des propriétés physiques différentes : $T_{\text{éb}}(\text{butane}) = \SI{-0,5}{\celsius}$ alors que $T_{\text{éb}}(\text{2-méthylpropane}) = \SI{-11,7}{\celsius}$.
\item \textbf{Vrai.} Butane \ce{CH3-CH2-CH2-CH3} et 2-méthylpropane \ce{CH3-CH(CH3)-CH3} ont la même formule brute \ce{C4H10} mais des enchaînements d'atomes différents : ce sont des isomères de constitution (de chaîne).
\item \textbf{Vrai.} $\ce{CH4 + 2 O2 -> CO2 + 2 H2O}$ : la combustion complète donne toujours \ce{CO2} et \ce{H2O}.
\item \textbf{Faux.} Les alcanes sont inertes vis-à-vis du dichlore dans l'obscurité et à température ambiante ; la substitution nécessite un apport d'énergie : lumière (UV) ou chauffage.
\end{enumerate}
\end{corrige}

\begin{corrige}[Exercice 3 — Texte à trous]
Dans une molécule d'alcane, chaque atome de carbone est \textbf{tétragonal} : il forme quatre liaisons covalentes dirigées vers les sommets d'un tétraèdre. Les alcanes ont pour formule générale \textbf{$C_nH_{2n+2}$}. Ce sont des molécules \textbf{apolaires}, ce qui explique leur insolubilité dans l'eau. Des composés de même formule brute mais de formules développées différentes sont des \textbf{isomères}. En présence de lumière, le dichlore réagit avec un alcane par une réaction de \textbf{substitution}, dont le mécanisme est de type \textbf{radicalaire}.
\end{corrige}

\begin{corrige}[Exercice 4 — Définitions]
\begin{enumerate}
\item \textbf{Alcane :} hydrocarbure saturé acyclique ne comportant que des liaisons simples \ce{C-C} et \ce{C-H}, de formule générale $C_nH_{2n+2}$.
\item \textbf{Carbone tétragonal :} atome de carbone hybridé $sp^3$ qui forme quatre liaisons covalentes simples dirigées vers les sommets d'un tétraèdre régulier (angles d'environ \SI{109,5}{\degree}).
\item \textbf{Isomères de constitution :} composés qui ont la même formule brute mais des formules développées différentes (enchaînement des atomes différent).
\item \textbf{Réaction de substitution :} réaction au cours de laquelle un atome (ou groupe d'atomes) d'une molécule est remplacé par un autre atome (ou groupe d'atomes). Ex. : $\ce{CH4 + Cl2 ->[h\nu] CH3Cl + HCl}$.
\item \textbf{Densité d'un gaz par rapport à l'air :} rapport de la masse d'un volume de gaz à la masse du même volume d'air dans les mêmes conditions ; elle se calcule par $d = \dfrac{M_{\text{gaz}}}{M_{\text{air}}} = \dfrac{M_{\text{gaz}}}{29}$ (sans unité).
\end{enumerate}
\end{corrige}

\begin{corrige}[Exercice 5 — Formules et nomenclature]
\begin{enumerate}
\item Formule générale $C_nH_{2n+2}$ :
\begin{itemize}
\item $n = 3$ : \ce{C3H8} : \textbf{propane} ;
\item $n = 6$ : \ce{C6H14} : \textbf{hexane} ;
\item $n = 8$ : \ce{C8H18} : \textbf{octane}.
\end{itemize}
\item Formules semi-développées :
\begin{itemize}
\item 2-méthylbutane : \ce{CH3-CH(CH3)-CH2-CH3} ;
\item 2,2-diméthylpropane : \ce{CH3-C(CH3)2-CH3} ;
\item 3-éthylpentane : \ce{CH3-CH2-CH(CH2-CH3)-CH2-CH3}.
\end{itemize}
\item Nomenclature :
\begin{itemize}
\item \ce{CH3-CH(CH3)-CH2-CH3} : chaîne principale de 4 carbones (butane), un méthyle en position 2 $\rightarrow$ \textbf{2-méthylbutane} ;
\item \ce{CH3-CH2-CH(CH3)-CH2-CH2-CH3} : chaîne principale de 6 carbones (hexane), un méthyle en position 3 $\rightarrow$ \textbf{3-méthylhexane}.
\end{itemize}
\end{enumerate}
\textbf{Point de vigilance :} pour le dernier composé, on choisit la chaîne la plus longue (6 C) et non la chaîne horizontale dessinée ; la numérotation donne le locant 3 dans les deux sens, le résultat est identique par symétrie.
\end{corrige}

\begin{corrige}[Exercice 6 — Combustion du butane]
\begin{enumerate}
\item Équation de la combustion complète :
\[ \ce{2 C4H10 + 13 O2 -> 8 CO2 + 10 H2O} \]
Vérification : C : $2\times4 = 8$ ; H : $2\times10 = 20 = 10\times2$ ; O : $13\times2 = 26 = 16 + 10$. L'équation est équilibrée.
\item Masse molaire du butane : $M(\ce{C4H10}) = 4\times12 + 10\times1 = \SI{58}{g.mol^{-1}}$.
Quantité de butane :
\[ n(\ce{C4H10}) = \frac{m}{M} = \frac{29}{58} = \SI{0,50}{mol} \]
D'après l'équation, $\dfrac{n(\ce{O2})}{n(\ce{C4H10})} = \dfrac{13}{2}$, donc :
\[ n(\ce{O2}) = \frac{13}{2}\times 0,50 = \SI{3,25}{mol} \]
Volume de dioxygène dans les CNTP :
\[ V(\ce{O2}) = n \times V_m = 3,25 \times 22,4 = \SI{72,8}{L} \]
\item Produits dangereux d'une combustion incomplète : le \textbf{monoxyde de carbone \ce{CO}} (gaz toxique, mortel par asphyxie) et le \textbf{carbone \ce{C}} (suie noire). D'où la nécessité d'une bonne ventilation de la cuisine.
\end{enumerate}
\end{corrige}

\begin{corrige}[Exercice 7 — Chloration de l'éthane]
\begin{enumerate}
\item Monochloration :
\[ \ce{C2H6 + Cl2 ->[h\nu] C2H5Cl + HCl} \]
Le produit organique \ce{CH3-CH2Cl} est le \textbf{chloroéthane}.
\item Dichloration :
\[ \ce{C2H6 + 2 Cl2 ->[h\nu] C2H4Cl2 + 2 HCl} \]
On peut obtenir \textbf{deux} dérivés dichlorés différents (isomères de position) :
\begin{itemize}
\item \ce{CH3-CHCl2} : \textbf{1,1-dichloroéthane} (les deux Cl sur le même carbone) ;
\item \ce{CH2Cl-CH2Cl} : \textbf{1,2-dichloroéthane} (un Cl sur chaque carbone).
\end{itemize}
\item La lumière (rayonnement UV) apporte l'énergie nécessaire à la \textbf{rupture homolytique} de la liaison \ce{Cl-Cl} : $\ce{Cl2 ->[h\nu] 2 Cl^.}$. C'est l'étape d'\textbf{initiation} du mécanisme radicalaire en chaîne ; sans lumière ni chauffage, la réaction ne se produit pas.
\end{enumerate}
\end{corrige}

\begin{corrige}[Exercice 8 — Un alcane gazeux]
\begin{enumerate}
\item Densité par rapport à l'air : $d = \dfrac{M(A)}{M_{\text{air}}}$, donc :
\[ M(A) = d \times M_{\text{air}} = 2 \times 29 = \SI{58}{g.mol^{-1}} \]
Pour un alcane : $M = 12n + (2n+2)\times1 = 14n + 2$.
\[ 14n + 2 = 58 \implies 14n = 56 \implies n = 4 \]
La formule brute de $A$ est \ce{C4H10} (c'est bien un gaz, car $n \le 4$).
\item Les deux isomères de \ce{C4H10} :
\begin{itemize}
\item \ce{CH3-CH2-CH2-CH3} : \textbf{butane} ;
\item \ce{CH3-CH(CH3)-CH3} : \textbf{2-méthylpropane}.
\end{itemize}
\end{enumerate}
\end{corrige}

\begin{corrige}[Exercice 9 — Les isomères du pentane et combustion (type Probatoire)]
\textbf{Partie A : isomérie.}
\begin{enumerate}
\item Les trois isomères de \ce{C5H12} :
\begin{itemize}
\item \ce{CH3-CH2-CH2-CH2-CH3} ;
\item \ce{CH3-CH(CH3)-CH2-CH3} ;
\item \ce{CH3-C(CH3)2-CH3}.
\end{itemize}
\item Noms UICPA : \textbf{pentane} ; \textbf{2-méthylbutane} ; \textbf{2,2-diméthylpropane}.
\item Pour le pentane, la formule développée (représentation en zigzag de la géométrie tétraédrique) est :
\[ \chemfig{CH_3-[:-30]CH_2-[:30]CH_2-[:-30]CH_2-[:30]CH_3} \]
Autour de \textbf{chaque} atome de carbone, la géométrie est \textbf{tétraédrique} (carbone tétragonal $sp^3$, angles voisins de \SI{109,5}{\degree}) ; la molécule n'est donc pas plane mais en zigzag.
\end{enumerate}
\textbf{Partie B : combustion.}
\begin{enumerate}
\item Équation de combustion complète :
\[ \ce{C5H12 + 8 O2 -> 5 CO2 + 6 H2O} \]
\item $M(\ce{C5H12}) = 5\times12 + 12\times1 = \SI{72}{g.mol^{-1}}$.
\[ n(\ce{C5H12}) = \frac{m}{M} = \frac{10}{72} \approx \SI{0,139}{mol} \]
\item Tableau d'avancement (le dioxygène est supposé en quantité suffisante, le pentane est le réactif limitant) :
\begin{center}
\begin{tabularx}{\linewidth}{|l|X|X|X|X|}\hline
\rowcolor{popblueL} Équation & \ce{C5H12} & $+~8\,\ce{O2}$ & $\rightarrow 5\,\ce{CO2}$ & $+~6\,\ce{H2O}$ \\\hline
Initial (mol) & $0,139$ & $n_0$ & $0$ & $0$ \\\hline
Intermédiaire & $0,139 - x$ & $n_0 - 8x$ & $5x$ & $6x$ \\\hline
Final (mol) & $0$ & $n_0 - 8x_{\max}$ & $5x_{\max}$ & $6x_{\max}$ \\\hline
\end{tabularx}
\end{center}
$x_{\max} = \SI{0,139}{mol}$. Le dioxygène nécessaire :
\[ n(\ce{O2}) = 8\,x_{\max} = 8 \times 0,139 \approx \SI{1,11}{mol} \]
\[ V(\ce{O2}) = 1,11 \times 22,4 \approx \SI{24,9}{L} \]
\item $n(\ce{CO2}) = 5\,x_{\max} = 5 \times 0,139 \approx \SI{0,694}{mol}$, d'où :
\[ V(\ce{CO2}) = 0,694 \times 22,4 \approx \SI{15,6}{L} \]
\end{enumerate}
\textbf{Barème indicatif (sur 10) :} A1 : 1,5 pt ; A2 : 1,5 pt ; A3 : 1 pt ; B1 : 1 pt ; B2 : 1 pt ; B3 : 2,5 pts (tableau 1 pt, calcul 1,5 pt) ; B4 : 1,5 pt.
\textbf{Points de vigilance :} ne pas oublier d'équilibrer l'équation avant tout calcul ; arrondir $10/72$ seulement en fin de calcul ; préciser les unités à chaque étape ; la géométrie tétraédrique doit être mentionnée explicitement en A3.
\end{corrige}

\begin{corrige}[Exercice 10 — Identification d'un alcane par combustion (type Probatoire)]
\begin{enumerate}
\item Équation générale :
\[ \ce{C_nH_{2n+2} + \frac{3n+1}{2} O2 -> n CO2 + (n+1) H2O} \]
\item $M(\ce{CO2}) = 12 + 2\times16 = \SI{44}{g.mol^{-1}}$.
\[ n(\ce{CO2}) = \frac{m'}{M} = \frac{8,8}{44} = \SI{0,20}{mol} \]
\item D'après l'équation, 1 mol d'alcane donne $n$ mol de \ce{CO2}, donc :
\[ n(A) = \frac{n(\ce{CO2})}{n} = \frac{0,20}{n} \ \text{mol} \]
La masse molaire de l'alcane s'écrit alors :
\[ M(A) = \frac{m}{n(A)} = \frac{2,9}{0,20/n} = 14,5\,n \ \si{g.mol^{-1}} \]
Or pour un alcane, $M(A) = 14n + 2$. D'où :
\[ 14,5\,n = 14n + 2 \implies 0,5\,n = 2 \implies n = 4 \]
La formule brute de $A$ est donc \ce{C4H10} (vérification : $M = 14\times4 + 2 = \SI{58}{g.mol^{-1}}$ et $n(A) = 2,9/58 = \SI{0,05}{mol}$, qui donne bien $4\times0,05 = \SI{0,20}{mol}$ de \ce{CO2}).
\item Isomères de \ce{C4H10} :
\begin{itemize}
\item \ce{CH3-CH2-CH2-CH3} : \textbf{butane} ;
\item \ce{CH3-CH(CH3)-CH3} : \textbf{2-méthylpropane}.
\end{itemize}
\item Le principal constituant du gaz domestique est le \textbf{butane}. Dans une pièce mal ventilée, le manque de dioxygène provoque une \textbf{combustion incomplète} qui produit du monoxyde de carbone \ce{CO}, gaz incolore, inodore et mortel (il bloque le transport du dioxygène par l'hémoglobine) :
\[ \ce{2 C4H10 + 9 O2 -> 8 CO + 10 H2O} \]
Il peut aussi se former de la suie \ce{C}. D'où l'importance d'aérer les cuisines.
\end{enumerate}
\textbf{Barème indicatif (sur 10) :} Q1 : 1 pt ; Q2 : 1 pt ; Q3 : 3 pts (expression de $n(A)$ : 1 pt, équation en $n$ : 1 pt, conclusion : 1 pt) ; Q4 : 2 pts ; Q5 : 3 pts (identification : 1 pt, danger justifié : 1 pt, équation : 1 pt).
\textbf{Points de vigilance :} bien distinguer $n$ (indice de la formule brute) et $n(A)$ (quantité de matière) ; la démonstration du Q3 doit être rédigée clairement, sans partir de la conclusion ; l'équation de combustion incomplète doit être équilibrée.
\end{corrige}

\begin{corrige}[Exercice 11 — Préparation du méthane au laboratoire (type Probatoire)]
\begin{enumerate}
\item Équation-bilan équilibrée :
\[ \ce{Al4C3(s) + 12 H2O(l) -> 4 Al(OH)3(s) + 3 CH4(g)} \]
Vérification : Al : 4 ; C : 3 ; H : $12\times2 = 24 = 4\times3 + 3\times4$ ; O : $12 = 4\times3$.
\item Schéma du dispositif :
\begin{popfigure}
\begin{tikzpicture}[scale=0.95, every node/.style={font=\small}]
% Ballon de réaction
\draw[thick, popblue] (0,0) circle (0.9);
\draw[thick, popblue] (-0.25,0.85) -- (-0.25,1.5) -- (0.25,1.5) -- (0.25,0.85);
\node[popdark] at (0,-0.15) {\ce{Al4C3} + eau};
% Entonnoir à robinet
\draw[thick, popdark] (-0.15,1.5) -- (-0.15,2.3) -- (-0.55,3.1) -- (0.25,3.1) -- (-0.02,2.3);
\draw[thick, popdark] (-0.02,2.3) -- (-0.02,1.5);
\draw[thick, poporange] (-0.28,2.05) -- (0.1,2.05);
\node[right, popdark] at (0.3,2.9) {entonnoir à robinet (eau)};
% Tube de sortie
\draw[thick, popgreen] (0.25,1.3) -- (1.6,1.3) -- (1.6,-0.4) -- (4.3,-0.4) -- (4.3,0.9);
% Cuvette à eau
\draw[thick, popblue] (2.6,-1.2) rectangle (6.2,-0.1);
\fill[popblue!25] (2.62,-1.18) rectangle (6.18,-0.35);
\node[popdark] at (4.4,-1.5) {cuvette à eau (bac pneumatique)};
% Éprouvette inversée
\draw[thick, poppurple] (3.9,1.6) -- (3.9,-0.6) (4.7,1.6) -- (4.7,-0.6);
\draw[thick, poppurple] (3.9,1.6) -- (4.7,1.6);
\fill[popblue!25] (3.92,-0.58) rectangle (4.68,0.55);
\draw[popgreen, thick, ->] (4.3,0.9) -- (4.3,0.6);
\node[popdark] at (4.3,1.95) {éprouvette graduée inversée};
\node[popdark] at (4.3,0.25) {\ce{CH4}};
\node[popdark] at (0,-1.5) {ballon};
\end{tikzpicture}
\legende{Préparation du méthane par action de l'eau sur le carbure d'aluminium et recueil par déplacement d'eau.}
\end{popfigure}
\item Le méthane peut être recueilli par déplacement d'eau car il est \textbf{très peu soluble dans l'eau} (propriété commune aux alcanes, molécules apolaires) : il ne se dissout pas et chasse l'eau de l'éprouvette.
\item Quantité de méthane :
\[ n(\ce{CH4}) = \frac{V}{V_m} = \frac{6,72}{22,4} = \SI{0,300}{mol} \]
\item Tableau d'avancement (eau en excès, \ce{Al4C3} réactif limitant) :
\begin{center}
\begin{tabularx}{\linewidth}{|l|X|X|X|X|}\hline
\rowcolor{popblueL} Équation & \ce{Al4C3} & $+~12\,\ce{H2O}$ & $\rightarrow 4\,\ce{Al(OH)3}$ & $+~3\,\ce{CH4}$ \\\hline
Initial (mol) & $n_0$ & excès & $0$ & $0$ \\\hline
Intermédiaire & $n_0 - x$ & excès & $4x$ & $3x$ \\\hline
Final (mol) & $0$ & excès & $4x_{\max}$ & $3x_{\max}$ \\\hline
\end{tabularx}
\end{center}
On veut $3\,x_{\max} = 0,300$ mol, donc $x_{\max} = \SI{0,100}{mol}$ et $n_0 = x_{\max} = \SI{0,100}{mol}$ de \ce{Al4C3}.
Masse molaire : $M(\ce{Al4C3}) = 4\times27 + 3\times12 = 108 + 36 = \SI{144}{g.mol^{-1}}$.
\[ m(\ce{Al4C3}) = n \times M = 0,100 \times 144 = \SI{14,4}{g} \]
\item Combustion complète du méthane :
\[ \ce{CH4 + 2 O2 -> CO2 + 2 H2O} \]
Il faut $n(\ce{O2}) = 2 \times n(\ce{CH4}) = 2 \times 0,300 = \SI{0,600}{mol}$, soit :
\[ V(\ce{O2}) = 0,600 \times 22,4 = \SI{13,44}{L} \]
\end{enumerate}
\textbf{Barème indicatif (sur 10) :} Q1 : 1,5 pt ; Q2 : 2 pts (dispositif cohérent : 1 pt, légendes : 1 pt) ; Q3 : 1 pt ; Q4 : 1 pt ; Q5 : 2,5 pts (tableau : 1 pt, calcul : 1,5 pt) ; Q6 : 2 pts (équation : 1 pt, calcul : 1 pt).
\textbf{Points de vigilance :} le schéma doit montrer l'entonnoir plongeant dans le ballon et l'éprouvette \emph{inversée} remplie d'eau au départ ; ne pas écrire \ce{Al4C3} en solution (c'est un solide) ; vérifier que l'eau est bien indiquée en excès pour justifier $n_0 = x_{\max}$.
\end{corrige}

\newpage

\coursec{Fiche bilan}

\coursec{Structure des alcanes : le carbone tétragonal}

\begin{savaistu}[Le gaz butane dans nos foyers]
Au Cameroun, la bouteille de gaz butane (mélange de butane et de propane) est indispensable pour la cuisson. Ce gaz, stocké sous pression à l'état liquide, s'enflamme au contact d'une étincelle pour produire une flamme bleue très chaude. Pourquoi le butane brûle-t-il si bien ? Pourquoi est-il gazeux à température ambiante alors que l'huile de palme est liquide ? La réponse se trouve dans la structure moléculaire des \cle{alcanes}, la famille chimique du butane.
\end{savaistu}

\courssub{Le carbone tétraédrique (hybridation \texorpdfstring{$\mathrm{sp^3}$}{sp3})}

\begin{definition}[Carbone tétragonal]
Un atome de carbone lié à quatre atomes par des liaisons simples \ce{C-C} ou \ce{C-H} adopte une géométrie \cle{tétraédrique}. Les quatre liaisons pointent vers les sommets d'un tétraèdre régulier ; l'angle entre deux liaisons vaut \SI{109,5}{\degree}. On dit que le carbone est \cle{hybridé $\mathrm{sp^3}$}.
\end{definition}

\begin{popfigure}
\begin{center}
\begin{tikzpicture}[scale=1.2]
% Tétraèdre pour le méthane
\coordinate (C) at (0,0,0);
\coordinate (H1) at (1,1,1);
\coordinate (H2) at (1,-1,-1);
\coordinate (H3) at (-1,1,-1);
\coordinate (H4) at (-1,-1,1);

% Arêtes du tétraèdre (lignes pointillées pour la profondeur)
\draw[popmuted,dashed,thin] (H1) -- (H2) (H1) -- (H3) (H1) -- (H4) (H2) -- (H3) (H2) -- (H4) (H3) -- (H4);
% Liaisons C-H
\draw[popblue,thick] (C) -- (H1) node[above right] {\ce{H}};
\draw[popblue,thick] (C) -- (H2) node[below left] {\ce{H}};
\draw[popblue,thick] (C) -- (H3) node[above left] {\ce{H}};
\draw[popblue,thick] (C) -- (H4) node[below right] {\ce{H}};
% Atome de carbone
\fill[popdark] (C) circle (3pt) node[below=2pt] {\textbf{C}};
% Angle
\draw[poporange,<->] ($(C)!0.8!(H1)$) arc (45:-35:0.8) node[midway,right] {\SI{109,5}{\degree}};
\end{tikzpicture}
\end{center}
\legende{Représentation de Cram du méthane \ce{CH4} : le carbone au centre, les quatre hydrogènes aux sommets d'un tétraèdre. Les traits pleins sortent du plan, les pointillés rentrent.}
\end{popfigure}

\begin{experience}[Modèle moléculaire du méthane et de l'éthane]
\textbf{Matériel} : kit de modèles moléculaires (boules et tiges) ou pâte à modeler et cure-dents.
\textbf{Protocole} :
\begin{enumerate}
  \item Construire le méthane : une boule noire (C) au centre, quatre boules blanches (H) aux sommets d'un tétraèdre.
  \item Construire l'éthane \ce{C2H6} : assembler deux tétraèdres \ce{CH3} par une liaison \ce{C-C}.
  \item Faire pivoter un groupe \ce{CH3} par rapport à l'autre autour de l'axe \ce{C-C}.
\end{enumerate}
\textbf{Observations} :
\begin{itemize}
  \item Toutes les liaisons \ce{C-H} sont équivalentes ; la molécule n'est pas plane.
  \item Dans l'éthane, la rotation autour de la liaison \ce{C-C} est libre à température ambiante.
\end{itemize}
\textbf{Interprétation} : La géométrie tétraédrique minimise les répulsions entre les paires d'électrons de liaison (théorie VSEPR). La rotation autour des liaisons simples \ce{C-C} génère différentes \cle{conformations}.
\end{experience}

\begin{definition}[Conformations]
Les \cle{conformations} sont les différents arrangements spatiaux d'une \emph{même} molécule obtenus par rotation autour d'une liaison simple \ce{C-C}. Pour l'éthane, on distingue principalement :
\begin{itemize}
  \item la conformation \cle{éclipsée} (hydrogènes face à face, énergie maximale) ;
  \item la conformation \cle{décalée} ou \cle{anti} (hydrogènes décalés de \SI{60}{\degree}, énergie minimale, plus stable).
\end{itemize}
\textbf{Attention} : Les conformations ne sont \textbf{pas} des isomères ; elles s'interconvertissent rapidement sans rompre de liaison.
\end{definition}

\begin{popfigure}
\begin{center}
\begin{tikzpicture}[scale=0.9]
% Conformation éclipsée (projection de Newman)
\begin{scope}[xshift=-4cm]
\draw[popline] (0,0) circle (1.5);
\foreach \i in {0,120,240} {
  \draw[popblue,thick] (0,0) -- (\i:1.5) node[anchor=\i+180] {\ce{H}};
  \draw[poporange,thick] (0,0) -- (\i:0.7) node[anchor=\i] {\ce{H}};
}
\fill[popdark] (0,0) circle (2pt) node[below=4pt] {C avant};
\node at (0,-2) {\textbf{Éclipsée}};
\node at (0,-2.5) {\small (énergie haute)};
\end{scope}

% Conformation décalée
\begin{scope}[xshift=4cm]
\draw[popline] (0,0) circle (1.5);
\foreach \i in {0,120,240} {
  \draw[popblue,thick] (0,0) -- (\i:1.5) node[anchor=\i+180] {\ce{H}};
  \draw[poporange,thick] (0,0) -- (\i+60:0.7) node[anchor=\i+60+180] {\ce{H}};
}
\fill[popdark] (0,0) circle (2pt) node[below=4pt] {C avant};
\node at (0,-2) {\textbf{Décalée (anti)}};
\node at (0,-2.5) {\small (énergie basse)};
\end{scope}
\end{tikzpicture}
\end{center}
\legende{Projections de Newman de l'éthane : vue le long de l'axe \ce{C-C}. Carbone avant (traits pleins vers H), carbone arrière (traits courts vers H).}
\end{popfigure}

\coursec{Formule générale, formules développées et isomérie}

\courssub{Formule générale des alcanes acycliques}

\begin{propriete}[Formule générale]
Les alcanes acycliques (chaîne ouverte) forment une série homologue de formule générale :
\[
\boxed{\mathrm{C}_n\mathrm{H}_{2n+2} \quad (n \geq 1)}
\]
Chaque passage d'un homologue au suivant ajoute un groupe \ce{CH2} (masse molaire + \SI{14}{g.mol^{-1}}).
\end{propriete}

\begin{exemplebox}[Premiers membres de la série]
\begin{center}
\begin{tabularx}{\linewidth}{|c|c|c|c|}
\hline
\rowcolor{popblueL} \textbf{n} & \textbf{Formule brute} & \textbf{Nom usuel (IUPAC)} & \textbf{Formule semi-développée} \\
\hline
1 & \ce{CH4} & Méthane & \ce{CH4} \\
2 & \ce{C2H6} & Éthane & \ce{CH3-CH3} \\
3 & \ce{C3H8} & Propane & \ce{CH3-CH2-CH3} \\
4 & \ce{C4H10} & Butane & \ce{CH3-CH2-CH2-CH3} \\
5 & \ce{C5H12} & Pentane & \ce{CH3-CH2-CH2-CH2-CH3} \\
6 & \ce{C6H14} & Hexane & \ce{CH3-(CH2)4-CH3} \\
\hline
\end{tabularx}
\end{center}
\end{exemplebox}

\courssub{Isomérie de constitution}

\begin{definition}[Isomères de constitution]
Des \cle{isomères de constitution} ont la \emph{même formule brute} mais un \emph{enchaînement d'atomes différent} (chaîne carbonée ramifiée vs linéaire). Ils ont des propriétés physiques et chimiques distinctes.
\end{definition}

\begin{methode}[Recherche systématique des isomères d'un alcane $\mathrm{C}_n\mathrm{H}_{2n+2}$ ($n \leq 5$)]
\begin{enumerate}
  \item Écrire la chaîne linéaire (carbone $n$).
  \item Raccourcir la chaîne principale d'un carbone ($n-1$) et placer ce carbone en ramification (méthyl) sur un carbone \emph{non terminal} (sinon on retrouve la chaîne linéaire).
  \item Raccourcir encore ($n-2$) et placer les carbures restants en ramifications (méthyl, éthyl) en respectant la tétravalence du carbone.
  \item Vérifier qu'aucune structure n'est redondante (symétrie).
\end{enumerate}
\end{methode}

\begin{exemplebox}[Isomères de $\mathrm{C_4H_{10}}$ ($n=4$) et $\mathrm{C_5H_{12}}$ ($n=5$)]
\textbf{Butane ($\mathrm{C_4H_{10}}$) : 2 isomères}
\begin{center}
\begin{tikzpicture}[scale=0.7, baseline=(current bounding box.center)]
% n-butane
\node at (0,0) {\chemfig{CH_3-CH_2-CH_2-CH_3}};
\node at (0,-1) {\textbf{Butane} (n-butane)};
% isobutane
\node at (6,0) {\chemfig{CH_3-CH(-[2]CH_3)-CH_3}};
\node at (6,-1) {\textbf{2-Méthylpropane} (isobutane)};
\end{tikzpicture}
\end{center}

\textbf{Pentane ($\mathrm{C_5H_{12}}$) : 3 isomères}
\begin{center}
\begin{tikzpicture}[scale=0.7, baseline=(current bounding box.center)]
% n-pentane
\node at (0,0) {\chemfig{CH_3-CH_2-CH_2-CH_2-CH_3}};
\node at (0,-1) {\textbf{Pentane} (n-pentane)};
% isopentane
\node at (6,0) {\chemfig{CH_3-CH(-[2]CH_3)-CH_2-CH_3}};
\node at (6,-1) {\textbf{2-Méthylbutane} (isopentane)};
% neopentane
\node at (12,0) {\chemfig{C(-[0]CH_3)(-[2]CH_3)(-[4]CH_3)-CH_3}};
\node at (12,-1) {\textbf{2,2-Diméthylpropane} (néopentane)};
\end{tikzpicture}
\end{center}
\end{exemplebox}

\courssub{Les cyclanes (alcanes cycliques)}

\begin{definition}[Cyclane]
Un \cle{cyclane} est un alcane dont la chaîne carbonée forme un cycle. Sa formule générale est :
\[
\boxed{\mathrm{C}_n\mathrm{H}_{2n} \quad (n \geq 3)}
\]
Les cycles les plus stables sont le cyclopentane ($n=5$) et le cyclohexane ($n=6$) ; le cyclopropane et le cyclobutane sont tendus (angles \ce{C-C-C} contraints).
\end{definition}

\begin{exemplebox}[Exemples de cyclanes]
\begin{center}
\begin{tabularx}{\linewidth}{|c|c|c|}
\hline
\rowcolor{popblueL} \textbf{Nom} & \textbf{Formule brute} & \textbf{Représentation} \\
\hline
Cyclopropane & $\mathrm{C_3H_6}$ & \chemfig{*3(---)} \\
Cyclobutane & $\mathrm{C_4H_8}$ & \chemfig{*4(---)} \\
Cyclopentane & $\mathrm{C_5H_{10}}$ & \chemfig{*5(---)} \\
Cyclohexane & $\mathrm{C_6H_{12}}$ & \chemfig{*6(---)} (siège) \\
\hline
\end{tabularx}
\end{center}
\end{exemplebox}

\coursec{Nomenclature des alcanes (règles IUPAC)}

\begin{methode}[Nommer un alcane ramifié]
\begin{enumerate}
  \item \textbf{Chaîne principale} : identifier la plus longue chaîne carbonée continue (pas nécessairement horizontale). Elle donne le nom de base (\emph{-ane}).
  \item \textbf{Numérotation} : numéroter la chaîne de manière à donner les \emph{indices les plus petits} aux substituants (règle du « premier point de différence »).
  \item \textbf{Substituants} : identifier les groupes alkyles (méthyl \ce{-CH3}, éthyl \ce{-CH2CH3}, propyl \ce{-CH2CH2CH3}, isopropyl \ce{-CH(CH3)2}, tert-butyl \ce{-C(CH3)3}...).
  \item \textbf{Écriture du nom} : \emph{indices}-noms des substituants (par \textbf{ordre alphabétique}, préfixes \emph{di-}, \emph{tri-} ignorés pour l'ordre) -\emph{nom de la chaîne principale}.
  \item Séparer les chiffres par des virgules, les chiffres des lettres par des traits d'union.
\end{enumerate}
\end{methode}

\begin{exemplebox}[Applications de la nomenclature]
\begin{center}
\begin{tabularx}{\linewidth}{|c|X|c|}
\hline
\rowcolor{popblueL} \textbf{Structure} & \textbf{Analyse} & \textbf{Nom IUPAC} \\
\hline
\chemfig{CH_3-CH(-[2]CH_3)-CH_2-CH_3} & Chaîne 4C $\to$ butane ; méthyl en C2 & \textbf{2-Méthylbutane} \\
\hline
\chemfig{CH_3-CH(-[2]CH_2CH_3)-CH_2-CH_3} & Chaîne 5C $\to$ pentane ; éthyl en C3 & \textbf{3-Éthylpentane} \\
\hline
\chemfig{CH_3-C(-[2]CH_3)(-[6]CH_3)-CH_2-CH_3} & Chaîne 4C $\to$ butane ; 2 méthyls en C2 & \textbf{2,2-Diméthylbutane} \\
\hline
\chemfig{CH_3-CH(-[2]CH_3)-CH(-[6]CH_3)-CH_3} & Chaîne 4C $\to$ butane ; méthyls en C2 et C3 & \textbf{2,3-Diméthylbutane} \\
\hline
\end{tabularx}
\end{center}
\end{exemplebox}

\begin{attention}[Pièges fréquents en nomenclature]
\begin{itemize}
  \item Ne pas confondre « plus longue chaîne » et « chaîne écrite horizontalement ».
  \item L'ordre alphabétique : \textbf{éthyl} avant \textbf{méthyl} (e avant m), \textbf{isopropyl} avant \textbf{méthyl} (i avant m). Les préfixes \emph{di-, tri-, tétra-} ne comptent pas pour l'ordre.
  \item Indices les plus petits : pour \ce{CH3-CH(CH3)-CH2-CH(CH3)-CH3}, numéroter de gauche à droite donne 2,4 ; de droite à gauche donne 2,4 aussi $\to$ choisir la numérotation qui donne l'indice le plus petit au \emph{premier} substituant cité alphabétiquement (ici équivalent).
  \item Ne jamais écrire « 1-méthyl » : un substituant en bout de chaîne allonge la chaîne principale.
\end{itemize}
\end{attention}

\coursec{Propriétés physiques des alcanes}

\begin{propriete}[Propriétés physiques des alcanes acycliques]
\begin{itemize}
  \item \textbf{État} : Gaz pour $n=1$ à $4$ (méthane à butane) ; liquides pour $n=5$ à $17$ (essence, kérosène) ; solides cireux pour $n \geq 18$ (paraffines).
  \item \textbf{Solubilité} : \cle{Insolubles dans l'eau} (apolaires), miscibles entre eux et dans les solvants apolaires (éther, chloroforme, benzène).
  \item \textbf{Densité} : $d < 1$ (flottent sur l'eau) ; augmente avec $n$ (de \SI{0,42}{g.mL^{-1}} pour le méthane liquide à \SI{0,78}{g.mL^{-1}} pour l'hexane).
  \item \textbf{Température d'ébullition $T_{\text{éb}}$} : Croît \emph{régulièrement} avec $n$ (forces de Van der Waals plus intenses pour les chaînes plus longues).
  \item \textbf{Ininflammabilité} : Tous sont inflammables ; les vapeurs des liquides forment des mélanges explosifs avec l'air.
\end{itemize}
\end{propriete}

\begin{popfigure}
\begin{center}
\begin{tikzpicture}
\begin{axis}[
  width=\linewidth, height=7cm,
  xlabel={Nombre d'atomes de carbone $n$},
  ylabel={$T_{\text{éb}}$ (\si{\celsius})},
  grid=major, grid style={popline},
  xmin=0, xmax=10, ymin=-180, ymax=150,
  tick label style={font=\small},
  label style={font=\small},
]
\addplot[popcurve, mark=*, mark size=2pt, popblue] coordinates {
  (1,-161.5) (2,-88.6) (3,-42.1) (4,-0.5) (5,36.1) (6,68.7) (7,98.4) (8,125.7) (9,150.8)
};
\addplot[poporange, mark=square*, mark size=2pt, only marks] coordinates {
  (3,-33) (4,12) (5,49) (6,81) (7,118)
};
\legend{Alcanes linéaires, Cyclanes (référence)}
\end{axis}
\end{tikzpicture}
\end{center}
\legende{Évolution de la température d'ébullition des alcanes linéaires (bleu) et comparaison avec quelques cyclanes (orange).}
\end{popfigure}

\begin{savaistu}[Le gaz domestique au Cameroun]
La bouteille « gaz butane » distribuée par la SCDP (Société Camerounaise de Dépôts Pétroliers) contient en réalité un mélange \textbf{butane/propane} (souvent 70/30 ou 50/50 selon la saison). Le propane ($T_{\text{éb}} = \SI{-42}{\celsius}$) assure la vaporisation par temps frais (régions de l'Ouest, nuit), tandis que le butane ($T_{\text{éb}} = \SI{-0,5}{\celsius}$) fournit l'essentiel de l'énergie. L'odeur caractéristique (mercaptans ajoutés) permet de détecter les fuites.
\end{savaistu}

\coursec{Propriétés chimiques : combustion et substitution}

\courssub{Combustion (réaction de destruction)}

\begin{definition}[Combustion]
La \cle{combustion} est la réaction complète d'un alcane avec le dioxygène \ce{O2} (air) sous l'action d'une flamme ou d'une étincelle. C'est une réaction d'oxydo-réduction très exothermique.
\end{definition}

\begin{propriete}[Combustion complète]
Équation générale (à savoir équilibrer pour tout $n$) :
\[
\ce{C_nH_{2n+2} + \frac{3n+1}{2} O2 -> n CO2 + (n+1) H2O}
\]
Si le coefficient de \ce{O2} est fractionnaire, multiplier toute l'équation par 2.
\end{propriete}

\begin{methode}[Équilibrer une combustion complète]
\begin{enumerate}
  \item Écrire les réactifs et produits : \ce{C_nH_{2n+2} + O2 -> CO2 + H2O}.
  \item Équilibrer le \textbf{carbone} : $n$ \ce{CO2}.
  \item Équilibrer l'\textbf{hydrogène} : $n+1$ \ce{H2O}.
  \item Équilibrer l'\textbf{oxygène} : $2n + (n+1) = 3n+1$ atomes O $\to$ $\frac{3n+1}{2}$ \ce{O2}.
  \item Si fraction $\to$ multiplier par 2.
\end{enumerate}
\end{methode}

\begin{exemplebox}[Exemples d'équations équilibrées]
\begin{align*}
\text{Méthane :} \quad & \ce{CH4 + 2 O2 -> CO2 + 2 H2O} \\
\text{Éthane :} \quad & \ce{C2H6 + \frac{7}{2} O2 -> 2 CO2 + 3 H2O} \quad \text{ou} \quad \ce{2 C2H6 + 7 O2 -> 4 CO2 + 6 H2O} \\
\text{Butane :} \quad & \ce{C4H10 + \frac{13}{2} O2 -> 4 CO2 + 5 H2O} \quad \text{ou} \quad \ce{2 C4H10 + 13 O2 -> 8 CO2 + 10 H2O} \\
\text{Pentane :} \quad & \ce{C5H12 + 8 O2 -> 5 CO2 + 6 H2O}
\end{align*}
\end{exemplebox}

\begin{propriete}[Combustion incomplète]
En cas de \textbf{défaut d'oxygène} (pièce mal aérée, brûleur mal réglé), la combustion est incomplète et produit du \textbf{monoxyde de carbone \ce{CO}} (gaz incolore, inodore, \textbf{très toxique}) et/ou du \textbf{carbone \ce{C}} (suie, dépôts noirs).
\[
\ce{C_nH_{2n+2} + \frac{n+1}{2} O2 -> n CO + (n+1) H2O} \quad \text{(produit CO)}
\]
\[
\ce{C_nH_{2n+2} + (n+1) O2 -> n C + (n+1) H2O} \quad \text{(produit suie)}
\]
\end{propriete}

\begin{attention}[Danger du monoxyde de carbone \ce{CO}]
\ce{CO} se fixe sur l'hémoglobine (affinité 200-250 fois supérieure à \ce{O2}) $\to$ \textbf{asphyxie} (maux de tête, nausées, mort). \textbf{Ne jamais utiliser un réchaud à gaz ou un groupe électrogène dans une pièce close sans aération permanente.} Au Cameroun, des décès surviennent chaque année par inhalation de \ce{CO} (groupe électrogène dans le couloir, réchaud dans la chambre).
\end{attention}

\courssub{Substitution halogénée (réaction de fonctionnalisation)}

\begin{definition}[Réaction de substitution]
Une \cle{substitution} est une réaction où un atome (ou groupe d'atomes) est remplacé par un autre atome (ou groupe) sans modification du squelette carboné. Pour les alcanes : \ce{H} remplacé par \ce{X} (halogène).
\end{definition}

\begin{propriete}[Halogénation des alcanes]
\begin{itemize}
  \item \textbf{Réactifs} : \ce{Cl2} (dichlore) ou \ce{Br2} (dibrome). \ce{F2} trop violent, \ce{I2} trop lent.
  \item \textbf{Condition} : \textbf{Lumière UV} (soleil, lampe à mercure) ou \textbf{chaleur} (\SIrange{250}{400}{\celsius}). La lumière casse la liaison \ce{X-X} $\to$ radicaux libres \ce{X^.}.
  \item \textbf{Produits} : Mélange de dérivés halogénés mono-, di-, tri-, tétra-substitués + \ce{HX}.
\end{itemize}
\end{propriete}

\begin{exemplebox}[Chloration du méthane (mécanisme radicalaire simplifié)]
\textbf{Initiation} (lumière) : $\ce{Cl2 ->[h\nu] 2 Cl^.}$
\textbf{Propagation} :
\begin{align*}
\ce{Cl^. + CH4 -> HCl + ^.CH3} \\
\ce{^.CH3 + Cl2 -> CH3Cl + Cl^.}
\end{align*}
\textbf{Terminaison} : $\ce{2 Cl^. -> Cl2}$, $\ce{Cl^. + ^.CH3 -> CH3Cl}$, $\ce{2 ^.CH3 -> C2H6}$.
\textbf{Bilan global (monosubstitution)} :
\[
\ce{CH4 + Cl2 ->[h\nu] CH3Cl + HCl}
\]
La réaction se poursuit : \ce{CH3Cl -> CH2Cl2 -> CHCl3 -> CCl4}.
\end{exemplebox}

\begin{savaistu}[Dérivés halogénés importants et enjeux environnementaux]
\begin{itemize}
  \item \textbf{Chloroforme \ce{CHCl3}} : ancien anesthésique (abandonné : toxicité hépatique), solvant pour le caoutchouc, extraction d'huiles essentielles.
  \item \textbf{Tétrachlorure de carbone \ce{CCl4}} : ancien extincteur, solvant, précurseur des fréons (interdit : cancérogène, hépatotoxique).
  \item \textbf{Fréons (CFC, HCFC)} : \ce{CCl2F2} (fréon-12), \ce{CCl3F} (fréon-11). Utilisés comme fluides frigorigènes (frigos, clim) et propulseurs aérosols. \textbf{Problème} : dans la stratosphère, UV libère \ce{Cl^.} qui détruit catalytiquement l'ozone \ce{O3} ($\ce{Cl^. + O3 -> ClO^. + O2}$, $\ce{ClO^. + O -> Cl^. + O2}$). Protocole de Montréal (1987) : élimination progressive. Remplacés par HFC (ex. \ce{CH2FCF3}, R-134a) sans chlore, mais gaz à effet de serre puissants.
\end{itemize}
\end{savaistu}

\coursec{Préparation du méthane au laboratoire}

\begin{experience}[Préparation du méthane par hydrolyse du carbure d'aluminium]
\textbf{Principe} : Le carbure d'aluminium \ce{Al4C3} réagit violemment avec l'eau pour donner de l'hydroxyde d'aluminium \ce{Al(OH)3} (précipité blanc) et du méthane \ce{CH4} (gaz).
\[
\ce{Al4C3(s) + 12 H2O(l) -> 4 Al(OH)3(s) + 3 CH4(g) ^}
\]

\textbf{Matériel} :
\begin{itemize}
  \item Tube à essai muni d'un bouchon à deux trous (ou erlenmeyer + bouchon percé)
  \item Entonnoir à thistle (entonnoir tubulé) pour l'introduction de l'eau
  \item Tube de sortie de gaz vers une cuve pneumatique (recueil par déplacement d'eau) ou vers un brûleur
  \item Carbure d'aluminium \ce{Al4C3} (solide grisâtre, attention : réagit avec l'humidité de l'air)
  \item Eau distillée
\end{itemize}

\textbf{Dispositif (schéma)} :
\begin{center}
\begin{tikzpicture}[scale=0.8, every node/.style={font=\small}]
% Erlenmeyer
\draw[popline,thick] (-1.5,-2) -- (-1,2) -- (1,2) -- (1.5,-2);
\draw[popline,thick] (-1.5,-2) .. controls (-1.5,-2.5) and (1.5,-2.5) .. (1.5,-2);
% Bouchon
\draw[popline,thick,fill=popcream] (-1.2,2) rectangle (1.2,2.4);
% Entonnoir thistle
\draw[popline,thick] (-0.2,2.4) -- (-0.2,4) .. controls (-0.2,4.5) and (0.5,4.5) .. (0.5,4);
\draw[popline,thick] (0.5,4) -- (0.5,2.4);
\draw[popline,thick,fill=popblue!20] (-0.2,4) rectangle (0.5,4.5) node[midway] {Eau};
% Tube sortie
\draw[popline,thick] (0.7,2.4) -- (0.7,3) -- (3,3) -- (3,1);
\draw[popline,thick,->] (3,1) -- (3.5,1) node[right] {\ce{CH4}};
% Cuve pneumatique (simplifiée)
\draw[popline,thick] (4,-1) -- (4,2) -- (6,2) -- (6,-1);
\draw[popline,thick,fill=popblue!10] (4,0) -- (6,0) -- (6,2) -- (4,2);
\draw[popline,thick,->] (5,1.5) -- (5,1.8) node[above] {Eau};
% Légendes
\node at (0,-2.8) {Erlenmeyer + \ce{Al4C3}};
\node at (5,2.5) {Cuve pneumatique};
\node at (5,-1.5) {Récipient gradué};
\end{tikzpicture}
\end{center}

\textbf{Protocole} :
\begin{enumerate}
  \item Placer quelques grammes de \ce{Al4C3} dans l'erlenmeyer (manipuler vite : hygroscopique).
  \item Monter le dispositif, vérifier l'étanchéité (boucher la sortie gaz, verser un peu d'eau dans l'entonnoir : pas de bulles).
  \item Ouvrir la sortie gaz vers la cuve pneumatique remplie d'eau.
  \item Verser l'eau goutte à goutte dans l'entonnoir : réaction immédiate, dégagement de bulles.
  \item Recueillir plusieurs tubes de gaz (le premier contient de l'air résiduel).
  \item Tester l'inflammabilité : approcher une flamme du bec de sortie $\to$ flamme bleue non lumineuse.
\end{enumerate}

\textbf{Observations} :
\begin{itemize}
  \item Dégagement gazeux rapide, réaction exothermique (erlenmeyer tiède).
  \item Formation d'un précipité blanc gélatineux (\ce{Al(OH)3}) dans l'erlenmeyer.
  \item Gaz incolore, inodore, moins dense que l'air ($d_{\ce{CH4}} \approx 0,55 d_{\text{air}}$), insoluble dans l'eau.
  \item Flamme bleue, non fumivore (combustion complète).
\end{itemize}

\textbf{Interprétation} : Hydrolyse du carbure ionique \ce{Al4C3} (contenant l'anion \ce{C^4-}) par l'eau (acide de Brønsted) $\to$ \ce{CH4}. L'équation ionique : $\ce{C^4- + 4 H2O -> CH4 + 4 OH-}$.
\end{experience}

\begin{attention}[Sécurité et pièges]
\begin{itemize}
  \item \ce{Al4C3} réagit avec l'humidité de l'air : conserver au sec, manipuler vite, refermer le flacon.
  \item La réaction est exothermique : ne pas verser l'eau trop vite (projection d'\ce{Al(OH)3} chaud).
  \item \ce{CH4} est \textbf{inflammable} et forme des mélanges explosifs avec l'air (5-15\% en volume). \textbf{Interdit de recueillir au-dessus d'une flamme.} Éloigner toute source d'inflammation pendant le recueil.
  \item Vérifier l'étanchéité \emph{avant} d'introduire l'eau (risque de projection de solide chaud).
\end{itemize}
\end{attention}

\coursec{Bilan et entraînement}

\begin{aretenir}[Checklist Probatoire — Alcanes]
\begin{itemize}
  \item[$\square$] Je sais représenter le carbone tétragonal (tétraèdre, \SI{109,5}{\degree}, $\mathrm{sp^3}$) et distinguer conformation (rotation \ce{C-C}) et isomérie de constitution.
  \item[$\square$] Je connais les formules générales : alcanes $\mathrm{C}_n\mathrm{H}_{2n+2}$, cyclanes $\mathrm{C}_n\mathrm{H}_{2n}$.
  \item[$\square$] Je sais écrire \textbf{toutes} les formules semi-développées des isomères pour $n \leq 5$ (1 pour $n=1,2,3$ ; 2 pour $n=4$ ; 3 pour $n=5$).
  \item[$\square$] Je maîtrise la nomenclature IUPAC : chaîne la plus longue, indices minimaux, ordre alphabétique des substituants (méthyl, éthyl, propyl, isopropyl, tert-butyl).
  \item[$\square$] Je connais les propriétés physiques : insolubles dans l'eau, $d<1$, $T_{\text{éb}}$ croît avec $n$, état gazeux/liquide/solide selon $n$.
  \item[$\square$] Je sais équilibrer la combustion complète ($\ce{CO2}+\ce{H2O}$) et incomplète ($\ce{CO}$ ou $\ce{C}+\ce{H2O}$) pour tout alcane.
  \item[$\square$] Je sais écrire les équations de substitution \ce{Cl2}/\ce{Br2} (condition : lumière/UV) et nommer les produits (ex. \ce{CH3Cl} chlorométhane, \ce{CH2Cl2} dichlorométhane, \ce{CHCl3} chloroforme, \ce{CCl4} tétrachlorure de carbone).
  \item[$\square$] Je cite des applications/risques : \ce{CO} toxique (aération), fréons (couche d'ozone), \ce{CCl4} cancérogène.
  \item[$\square$] Je sais décrire et schématiser la préparation du \ce{CH4} par \ce{Al4C3 + 12 H2O -> 4 Al(OH)3 + 3 CH4} (dispositif : erlenmeyer, entonnoir thistle, recueil pneumatique).
\end{itemize}
\end{aretenir}

\begin{exoresolu}[Nomenclature et isomérie]
\textbf{Énoncé} : Donner la formule semi-développée et le nom IUPAC de tous les isomères de constitution de l'alcane de formule brute $\mathrm{C_6H_{14}}$.
\textbf{Solution détaillée} :
\begin{enumerate}
  \item Chaîne linéaire (6C) : \chemfig{CH_3-CH_2-CH_2-CH_2-CH_2-CH_3} $\to$ \textbf{Hexane}.
  \item Chaîne 5C + 1 méthyl : positions 2 et 3 (position 4 = 2 par symétrie, position 1 ou 5 = hexane).
    \begin{itemize}
      \item Méthyl en C2 : \chemfig{CH_3-CH(-[2]CH_3)-CH_2-CH_2-CH_3} $\to$ \textbf{2-Méthylpentane}.
      \item Méthyl en C3 : \chemfig{CH_3-CH_2-CH(-[2]CH_3)-CH_2-CH_3} $\to$ \textbf{3-Méthylpentane}.
    \end{itemize}
  \item Chaîne 4C + 2 carbones en ramifications :
    \begin{itemize}
      \item Deux méthyls sur le même carbone (C2) : \chemfig{CH_3-C(-[2]CH_3)(-[6]CH_3)-CH_2-CH_3} $\to$ \textbf{2,2-Diméthylbutane}.
      \item Deux méthyls sur carbones adjacents (C2 et C3) : \chemfig{CH_3-CH(-[2]CH_3)-CH(-[6]CH_3)-CH_3} $\to$ \textbf{2,3-Diméthylbutane}.
    \end{itemize}
  \item Chaîne 3C impossible (nécessiterait 3 carbones en ramifications $\to$ tétravalence dépassée).
\end{enumerate}
\textbf{Bilan} : 5 isomères pour $\mathrm{C_6H_{14}}$.
\end{exoresolu}

\begin{exoresolu}[Combustion et stœchiométrie]
\textbf{Énoncé} : On brûle complètement \SI{2,80}{g} de butane (\ce{C4H10}, $M = \SI{58,1}{g.mol^{-1}}$) dans un excès de dioxygène. Calculer le volume de \ce{CO2} formé aux CNTP ($V_{\text{m}} = \SI{22,4}{L.mol^{-1}}$).
\textbf{Solution détaillée} :
\begin{enumerate}
  \item Quantité de matière de butane : $n_{\ce{C4H10}} = \frac{m}{M} = \frac{2,80}{58,1} = \SI{0,0482}{mol}$.
  \item Équation équilibrée : $\ce{2 C4H10 + 13 O2 -> 8 CO2 + 10 H2O}$.
    Rapport stœchiométrique : $\frac{n_{\ce{CO2}}}{n_{\ce{C4H10}}} = \frac{8}{2} = 4$.
  \item $n_{\ce{CO2}} = 4 \times n_{\ce{C4H10}} = 4 \times 0,0482 = \SI{0,193}{mol}$.
  \item Volume aux CNTP : $V_{\ce{CO2}} = n_{\ce{CO2}} \times V_{\text{m}} = 0,193 \times 22,4 = \SI{4,32}{L}$.
\end{enumerate}
\textbf{Réponse} : On obtient \SI{4,32}{L} de dioxyde de carbone aux CNTP.
\end{exoresolu}

\begin{exercice}[Entraînement — Nomenclature et propriétés]
\begin{enumerate}
  \item Écrire la formule semi-développée et donner le nom IUPAC de l'isomère de $\mathrm{C_5H_{12}}$ qui a un carbone quaternaire.
  \item Nommer : \ce{CH3-CH2-CH(CH3)-CH2-CH2-CH3} et \ce{CH3-C(CH3)2-CH2-CH3}.
  \item Parmi les affirmations suivantes, cocher la/les vraie(s) :
    \begin{itemize}
      \item[a)] Le cyclohexane a pour formule brute $\mathrm{C_6H_{14}}$.
      \item[b)] La combustion incomplète du propane peut produire du \ce{CO}.
      \item[c)] La chloration du méthane à l'obscurité totale se fait rapidement.
      \item[d)] Le 2-méthylbutane et le pentane ont la même température d'ébullition.
    \end{itemize}
  \item Écrire l'équation de la combustion complète du 2,2-diméthylpropane ($\mathrm{C_5H_{12}}$).
  \item On réalise la substitution du chlore sur l'éthane. Écrire l'équation de la monosubstitution et donner le nom du produit organique formé.
\end{enumerate}
\end{exercice}

\begin{corrige}[Exercice n]
\begin{enumerate}
  \item \chemfig{C(-[0]CH_3)(-[2]CH_3)(-[4]CH_3)-CH_3} $\to$ \textbf{2,2-Diméthylpropane} (néopentane).
  \item \ce{CH3-CH2-CH(CH3)-CH2-CH2-CH3} : chaîne 6C (hexane), méthyl en C3 $\to$ \textbf{3-Méthylhexane}. \\
        \ce{CH3-C(CH3)2-CH2-CH3} : chaîne 4C (butane), deux méthyls en C2 $\to$ \textbf{2,2-Diméthylbutane}.
  \item Vraies : \textbf{b)} seulement. a) Faux : cyclohexane $\mathrm{C_6H_{12}}$. c) Faux : lumière/UV indispensable. d) Faux : $T_{\text{éb}}$(2-méthylbutane) \SI{27,8}{\celsius} < $T_{\text{éb}}$(pentane) \SI{36,1}{\celsius} (ramification $\downarrow$ surface de contact $\downarrow$ forces Van der Waals).
  \item $\ce{C5H12 + 8 O2 -> 5 CO2 + 6 H2O}$.
  \item $\ce{C2H6 + Cl2 ->[h\nu] C2H5Cl + HCl}$ ; produit : \textbf{chloroéthane}.
\end{enumerate}
\end{corrige}

\findecours

\end{document}
